% Probabilités — Mathématiques Tle Tech — Cours signé OnBuch+
% Programme officiel MINESEC. Compiler avec : tectonic N05-probabilites.tex
\newcommand{\DOCMATIERE}{Mathématiques}
\newcommand{\DOCNIVEAU}{Tle Tech}
\newcommand{\DOCMODULE}{Mathématiques – classe de Terminale technique}
\newcommand{\DOCLECON}{N05}
\newcommand{\DOCTITRE}{Probabilités}
% OnBuch+ — préambule « cours signés » ENRICHI (compilé avec tectonic / XeLaTeX)
% Système visuel « Pop » d'OnBuch+ : fond crème (jamais blanc), bordures encre
% épaisses, ombres « dures » décalées, titres Archivo Black en capitales.
% Version MATHÉMATIQUES : graphes (pgfplots/tikz), boîtes pédagogiques
% (définition, propriété, méthode, attention, le savais-tu, exercice résolu,
% exercice, TP, bilan), page de garde et signature OnBuch+ ancrée en bas.
%
% Macros attendues dans main.tex AVANT \input{preamble} :
%   \DOCMATIERE  \DOCNIVEAU  \DOCMODULE  \DOCLECON (numéro)  \DOCTITRE
\documentclass[11pt]{article}

\usepackage{fontspec}
\usepackage[french]{babel}
\usepackage[a4paper,margin=2cm,top=3.4cm,bottom=2.8cm,headheight=1.4cm,footskip=1.3cm]{geometry}
\usepackage{amsmath,amssymb,amsfonts}
\usepackage{mathtools} % \xmapsto, \coloneqq... souvent utilisés par les modèles
\usepackage{cancel} % simplifications barrées
% \abs{z} (module, valeur absolue) et \norm{u} (norme), utilisés par les modèles
\providecommand{\abs}{}\let\abs\relax\DeclarePairedDelimiter{\abs}{\lvert}{\rvert}
\providecommand{\norm}{}\let\norm\relax\DeclarePairedDelimiter{\norm}{\lVert}{\rVert}
\usepackage[table]{xcolor} % [table] : fournit aussi \rowcolors (alternance de lignes)
\usepackage{pagecolor}
\usepackage{tikz}
\usetikzlibrary{arrows.meta,positioning,calc,shapes.geometric,shapes.misc,shapes.arrows,decorations.pathmorphing,decorations.pathreplacing,decorations.markings,patterns,shadows,mindmap,trees,fit,backgrounds,matrix,intersections,angles,quotes}
\usepackage{pgfplots}
\usepackage{pgf-pie} % diagrammes circulaires \pie (statistiques)
\pgfplotsset{compat=1.18}
\usepgfplotslibrary{fillbetween,groupplots} % aires entre courbes (calcul intégral)
\usepackage{enumitem}
\usepackage{fancyhdr}
% [most] charge toutes les bibliothèques tcolorbox (listings, minted, raster,
% documentation...) et gonfle inutilement la mémoire TeX sur un document long
% (~300 boîtes) : on ne charge que ce qui est réellement utilisé.
\usepackage[skins,breakable]{tcolorbox}
\usepackage{graphicx}
\usepackage{colortbl}
\usepackage{booktabs}
\usepackage{tabularx}
\usepackage{diagbox} % case d'en-tête diagonale des tableaux à double entrée
% Colonnes extensibles centrée / gauche / droite (C, L, R), souvent utilisées par les modèles
\newcolumntype{C}{>{\centering\arraybackslash}X}
\newcolumntype{L}{>{\raggedright\arraybackslash}X}
\newcolumntype{R}{>{\raggedleft\arraybackslash}X}
\usepackage{array}
\usepackage{multicol}
\usepackage{siunitx}
% parse-numbers=false : accepte aussi \SI{2,5 \times 10^{-3}}{...}
\sisetup{locale=FR,output-decimal-marker={,},group-minimum-digits=4,parse-numbers=false}
\DeclareSIUnit\ohm{\ensuremath{\symup{\Omega}}} % U+2126 absent de la police
\usepackage{qrcode}
% \degree : commande fréquemment utilisée par les modèles pour un angle ou une
% température en dehors de \SI{}{} (non fournie par les paquets chargés).
\providecommand{\degree}{^{\circ}}
% \qed (fin de démonstration), utilisé par les modèles sans amsthm
\providecommand{\qed}{\hfill$\square$}
% \barre : double barre des valeurs interdites dans un tableau de variations (tkz-tab)
\providecommand{\barre}{\ensuremath{\|}}
% environnement proof (amsthm non chargé) utilisé par les modèles
\ifdefined\proof\else\newenvironment{proof}[1][Démonstration]{\par\noindent\textit{#1.}\ }{\qed\par}\fi
\usepackage{lastpage}
\usepackage{hyperref}
\hypersetup{hidelinks}

% ============================================================
% Palette « Pop » OnBuch+ — mêmes hex que la classe P (plus_ui.dart)
% ============================================================
\definecolor{poporange}{HTML}{F59321}
\definecolor{popdark}{HTML}{E07A0C}
\definecolor{popcream}{HTML}{FFF6E8}
\definecolor{popink}{HTML}{1C1714}
\definecolor{poppurple}{HTML}{7A5AE0}
\definecolor{popgreen}{HTML}{2E9E6A}
\definecolor{poppink}{HTML}{E0457B}
\definecolor{popblue}{HTML}{2F7DE1}
\definecolor{popgold}{HTML}{C98A12}
\definecolor{popmuted}{HTML}{8A7F76}
\definecolor{popline}{HTML}{EADFD2}
\definecolor{popgray}{HTML}{8A7F76}
\colorlet{popgrey}{popgray}
% Alias tolérants (le modèle nomme parfois les couleurs différemment)
\colorlet{popwhite}{white}
\colorlet{popblack}{popink}
\colorlet{popred}{poppink}
\colorlet{popyellow}{popgold}
% Teintes claires pour fonds de tableaux / zones
\colorlet{popblueL}{popblue!10!white}
\colorlet{poporangeL}{poporange!14!white}
\colorlet{popgreenL}{popgreen!12!white}
\colorlet{poppurpleL}{poppurple!12!white}
\colorlet{poppinkL}{poppink!12!white}
\colorlet{popgoldL}{popgold!14!white}

\pagecolor{popcream}

% ============================================================
% Typographie
% ============================================================
\defaultfontfeatures{Ligatures=TeX}
\setmainfont{PlusJakartaSans-Medium.ttf}[Path=../../../fonts/,BoldFont=PlusJakartaSans-Bold.ttf]
% Maths en sans-serif (Fira Math) avec chiffres et lettres droites de Plus
% Jakarta, pour que les formules se fondent dans le texte.
\usepackage[math-style=ISO,bold-style=ISO]{unicode-math}
\setmathfont{FiraMath-Regular.otf}[Path=../../../fonts/]
\setmathfont{PlusJakartaSans-Medium.ttf}[Path=../../../fonts/,range={up/num,up/{Latin,latin},\mathup}]
\usepackage{icomma} % virgule décimale sans espace en mode math
% Caractères mathématiques Unicode absents des polices de texte (Plus Jakarta,
% Archivo Black) : rendus par leur équivalent mathématique, en texte comme en
% formule — sinon ils s'affichent en carré vide (ex. « Arithmétique dans ℤ »).
\usepackage{newunicodechar}
\newunicodechar{ℝ}{\ensuremath{\mathbb{R}}}
\newunicodechar{ℤ}{\ensuremath{\mathbb{Z}}}
\newunicodechar{ℂ}{\ensuremath{\mathbb{C}}}
\newunicodechar{ℕ}{\ensuremath{\mathbb{N}}}
\newunicodechar{ℚ}{\ensuremath{\mathbb{Q}}}
\newunicodechar{𝔻}{\ensuremath{\mathbb{D}}}
\newunicodechar{α}{\ensuremath{\alpha}}
\newunicodechar{β}{\ensuremath{\beta}}
\newunicodechar{γ}{\ensuremath{\gamma}}
\newunicodechar{δ}{\ensuremath{\delta}}
\newunicodechar{ε}{\ensuremath{\varepsilon}}
\newunicodechar{θ}{\ensuremath{\theta}}
\newunicodechar{λ}{\ensuremath{\lambda}}
\newunicodechar{μ}{\ensuremath{\mu}}
\newunicodechar{σ}{\ensuremath{\sigma}}
\newunicodechar{τ}{\ensuremath{\tau}}
\newunicodechar{φ}{\ensuremath{\varphi}}
\newunicodechar{ψ}{\ensuremath{\psi}}
\newunicodechar{ω}{\ensuremath{\omega}}
\newunicodechar{Σ}{\ensuremath{\Sigma}}
% majuscules produites par \MakeUppercase dans les titres (α-aminés -> Α)
\newunicodechar{Α}{\ensuremath{\alpha}}
\newunicodechar{Β}{\ensuremath{\beta}}
\newunicodechar{Γ}{\ensuremath{\Gamma}}
\newunicodechar{Δ}{\ensuremath{\Delta}}
\newunicodechar{Ε}{\ensuremath{\varepsilon}}
\newunicodechar{Θ}{\ensuremath{\Theta}}
\newunicodechar{Λ}{\ensuremath{\Lambda}}
\newunicodechar{Μ}{\ensuremath{\mu}}
\newunicodechar{Τ}{\ensuremath{\tau}}
\newunicodechar{Φ}{\ensuremath{\Phi}}
\newunicodechar{Ψ}{\ensuremath{\Psi}}
\newunicodechar{Ω}{\ensuremath{\Omega}}
\newunicodechar{↦}{\ensuremath{\mapsto}}
\newunicodechar{∈}{\ensuremath{\in}}
\newunicodechar{∉}{\ensuremath{\notin}}
\newunicodechar{∧}{\ensuremath{\wedge}}
\newunicodechar{⇔}{\ensuremath{\Leftrightarrow}}
\newunicodechar{⟺}{\ensuremath{\Longleftrightarrow}}
\newunicodechar{⇒}{\ensuremath{\Rightarrow}}
\newunicodechar{⟹}{\ensuremath{\Longrightarrow}}
\newunicodechar{∘}{\ensuremath{\circ}}
\newunicodechar{∪}{\ensuremath{\cup}}
\newunicodechar{∩}{\ensuremath{\cap}}
\newunicodechar{≡}{\ensuremath{\equiv}}
\newunicodechar{⊂}{\ensuremath{\subset}}
\newunicodechar{⊥}{\ensuremath{\perp}}
\newunicodechar{∃}{\ensuremath{\exists}}
\newunicodechar{∀}{\ensuremath{\forall}}
\newunicodechar{∛}{\ensuremath{\sqrt[3]{\,}}}
\newunicodechar{∜}{\ensuremath{\sqrt[4]{\,}}}
\newunicodechar{⃗}{\ensuremath{{}^{\to}}}
\newunicodechar{ⁿ}{\ifmmode^{n}\else\textsuperscript{\ensuremath{n}}\fi}
\newunicodechar{ˣ}{\ifmmode^{x}\else\textsuperscript{\ensuremath{x}}\fi}
\newunicodechar{ᵇ}{\ifmmode^{b}\else\textsuperscript{\ensuremath{b}}\fi}
\newunicodechar{ʳ}{\ifmmode^{r}\else\textsuperscript{\ensuremath{r}}\fi}
\newunicodechar{ᵘ}{\ifmmode^{u}\else\textsuperscript{\ensuremath{u}}\fi}
\newunicodechar{ʸ}{\ifmmode^{y}\else\textsuperscript{\ensuremath{y}}\fi}
\newunicodechar{ᵃ}{\ifmmode^{a}\else\textsuperscript{\ensuremath{a}}\fi}
\newunicodechar{ᵉ}{\ifmmode^{e}\else\textsuperscript{\ensuremath{e}}\fi}
\newunicodechar{ᵏ}{\ifmmode^{k}\else\textsuperscript{\ensuremath{k}}\fi}
\newunicodechar{ᵖ}{\ifmmode^{p}\else\textsuperscript{\ensuremath{p}}\fi}
\newunicodechar{ᴺ}{\ifmmode^{N}\else\textsuperscript{\ensuremath{N}}\fi}
\newunicodechar{ᶜ}{\ifmmode^{c}\else\textsuperscript{\ensuremath{c}}\fi}
\newunicodechar{ʰ}{\ifmmode^{h}\else\textsuperscript{\ensuremath{h}}\fi}
\newunicodechar{ᵠ}{\ifmmode^{q}\else\textsuperscript{\ensuremath{q}}\fi}
\newunicodechar{ₙ}{\ifmmode_{n}\else\textsubscript{\ensuremath{n}}\fi}
\newunicodechar{ₐ}{\ifmmode_{a}\else\textsubscript{\ensuremath{a}}\fi}
\newunicodechar{ᵢ}{\ifmmode_{i}\else\textsubscript{\ensuremath{i}}\fi}
\newunicodechar{ᵣ}{\ifmmode_{r}\else\textsubscript{\ensuremath{r}}\fi}
\newunicodechar{ₖ}{\ifmmode_{k}\else\textsubscript{\ensuremath{k}}\fi}
\newunicodechar{ᵦ}{\ifmmode_{\beta}\else\textsubscript{\ensuremath{\beta}}\fi}
\newunicodechar{ₓ}{\ifmmode_{x}\else\textsubscript{\ensuremath{x}}\fi}
\newfontfamily\popdisplay{ArchivoBlack-Regular.ttf}[Path=../../../fonts/]
\newfontfamily\poplogo{SpaceGrotesk-Bold.ttf}[Path=../../../fonts/]
\newfontfamily\popsemi{PlusJakartaSans-SemiBold.ttf}[Path=../../../fonts/]
\newfontfamily\popxbold{PlusJakartaSans-ExtraBold.ttf}[Path=../../../fonts/]
\newfontfamily\popxboldls{PlusJakartaSans-ExtraBold.ttf}[Path=../../../fonts/,LetterSpace=3.0]

\setlist[enumerate]{leftmargin=1.7em,itemsep=5pt,topsep=4pt}
\setlist[itemize]{leftmargin=1.7em,itemsep=4pt,topsep=4pt,label={\color{poporange}\textbullet}}
\setlength{\parindent}{0pt}
\setlength{\parskip}{5pt}
\renewcommand{\arraystretch}{1.25}

% pgfplots : style Pop par défaut
\pgfplotsset{
  every axis/.append style={axis line style={popink,line width=1pt},tick style={popink},
    grid style={popline,line width=0.5pt},label style={font=\small},tick label style={font=\footnotesize}},
  popcurve/.style={poporange,line width=1.8pt,no markers,smooth},
}
% Même style disponible dans un \draw TikZ simple (hors pgfplots)
\tikzset{popcurve/.style={poporange,line width=1.8pt}}

% ============================================================
% Pill / squiggle
% ============================================================
\newcommand{\poppill}[3]{%
  \tikz[baseline=-0.55ex]\node[draw=popink,line width=1.2pt,rounded corners=2.6pt,%
    fill=#2,inner xsep=6.5pt,inner ysep=3pt]{%
    \popxboldls\fontsize{7}{7}\selectfont\color{#3}\MakeUppercase{#1}};%
}
\newcommand{\popsquiggle}[1]{%
  \tikz[baseline=-0.4ex]{
    \draw[#1,line width=2.6pt,line cap=round]
      plot [smooth] coordinates {(0,0.09) (0.34,0.01) (0.68,0.21) (1.02,0.01) (1.36,0.10)};
  }%
}
% Mot-clé surligné (marqueur orange sous le texte)
\newcommand{\cle}[1]{{\popxbold\color{popdark}#1}}

% ============================================================
% En-tête / pied
% ============================================================
\pagestyle{fancy}
\fancyhf{}
\renewcommand{\headrulewidth}{0pt}
\renewcommand{\footrulewidth}{0pt}
\lhead{\raisebox{-0.15ex}{\poplogo\fontsize{13}{13}\selectfont\color{popink}OnBuch\color{poporange}+}}
\rhead{\poppill{\DOCMATIERE\ \textbullet\ \DOCNIVEAU\ \textbullet\ Leçon \DOCLECON}{white}{popink}}
\lfoot{\popxbold\fontsize{7.6}{7.6}\selectfont\color{popmuted}\MakeUppercase{Cours signé OnBuch+}\ \textbullet\ {\popsemi\DOCTITRE}}
\rfoot{\tikz[baseline=-0.6ex]\node[rounded corners=8.5pt,fill=popink,minimum height=17pt,minimum width=17pt,inner xsep=5pt,inner sep=0]{\popxbold\fontsize{8}{8}\selectfont\color{popcream}\thepage\,/\,\pageref*{LastPage}};}

% ============================================================
% Titres
% ============================================================
\newcommand{\chaptitre}[2]{%
  \begin{tcolorbox}[enhanced,colback=poporange,colframe=popink,boxrule=2.6pt,arc=11pt,
      drop shadow=popink,left=15pt,right=15pt,top=12pt,bottom=11pt,before skip=3pt,after skip=18pt]
    \poppill{#2}{popink}{popcream}\par\vspace{9pt}
    {\raggedright\hyphenpenalty=10000\exhyphenpenalty=10000\popdisplay\fontsize{22}{24}\selectfont\color{popcream}\MakeUppercase{#1}\par}
  \end{tcolorbox}
}
% Section numérotée : pastille orange avec le numéro + titre Archivo
\newcounter{popsec}
\newcommand{\coursec}[1]{%
  \stepcounter{popsec}\setcounter{popsub}{0}%
  \par\vspace{16pt}\noindent
  \begin{minipage}{\linewidth}
  \tikz[baseline=-0.7ex]\node[draw=popink,line width=1.6pt,fill=poporange,rounded corners=4pt,minimum size=22pt,inner sep=0]{\popdisplay\fontsize{12}{12}\selectfont\color{popcream}\thepopsec};%
  \hspace{8pt}{\popdisplay\fontsize{15}{17}\selectfont\color{popink}\MakeUppercase{#1}}\par
  \vspace{4pt}\noindent{\color{poporange}\rule{36pt}{3.6pt}}
  \end{minipage}\par\nopagebreak\vspace{8pt}
}
\newcounter{popsub}
\newcommand{\courssub}[1]{%
  \stepcounter{popsub}%
  \par\vspace{9pt}\noindent{\popxbold\fontsize{11.5}{13}\selectfont\color{popink}{\color{poporange}\thepopsec.\thepopsub}\hspace{6pt}#1}\par\nopagebreak\vspace{4pt}
}

% ============================================================
% Boîtes pédagogiques — même recette : fond blanc, bordure épaisse colorée,
% ombre dure encre, badge plein. Titre optionnel [#1].
% ============================================================
\tcbset{popbase/.style={breakable,enhanced,colback=white,boxrule=2.3pt,arc=9pt,
  shadow={3pt}{-3pt}{0pt}{popink},left=14pt,right=14pt,top=11pt,bottom=11pt,before skip=11pt,after skip=15pt,
  fontupper=\popsemi\fontsize{10.6}{14.5}\selectfont\color{popink},
  fontlower=\popsemi\fontsize{10.6}{14.5}\selectfont\color{popink}}}
% Coche dessinée (le glyphe ✓ n'existe pas dans les polices du document)
\newcommand{\popcheck}[1]{\tikz[baseline=-0.5ex]\draw[#1,line width=1.6pt,line cap=round,line join=round](-0.13,0.0)--(-0.04,-0.09)--(0.13,0.1);}
\newcommand{\popboxhead}[3]{\poppill{#1}{#2}{white}\ifx\relax#3\relax\else\hspace{6pt}{\popxbold\fontsize{10.6}{12}\selectfont\color{popink}#3}\fi\par\vspace{7pt}}

\newtcolorbox{aretenir}[1][]{popbase,colframe=popgreen,before upper={\popboxhead{À retenir}{popgreen}{#1}}}
\newtcolorbox{exemplebox}[1][]{popbase,colframe=poporange,before upper={\popboxhead{Exemple}{poporange}{#1}}}
\newtcolorbox{definition}[1][]{popbase,colframe=popblue,before upper={\popboxhead{Définition}{popblue}{#1}}}
\newtcolorbox{propriete}[1][]{popbase,colframe=poppurple,before upper={\popboxhead{Propriété}{poppurple}{#1}}}
\newtcolorbox{methode}[1][]{popbase,colframe=popgold,before upper={\popboxhead{Méthode}{popgold}{#1}}}
\newtcolorbox{attention}[1][]{popbase,colframe=poppink,before upper={\popboxhead{Attention}{poppink}{#1}}}
\newtcolorbox{experience}[1][]{popbase,colframe=popblue,colback=popblueL,before upper={\popboxhead{Expérience}{popblue}{#1}}}
% « Le savais-tu ? » : carte violette pleine (contexte, culture, Cameroun)
\newtcolorbox{savaistu}[1][]{popbase,colback=poppurple,colframe=popink,
  fontupper=\popsemi\fontsize{10.4}{14.2}\selectfont\color{white},
  before upper={\poppill{Le savais-tu\,?}{white}{poppurple}\ifx\relax#1\relax\else\hspace{6pt}{\popxbold\color{white}#1}\fi\par\vspace{7pt}}}
% Exercice résolu : énoncé en haut, \tcblower puis la solution sur fond crème
\newcounter{popexo}\newcounter{popexr}
\newtcolorbox{exoresolu}[1][]{popbase,colframe=poporange,colbacklower=popcream,
  segmentation style={popink,line width=1.2pt,dashed},
  before upper={\stepcounter{popexr}\popboxhead{Exercice résolu \thepopexr}{poporange}{#1}},
  before lower={\poppill{Solution}{popink}{popcream}\par\vspace{6pt}}}
% Exercice (non résolu en place) — la correction est dans la partie Corrigés
\newtcolorbox{exercice}[1][]{popbase,colframe=popink,
  before upper={\stepcounter{popexo}\popboxhead{Exercice \thepopexo}{popink}{#1}}}
% Corrigé
\newtcolorbox{corrige}[1][]{popbase,colframe=popgreen,colback=popgreenL,
  before upper={\popboxhead{Corrigé}{popgreen}{#1}}}
% Objectifs / prérequis / bilan
\newtcolorbox{objectifs}{popbase,colframe=popink,colback=poporangeL,
  before upper={\popboxhead{Objectifs de la leçon}{popink}{}}}
\newtcolorbox{prerequis}{popbase,colframe=popink,colback=popblueL,
  before upper={\popboxhead{Prérequis}{popblue}{}}}
\newtcolorbox{bilan}{popbase,colframe=popink,colback=white,boxrule=2.8pt,
  before upper={\popboxhead{Fiche bilan}{poporange}{L'essentiel à connaître pour l'examen}}}
% Figure encadrée avec légende
\newtcolorbox{popfigure}[1][]{enhanced,breakable=false,colback=white,colframe=popline,boxrule=1.4pt,arc=8pt,
  left=10pt,right=10pt,top=10pt,bottom=8pt,before skip=10pt,after skip=12pt,halign=center,
  lower separated=false,halign lower=center,
  fontlower=\popsemi\fontsize{9}{11.5}\selectfont\color{popmuted},
  #1}
\newcommand{\legende}[1]{\par\vspace{6pt}{\popsemi\fontsize{9}{11.5}\selectfont\color{popmuted}#1}}

% ============================================================
% Signature OnBuch+
% ============================================================
\newcommand{\poplogoinline}[1]{{\poplogo\fontsize{#1}{#1}\selectfont\color{popink}OnBuch\color{poporange}+}}
\newcommand{\signatureonbuch}{%
  \begin{tcolorbox}[enhanced,colback=popink,colframe=popink,boxrule=2.2pt,arc=10pt,
      drop shadow=poporange,left=14pt,right=14pt,top=10pt,bottom=10pt,before skip=0pt,after skip=0pt]
    \tikz[baseline=-0.7ex]\node[circle,fill=popgreen,draw=popcream,line width=1.3pt,minimum size=22pt,inner sep=0]{\popcheck{white}};%
    \hspace{9pt}\begin{minipage}[c]{0.62\linewidth}
      {\popxbold\fontsize{12}{13}\selectfont\color{popcream}Signé\ \textbullet\ {\poplogo OnBuch\color{poporange}+}}\par\vspace{2pt}
      {\raggedright\popsemi\fontsize{8}{10.5}\selectfont\color{popline}\MakeUppercase{Équipe pédagogique OnBuch+}\\\MakeUppercase{Conforme au programme officiel MINESEC}\par}
    \end{minipage}\hfill
    {\poplogo\fontsize{22}{22}\selectfont\color{popcream}OnBuch\color{poporange}+}
  \end{tcolorbox}%
}

% ============================================================
% Page de garde
% #1 titre · #2 matière · #3 niveau · #4 module · #5 n° leçon · #6 durée ·
% #7 description / objectifs (texte ou itemize)
% ============================================================
\newcommand{\pagegarde}[7]{%
  \thispagestyle{empty}
  \newgeometry{a4paper,margin=1.7cm,top=1.6cm,bottom=1.4cm}%
  \begin{tikzpicture}[remember picture,overlay]
    \fill[poporange!18] ([xshift=-3.2cm,yshift=-3.6cm]current page.north east) circle (4.6cm);
    \fill[poppurple!14] ([xshift=2.4cm,yshift=7.6cm]current page.south west) circle (3.8cm);
  \end{tikzpicture}%
  {\poplogo\fontsize{30}{30}\selectfont\color{popink}OnBuch\color{poporange}+}\hfill
  \raisebox{0.6ex}{\poppill{Cours signé \textbullet\ Premium}{popink}{popcream}}\par
  \vspace{1pt}\popsquiggle{poppurple}\par
  \vspace{8pt}
  \begin{tcolorbox}[enhanced,colback=poporange,colframe=popink,boxrule=2.8pt,arc=13pt,
      drop shadow=popink,left=18pt,right=18pt,top=12pt,bottom=13pt,after skip=10pt]
    \poppill{#2 \textbullet\ #3}{popink}{popcream}\hspace{5pt}\poppill{#4}{white}{popink}\par\vspace{8pt}
    {\popdisplay\fontsize{14}{14}\selectfont\color{popink}LEÇON #5}\par\vspace{4pt}
    {\raggedright\hyphenpenalty=10000\exhyphenpenalty=10000\popdisplay\fontsize{25}{27}\selectfont\color{popcream}\MakeUppercase{#1}\par}
    \vspace{8pt}
    {\popxbold\fontsize{9.5}{9.5}\selectfont\color{popink}\MakeUppercase{Durée conseillée : #6}}
  \end{tcolorbox}
  \begin{tcolorbox}[enhanced,colback=white,colframe=popink,boxrule=2.2pt,arc=10pt,
      drop shadow=popink,left=16pt,right=16pt,top=10pt,bottom=10pt,after skip=8pt]
    \poppill{Dans cette leçon}{poppurple}{white}\par\vspace{5pt}
    \popsemi\fontsize{9.3}{12.6}\selectfont\color{popink}#7
  \end{tcolorbox}
  \noindent\begin{minipage}[t]{0.487\linewidth}
    \begin{tcolorbox}[enhanced,colback=popgreen,colframe=popink,boxrule=2.2pt,arc=10pt,
        drop shadow=popink,left=11pt,right=11pt,top=10pt,bottom=10pt,height=3.1cm,valign=center]
      \tcbox[colback=white,colframe=popink,boxrule=1.3pt,arc=5pt,boxsep=0pt,nobeforeafter,
        left=3pt,right=3pt,top=3pt,bottom=3pt,valign=center]{\qrcode[height=1.9cm]{https://play.google.com/store/apps/details?id=com.luvvix.onbuch&hl=fr}}%
      \hspace{8pt}\begin{minipage}[c]{0.5\linewidth}
        {\popdisplay\fontsize{11}{12.5}\selectfont\color{white}\MakeUppercase{Télécharge OnBuch}}\par\vspace{4pt}
        {\raggedright\popsemi\fontsize{8.4}{11}\selectfont\color{white}Gratuit \textbullet\ Google Play\\Tuteur IA, annales, concours, résultats\par}
      \end{minipage}
    \end{tcolorbox}
  \end{minipage}\hfill
  \begin{minipage}[t]{0.487\linewidth}
    \begin{tcolorbox}[enhanced,colback=poppurple,colframe=popink,boxrule=2.2pt,arc=10pt,
        drop shadow=popink,left=11pt,right=11pt,top=10pt,bottom=10pt,height=3.1cm,valign=center]
      \tikz[baseline=-0.7ex]\node[circle,fill=white,draw=popink,line width=1.3pt,minimum size=1.6cm,inner sep=0]{\popdisplay\fontsize{24}{24}\selectfont\color{poppurple}+};%
      \hspace{8pt}\begin{minipage}[c]{0.6\linewidth}
        {\popdisplay\fontsize{11}{12.5}\selectfont\color{white}\MakeUppercase{Passe à OnBuch+}}\par\vspace{4pt}
        {\raggedright\popsemi\fontsize{8.4}{11}\selectfont\color{white}Tous les cours signés, Tuteur IA illimité, fiches PDF — onglet OnBuch+ de l'app\par}
      \end{minipage}
    \end{tcolorbox}
  \end{minipage}
  \vspace{8pt}
  \popcontenu
  \vfill
  \signatureonbuch
  \restoregeometry
  \newpage
}

% Rangée « Ce cours contient » (page de garde)
\newcommand{\poptile}[3]{% #1 couleur #2 gros texte #3 légende
  \begin{tcolorbox}[enhanced,colback=#1,colframe=popink,boxrule=2pt,arc=9pt,shadow={2.5pt}{-2.5pt}{0pt}{popink},
      left=6pt,right=6pt,top=9pt,bottom=9pt,halign=center,valign=center,height=1.75cm,width=\linewidth,nobeforeafter]
    {\popdisplay\fontsize{12.5}{13}\selectfont\color{white}\MakeUppercase{#2}}\par\vspace{3pt}
    {\centering\popsemi\fontsize{7.8}{9.5}\selectfont\color{white}#3\par}
  \end{tcolorbox}}
\newcommand{\popcontenu}{%
  \par\noindent{\popxboldls\fontsize{7.5}{7.5}\selectfont\color{popmuted}CE COURS SIGNÉ CONTIENT}\par\vspace{7pt}
  \noindent\begin{minipage}[t]{0.235\linewidth}\poptile{popblue}{Cours}{complet, illustré, pas à pas}\end{minipage}\hfill
  \begin{minipage}[t]{0.235\linewidth}\poptile{poporange}{Méthodes}{et exercices résolus}\end{minipage}\hfill
  \begin{minipage}[t]{0.235\linewidth}\poptile{poppink}{Exercices}{type examen + corrigés}\end{minipage}\hfill
  \begin{minipage}[t]{0.235\linewidth}\poptile{popgreen}{Fiche bilan}{l'essentiel à retenir}\end{minipage}\par}

% Sommaire
\newtcolorbox{sommaire}{enhanced,colback=white,colframe=popink,boxrule=2.2pt,arc=10pt,
  drop shadow=popink,left=18pt,right=18pt,top=13pt,bottom=13pt,before skip=6pt,after skip=20pt,
  before upper={\poppill{Sommaire}{poppurple}{white}\par\vspace{9pt}\popsemi\fontsize{10.6}{14.5}\selectfont\color{popink}}}

% Signature de fin de cours — poussée tout en bas de la dernière page
\newcommand{\findecours}{%
  \par\vfill\nopagebreak
  \begin{center}\popsquiggle{poporange}\end{center}\vspace{4pt}
  \signatureonbuch
}
\AtBeginDocument{\def\llbracket{\mathopen{[\mkern-3.5mu[}}\def\rrbracket{\mathclose{]\mkern-3.5mu]}}} % intervalles d'entiers (glyphe ⟦ absent de la police)
% Glyphes absents de Fira Math : redéfinis à partir de glyphes disponibles
\AtBeginDocument{%
  \def\setminus{\mathbin{\backslash}}\let\smallsetminus\setminus
  \def\vdots{\mathord{\vbox{\baselineskip3.2pt\lineskiplimit0pt\kern4pt\hbox{.}\hbox{.}\hbox{.}}}}%
  \def\ddots{\mathinner{\mkern1mu\raise6.5pt\hbox{.}\mkern2mu\raise3.6pt\hbox{.}\mkern2mu\raise0.7pt\hbox{.}\mkern1mu}}%
  \def\triangle{\mathord{\scalebox{0.9}{$\symup{\Delta}$}}}%
  \def\nabla{\mathord{\raisebox{\depth}{\scalebox{1}[-1]{$\symup{\Delta}$}}}}%
  \def\checkmark{\popcheck{popdark}}%
  \def\star{\mathbin{\ast}}\def\bigstar{\mathord{\ast}}%
  \def\diamond{\mathbin{\scalebox{0.6}{$\Diamond$}}}%
  \def\bigcup{\mathop{\mathchoice{\raisebox{-0.3em}{\scalebox{1.7}{$\cup$}}}{\scalebox{1.25}{$\cup$}}{\cup}{\cup}}\displaylimits}%
  \def\bigcap{\mathop{\mathchoice{\raisebox{-0.3em}{\scalebox{1.7}{$\cap$}}}{\scalebox{1.25}{$\cap$}}{\cap}{\cap}}\displaylimits}%
  \def\lesseqqgtr{\mathrel{\lessgtr}}\def\bigoplus{\mathop{\oplus}\displaylimits}%
}
\providecommand{\ssi}{si et seulement si} % abréviation fréquente
\AtBeginDocument{\providecommand{\wideparen}{\overparen}} % arc de cercle

%% FIGFIX-BEGIN v5 — correctifs globaux des figures (docs/onbuch-generation/tools/figfix)
\usepackage{adjustbox}\usepackage{etoolbox}\tcbuselibrary{raster}
% toute figure TikZ/pgfplots plus large que la ligne est réduite à la largeur disponible
\BeforeBeginEnvironment{tikzpicture}{\begin{adjustbox}{max width=\linewidth}}
\AfterEndEnvironment{tikzpicture}{\end{adjustbox}}
% légendes pgfplots lisibles par-dessus les courbes
\pgfplotsset{every axis/.append style={legend style={fill=white,fill opacity=0.92,text opacity=1,draw=black!15,inner sep=3pt}}}
% étiquettes d'arêtes (syntaxe edge["..."]) avec fond blanc
\tikzset{every edge quotes/.append style={fill=white,inner sep=1.2pt}}
%% FIGFIX-END
\begin{document}

\pagegarde{Probabilités}{Mathématiques}{Tle Tech}{Mathématiques – classe de Terminale technique}{N05}{11 séances (fiche de progression harmonisée)}{Cette leçon aborde les notions fondamentales de probabilités (vocabulaire, expériences aléatoires, probabilités conditionnelles, variables aléatoires, épreuves de Bernoulli) et les coniques (définition, équations réduites, tangentes, construction) pour outiller l'élève face à des problèmes concrets de modélisation et de décision.
\vspace{4pt}
\begin{multicols}{2}\small
\begin{itemize}
\item À la fin de cette leçon, je sais définir une expérience aléatoire, un événement et calculer sa probabilité.
\item Je sais utiliser la formule des probabilités conditionnelles et l'appliquer à des situations réelles.
\item Je sais modéliser une suite d'épreuves indépendantes par une variable aléatoire suivant une loi de Bernoulli ou binomiale.
\item Je sais reconnaître une conique (ellipse, parabole, hyperbole) à partir de son équation réduite.
\item Je sais déterminer l'équation de la tangente en un point d'une conique.
\item Je sais construire une conique à l'aide de sa définition géométrique (foyers, directrices).
\end{itemize}\end{multicols}}

\begin{sommaire}
\begin{multicols}{2}
\begin{enumerate}
\item 1. Vocabulaire et expériences aléatoires
\item 2. Probabilité d'un événement et probabilités conditionnelles
\item 3. Variables aléatoires et épreuves de Bernoulli
\item 4. Présentation générale des coniques et propriétés
\item 5. Tangente en un point à une conique
\item 6. Construction d'une conique et applications intégrées
\item Activité d'intégration
\item Méthodes et exercices résolus
\item Exercices
\item Corrigés des exercices
\item Fiche bilan
\end{enumerate}
\end{multicols}
\end{sommaire}

\begin{objectifs}
\begin{itemize}
\item À la fin de cette leçon, je sais définir une expérience aléatoire, un événement et calculer sa probabilité.
\item Je sais utiliser la formule des probabilités conditionnelles et l'appliquer à des situations réelles.
\item Je sais modéliser une suite d'épreuves indépendantes par une variable aléatoire suivant une loi de Bernoulli ou binomiale.
\item Je sais reconnaître une conique (ellipse, parabole, hyperbole) à partir de son équation réduite.
\item Je sais déterminer l'équation de la tangente en un point d'une conique.
\item Je sais construire une conique à l'aide de sa définition géométrique (foyers, directrices).
\item Je sais rédiger une démonstration guidée pour les propriétés principales des coniques.
\item Je sais mobiliser ces outils pour résoudre un problème intégrateur mêlant probabilités et géométrie analytique.
\end{itemize}
\end{objectifs}

\begin{prerequis}
\begin{itemize}
\item Notions d'ensembles, d'union, d'intersection et de complémentaire (4ème).
\item Calcul des pourcentages et manipulation des fractions (4ème).
\item Repère cartésien, équation de la droite et du cercle (3ème/2nde).
\item Notion de fonction, tableau de variation et dérivée (1ère).
\item Résolution d'équations du second degré et inéquations (2nde).
\end{itemize}
\end{prerequis}

\newpage

\coursec{Vocabulaire et expériences aléatoires}

Au quartier Mokolo, un commerçant lance un jeu : il place une cible circulaire sur un mur et propose aux passants de lancer une balle en mousse. Chaque lancer coûte 100 FCFA et le gain est de 500 FCFA si la balle touche la cible. Combien de fois faut-il jouer pour espérer gagner ? Cette situation mêle hasard, géométrie de la cible et modélisation mathématique. Pour y répondre, nous devons d'abord poser le langage rigoureux des probabilités : qu'est-ce qu'une expérience aléatoire ? Comment décrire l'ensemble de tous les résultats possibles ? Comment manipuler les événements ?

\courssub{Expérience aléatoire et univers des possibles}

\begin{definition}[Expérience aléatoire et univers $\Omega$]
Une \cle{expérience aléatoire} est une expérience qui remplit deux conditions :
\begin{enumerate}
    \item On peut en lister \emph{à l'avance} tous les résultats possibles.
    \item On ne peut \emph{pas prédire avec certitude} lequel se produira lors d'une réalisation.
\end{enumerate}
L'ensemble de tous les résultats possibles est appelé \cle{univers} (ou espace des possibles) et se note \cle{$\Omega$} (oméga majuscule). Un résultat élémentaire se note souvent $\omega$ (oméga minuscule).
\end{definition}

\begin{exemplebox}[Exemples d'expériences aléatoires]
\begin{itemize}
    \item \textbf{Lancer un dé à 6 faces équilibrées.} $\Omega = \{1; 2; 3; 4; 5; 6\}$. On connaît la liste, mais le résultat du prochain lancer est incertain.
    \item \textbf{Tirer une boule dans une urne contenant 3 boules blanches (B) et 2 boules noires (N).} $\Omega = \{B_1; B_2; B_3; N_1; N_2\}$ (si on numérote les boules pour les distinguer) ou simplement $\Omega = \{B; N\}$ si on ne s'intéresse qu'à la couleur.
    \item \textbf{Contrôler la tension d'une pile neuve.} $\Omega = [0; +\infty[$ (intervalle de valeurs réelles). C'est une expérience à résultats \emph{continus}.
\end{itemize}
\end{exemplebox}

\begin{attention}[Expérience déterministe vs aléatoire]
Mesurer la longueur d'une pièce avec une règle graduée n'est \textbf{pas} une expérience aléatoire : le résultat est (presque) certain si l'opérateur est compétent. En revanche, \emph{mesurer le temps de réaction d'un opérateur} face à un signal lumineux est une expérience aléatoire (variabilité physiologique).
\end{attention}

\courssub{Les événements : définition et opérations}

\begin{definition}[Événement]
Un \cle{événement} est une partie de l'univers $\Omega$, c'est-à-dire un ensemble de résultats possibles. On note généralement les événements par des lettres majuscules : $A, B, C, \dots$
\begin{itemize}
    \item L'événement \cle{certain} est $\Omega$ lui-même (il se réalise à coup sûr).
    \item L'événement \cle{impossible} est l'ensemble vide $\varnothing$ (il ne peut jamais se réaliser).
    \item L'événement \cle{contraire} (ou complémentaire) de $A$ est $\overline{A} = \Omega \setminus A$ : « $A$ ne se réalise pas ».
\end{itemize}
\end{definition}

\begin{propriete}[Opérations sur les événements]
Soient $A$ et $B$ deux événements de $\Omega$.
\begin{itemize}
    \item \textbf{Union} $A \cup B$ : « $A$ \textbf{ou} $B$ se réalise » (au moins l'un des deux).
    \item \textbf{Intersection} $A \cap B$ : « $A$ \textbf{et} $B$ se réalisent » (les deux simultanément).
    \item \textbf{Différence} $A \setminus B$ : « $A$ se réalise \textbf{et} $B$ ne se réalise pas » ($A \cap \overline{B}$).
    \item \textbf{Incompatibilité} : $A$ et $B$ sont \cle{incompatibles} (ou mutuellement exclusifs) si $A \cap B = \varnothing$. Ils ne peuvent pas se réaliser ensemble.
\end{itemize}
\end{propriete}

\begin{popfigure}
\begin{center}
\begin{tikzpicture}[scale=1.2, every node/.style={font=\small}]
    % Univers rectangle
    \draw[popdark, thick] (-3,-2) rectangle (3,2) node[above left, popdark] {$\Omega$};
    % Cercle A
    \draw[popblue, thick, fill=popblueL] (-1,0) circle (1.2) node[left=0.2cm] {$A$};
    % Cercle B
    \draw[poporange, thick, fill=poporangeL] (1,0) circle (1.2) node[right=0.2cm] {$B$};
    % Intersection hachurée
    \begin{scope}
        \clip (-1,0) circle (1.2);
        \clip (1,0) circle (1.2);
        \fill[poppurpleL, opacity=0.6] (-1,0) circle (1.2);
    \end{scope}
    % Légendes
    \node[poppurple] at (0,0) {$A \cap B$};
    \node[popblue] at (-2,1.2) {$A \setminus B$};
    \node[poporange] at (2,1.2) {$B \setminus A$};
    \node[popmuted] at (0,-1.5) {$\overline{A \cup B}$};
    % Flèches explicatives
    \draw[->, popdark] (-3.5, 1.5) -- (-1.5, 1.5) node[midway, above] {$A \cup B$};
    \draw[->, popdark] (3.5, -1.5) -- (1.5, -1.5) node[midway, below] {$\overline{A} \cap \overline{B}$};
\end{tikzpicture}
\end{center}
\legende{Diagramme de Venn : union, intersection, différence et complémentaire.}
\end{popfigure}

\begin{attention}[Incompatibles $\neq$ Contraires]
Deux événements \textbf{incompatibles} ne peuvent pas arriver ensemble ($A \cap B = \varnothing$).
Deux événements \textbf{contraires} sont incompatibles \textbf{ET} l'un des deux arrive obligatoirement ($A \cup B = \Omega$).
\begin{itemize}
    \item Au lancer d'un dé : $A=$« résultat pair » et $B=$« résultat impair » sont \textbf{contraires}.
    \item Au lancer d'un dé : $A=$« résultat = 1 » et $B=$« résultat = 2 » sont \textbf{incompatibles} mais pas contraires (on peut faire 3, 4, 5 ou 6).
\end{itemize}
\end{attention}

\courssub{Arbre des issues pour une expérience à plusieurs étapes}

Lorsque l'expérience se déroule en étapes successives (lancer deux dés, tirer deux boules sans remise, etc.), l'\cle{arbre des issues} (ou arbre de probabilités) est l'outil de représentation idéal.

\begin{methode}[Construire l'arbre des issues d'une expérience à deux étapes]
\begin{enumerate}
    \item \textbf{Identifier les étapes} (Étape 1, Étape 2, ...).
    \item \textbf{Dessiner le premier niveau} : depuis une racine, tracer une branche par issue possible de l'étape 1. Noter l'issue sur la branche.
    \item \textbf{Dessiner les niveaux suivants} : au bout de chaque branche de l'étape 1, tracer les branches correspondant aux issues possibles de l'étape 2 (attention : les issues peuvent dépendre de l'étape 1, ex: tirage sans remise).
    \item \textbf{Calculer les probabilités} : on écrit la probabilité sur chaque branche. Le long d'un chemin, on \textbf{multiplie} les probabilités (règle du produit). La somme des probabilités des branches issues d'un même nœud vaut 1.
    \item \textbf{Lire les résultats} : chaque chemin complet (de la racine à une feuille) correspond à un résultat élémentaire de l'expérience composée. Sa probabilité est le produit des probabilités sur le chemin.
\end{enumerate}
\end{methode}

\begin{popfigure}
\begin{center}
\begin{tikzpicture}[
    scale=0.9,
    level 1/.style={sibling distance=3.5cm, level distance=2.5cm},
    level 2/.style={sibling distance=0.7cm, level distance=2.5cm},
    edge from parent/.style={draw=popdark, thick, ->},
    every node/.style={font=\small, inner sep=2pt},
    branch/.style={edge from parent path={(\tikzparentnode) -- (\tikzchildnode)}},
    leaf/.style={fill=popcream, draw=popline, rounded corners, minimum width=1.2cm, minimum height=0.6cm},
    node1/.style={circle, draw=popdark, fill=popblueL, minimum size=6pt, inner sep=0}
]
\node[node1] (root) {}
    child[branch] { node[leaf] (1) {1}
        child[branch] { node[leaf] (11) {(1,1)} edge from parent node[left] {$\frac{1}{6}$} }
        child[branch] { node[leaf] (12) {(1,2)} edge from parent node[left] {$\frac{1}{6}$} }
        child[branch] { node[leaf] (13) {(1,3)} edge from parent node[left] {$\frac{1}{6}$} }
        child[branch] { node[leaf] (14) {(1,4)} edge from parent node[left] {$\frac{1}{6}$} }
        child[branch] { node[leaf] (15) {(1,5)} edge from parent node[left] {$\frac{1}{6}$} }
        child[branch] { node[leaf] (16) {(1,6)} edge from parent node[left] {$\frac{1}{6}$} }
        edge from parent node[left, pos=0.5] {$\frac{1}{6}$}
    }
    child[branch] { node[leaf] (2) {2}
        child[branch] { node[leaf] (21) {(2,1)} edge from parent node[left] {$\frac{1}{6}$} }
        child[branch] { node[leaf] (22) {$\cdots$} edge from parent node[left] {$\frac{1}{6}$} }
        edge from parent node[left, pos=0.5] {$\frac{1}{6}$}
    }
    child[branch] { node[leaf] (3) {3} edge from parent node[left, pos=0.5] {$\frac{1}{6}$} }
    child[branch] { node[leaf] (4) {4} edge from parent node[right, pos=0.5] {$\frac{1}{6}$} }
    child[branch] { node[leaf] (5) {5} edge from parent node[right, pos=0.5] {$\frac{1}{6}$} }
    child[branch] { node[leaf] (6) {6}
        child[branch] { node[leaf] (61) {$\cdots$} edge from parent node[right] {$\frac{1}{6}$} }
        child[branch] { node[leaf] (66) {(6,6)} edge from parent node[right] {$\frac{1}{6}$} }
        edge from parent node[right, pos=0.5] {$\frac{1}{6}$}
    };
\node[above=0.3cm of root, popdark, font=\bfseries] {Lancer 1};
\node[below=2.5cm of 1, popdark, font=\bfseries] {Lancer 2 (si 1 au 1er)};
\node[below=2.5cm of 2, popmuted, font=\small] {Autres branches similaires};
\end{tikzpicture}
\end{center}
\legende{Arbre des issues (partiel) pour deux lancers successifs d'un dé équilibré (36 issues équiprobables au total).}
\end{popfigure}

\courssub{Probabilité d'un événement : définition et premières propriétés}

Dans le cadre de l'enseignement technique (probabilités finies, issues équiprobables), on adopte la définition classique de Laplace.

\begin{definition}[Probabilité d'un événement (cas équiprobable)]
Si l'univers $\Omega$ est \textbf{fini} et que tous les résultats élémentaires ont la même chance de se produire (on dit qu'ils sont \cle{équiprobables}), alors pour tout événement $A \subset \Omega$ :
\[
P(A) = \frac{\text{Nombre de résultats favorables à } A}{\text{Nombre total de résultats possibles}} = \frac{\text{Card}(A)}{\text{Card}(\Omega)}.
\]
La probabilité est un nombre réel compris entre 0 et 1.
\end{definition}

\begin{propriete}[Propriétés fondamentales de la probabilité]
Pour tout événement $A$ de $\Omega$ :
\begin{enumerate}
    \item $0 \le P(A) \le 1$.
    \item $P(\Omega) = 1$ \quad (certitude).
    \item $P(\varnothing) = 0$ \quad (impossibilité).
    \item $P(\overline{A}) = 1 - P(A)$.
    \item Si $A \subset B$ alors $P(A) \le P(B)$.
    \item \textbf{Formule de l'union :} $P(A \cup B) = P(A) + P(B) - P(A \cap B)$.
    \item Si $A$ et $B$ sont incompatibles ($A \cap B = \varnothing$) : $P(A \cup B) = P(A) + P(B)$.
\end{enumerate}
\end{propriete}

\begin{exoresolu}[Calculer $P(\text{« somme = 7 »})$ au lancer de deux dés équilibrés]
\textbf{Énoncé :} On lance deux dés à 6 faces équilibrés (ou un même dé deux fois). On note $S$ la somme des deux résultats. Calculer $P(S = 7)$.

\textbf{Solution détaillée :}
\begin{enumerate}
    \item \textbf{Univers $\Omega$ :} Les couples ordonnés $(d_1; d_2)$ avec $d_1, d_2 \in \{1,\dots,6\}$. $\text{Card}(\Omega) = 6 \times 6 = 36$. Les 36 issues sont équiprobables (arbre ci-dessus).
    \item \textbf{Événement $A$ :} « La somme vaut 7 ». $A = \{(1,6); (2,5); (3,4); (4,3); (5,2); (6,1)\}$.
    \item \textbf{Cardinal :} $\text{Card}(A) = 6$.
    \item \textbf{Probabilité :} $P(A) = \dfrac{6}{36} = \dfrac{1}{6} \approx 0,1667$.
\end{enumerate}
\emph{Remarque :} C'est la somme la plus probable au lancer de deux dés.
\end{exoresolu}

\begin{exemplebox}[Tirage dans une urne : avec et sans remise]
Une urne contient 5 boules : 3 Blanches (B), 2 Rouges (R). On tire 2 boules successivement.
\begin{enumerate}
    \item \textbf{Avec remise :} Après le 1er tirage, on remet la boule. Les tirages sont indépendants. $\Omega$ a $5 \times 5 = 25$ issues équiprobables.
        $P(\text{2 Blanches}) = \frac{3}{5} \times \frac{3}{5} = \frac{9}{25} = 0,36$.
    \item \textbf{Sans remise :} On ne remet pas la 1ère boule. Les tirages ne sont pas indépendants. $\Omega$ a $5 \times 4 = 20$ issues équiprobables (couples ordonnés de boules distinctes).
        $P(\text{2 Blanches}) = \frac{3}{5} \times \frac{2}{4} = \frac{6}{20} = 0,30$.
\end{enumerate}
L'arbre des issues permet de visualiser la modification des probabilités au 2ème niveau dans le cas « sans remise ».
\end{exemplebox}

\courssub{Application : probabilité géométrique et jeu de Mokolo}

\begin{definition}[Probabilité géométrique]
Si l'univers $\Omega$ est une région du plan (ou de l'espace) de mesure finie et non nulle (aire, longueur, volume), et si le hasard est \emph{uniforme} (tout point a la même chance d'être atteint), alors pour tout événement $A$ correspondant à une sous-région mesurable :
\[
P(A) = \frac{\text{Mesure}(A)}{\text{Mesure}(\Omega)}.
\]
\end{definition}

Reprenons la situation d'accroche. La cible est un disque de rayon $r = 10$ cm placé au centre d'un panneau carré de côté $L = 1$ m = 100 cm. Le joueur lance la balle au hasard sur le panneau (hypothèse : le point d'impact est uniformément réparti sur le carré).

\begin{exoresolu}[Probabilité de gagner au jeu de Mokolo]
\textbf{Énoncé :} Calculer la probabilité de toucher la cible. En déduire l'espérance de gain pour une partie.

\textbf{Solution :}
\begin{enumerate}
    \item \textbf{Univers $\Omega$ :} L'ensemble des points du panneau carré. $\text{Aire}(\Omega) = L^2 = 10\,000 \text{ cm}^2$.
    \item \textbf{Événement $G$ :} « La balle touche la cible » $\iff$ le point d'impact appartient au disque de rayon $r=10$ cm.
    \item \textbf{Aire favorable :} $\text{Aire}(G) = \pi r^2 = \pi \times 100 \approx 314,16 \text{ cm}^2$.
    \item \textbf{Probabilité (géométrique) :} $P(G) = \frac{\text{Aire}(G)}{\text{Aire}(\Omega)} = \frac{100\pi}{10\,000} = \frac{\pi}{100} \approx 0,0314$.
    \item \textbf{Espérance de gain (aperçu) :} Gain = 500 FCFA si succès, 0 sinon. Mise = 100 FCFA.
        Gain net aléatoire $X$ : $X = 400$ avec $p \approx 0,0314$, $X = -100$ avec $p \approx 0,9686$.
        Espérance $E(X) = 400 \times 0,0314 + (-100) \times 0,9686 \approx 12,56 - 96,86 = -84,30$ FCFA.
        \emph{(Cette notion d'espérance mathématique sera formalisée dans la section sur les variables aléatoires).}
\end{enumerate}
\textbf{Conclusion :} Le jeu est très défavorable au joueur (perte moyenne de 84 FCFA par partie). Il faudrait jouer environ $1/0,0314 \approx 32$ fois pour toucher la cible une fois en moyenne, ce qui coûterait 3200 FCFA pour gagner 500 FCFA.
\end{exoresolu}

\begin{savaistu}[La loi des grands nombres : la fréquence stabilise la probabilité]
Si vous répétez un grand nombre de fois la même expérience aléatoire dans les mêmes conditions, la \cle{fréquence} d'apparition d'un événement (nombre de succès / nombre total d'essais) se rapproche de plus en plus de sa \cle{probabilité théorique}.
C'est ce qui justifie l'utilisation des statistiques passées (fréquences observées) pour estimer des probabilités inconnues (ex: taux de panne d'une machine, probabilité de pluie, fiabilité d'un composant électronique). Au Probatoire, on distingue bien la probabilité \emph{théorique} (calculée sur un modèle) et la fréquence \emph{observée} (issue d'une simulation ou d'un relevé réel).
\end{savaistu}

\begin{exercice}[Entraînement : Vocabulaire et arbre]
\begin{enumerate}
    \item Une entreprise fabrique des pièces. 95\% sont conformes (C), 5\% sont défectueuses (D). On prélève deux pièces au hasard (avec remise, car le stock est grand).
    \begin{itemize}
        \item[a)] Représenter l'expérience par un arbre des issues pondéré (probabilités sur les branches).
        \item[b)] Calculer la probabilité d'avoir exactement une pièce défectueuse.
        \item[c)] Calculer la probabilité d'avoir au moins une pièce conforme.
    \end{itemize}
    \item Dans un lot de 20 composants électroniques, 3 sont défectueux. On tire 2 composants \textbf{sans remise}.
    \begin{itemize}
        \item[a)] Construire l'arbre des issues (attention : les probabilités du 2ème niveau changent selon la 1ère branche).
        \item[b)] Calculer $P(\text{les deux sont conformes})$.
        \item[c)] Calculer $P(\text{au moins un est défectueux})$ en utilisant l'événement contraire.
    \end{itemize}
    \item Soit $\Omega = \{1,2,3,4,5,6,7,8,9,10\}$, $A = \{2,4,6,8,10\}$ (pairs), $B = \{3,6,9\}$ (multiples de 3).
    \begin{itemize}
        \item[a)] Déterminer $A \cup B$, $A \cap B$, $\overline{A}$, $A \setminus B$.
        \item[b)] Vérifier la formule $P(A \cup B) = P(A) + P(B) - P(A \cap B)$ en supposant les issues équiprobables.
    \end{itemize}
\end{enumerate}
\end{exercice}

\begin{corrige}[Exercice 1]
\begin{enumerate}
    \item \textbf{Arbre :} Racine $\to$ C ($0,95$) et D ($0,05$). Chaque nœud $\to$ C ($0,95$) et D ($0,05$).
    \item \textbf{Exactement 1 défectueuse :} Chemins (C,D) et (D,C). $P = 0,95 \times 0,05 + 0,05 \times 0,95 = 2 \times 0,0475 = 0,095$.
    \item \textbf{Au moins 1 conforme :} Contraire de « 2 défectueuses ». $P(\text{2 D}) = 0,05^2 = 0,0025$. $P(\ge 1 C) = 1 - 0,0025 = 0,9975$.
\end{enumerate}
\end{corrige}

\begin{corrige}[Exercice 2]
\begin{enumerate}
    \item \textbf{Arbre sans remise :} 1er niveau : C ($17/20$), D ($3/20$).
        Si 1er = C : reste 16 C, 3 D $\to$ 2ème niveau : C ($16/19$), D ($3/19$).
        Si 1er = D : reste 17 C, 2 D $\to$ 2ème niveau : C ($17/19$), D ($2/19$).
    \item $P(\text{C et C}) = \frac{17}{20} \times \frac{16}{19} = \frac{272}{380} = \frac{68}{95} \approx 0,7158$.
    \item $P(\ge 1 D) = 1 - P(\text{2 C}) = 1 - \frac{68}{95} = \frac{27}{95} \approx 0,2842$.
\end{enumerate}
\end{corrige}

\begin{corrige}[Exercice 3]
\begin{enumerate}
    \item $A \cup B = \{2,3,4,6,8,9,10\}$ (Card 7). $A \cap B = \{6\}$ (Card 1). $\overline{A} = \{1,3,5,7,9\}$ (Card 5). $A \setminus B = \{2,4,8,10\}$ (Card 4).
    \item $P(A) = 5/10 = 0,5$. $P(B) = 3/10 = 0,3$. $P(A \cap B) = 1/10 = 0,1$.
        $P(A \cup B) = 7/10 = 0,7$.
        Vérification : $0,5 + 0,3 - 0,1 = 0,7$. \checkmark
\end{enumerate}
\end{corrige}

\coursec{Probabilité d'un événement et probabilités conditionnelles}

\courssub{Rappel et transition : de l'expérience aléatoire à la mesure de l'incertitude}

Dans la section précédente, nous avons posé le vocabulaire fondamental : \cle{expérience aléatoire}, \cle{univers} $\Omega$, \cle{événement} (partie de $\Omega$), opérations sur les événements (réunion, intersection, complémentaire, incompatibilité). Nous savons décrire les issues possibles d'une situation aléatoire (lancer de dé, contrôle qualité, test médical). Il reste à quantifier la \cle{chance} qu'un événement se réalise. C'est l'objet de la \cle{probabilité}.

En Terminale technique, nous nous plaçons d'abord dans le cadre \textbf{fini et équiprobable} (toutes les issues élémentaires ont la même chance), ce qui couvre la grande majorité des applications industrielles (échantillonnage, jeux, fiabilité) et permet d'introduire rigoureusement le conditionnement.

\begin{definition}[Probabilité équiprobable]
Soit $\Omega$ un univers \textbf{fini} muni de l'équiprobabilité (chaque issue élémentaire $\omega \in \Omega$ a la même probabilité $1/|\Omega|$). Pour tout événement $A \subseteq \Omega$, la probabilité de $A$ est :
\[
P(A) = \frac{|A|}{|\Omega|} = \frac{\text{nombre de cas favorables à }A}{\text{nombre total de cas possibles}}.
\]
On a toujours $0 \le P(A) \le 1$, $P(\emptyset)=0$, $P(\Omega)=1$.
\end{definition}

\begin{propriete}[Propriétés usuelles (admises au Probatoire)]
Pour tous événements $A, B \subseteq \Omega$ :
\begin{itemize}
    \item $P(\overline{A}) = 1 - P(A)$.
    \item $P(A \cup B) = P(A) + P(B) - P(A \cap B)$.
    \item Si $A \cap B = \emptyset$ (incompatibles) : $P(A \cup B) = P(A) + P(B)$.
    \item Si $A \subseteq B$ : $P(A) \le P(B)$.
\end{itemize}
\end{propriete}

\begin{exemplebox}[Contrôle qualité en atelier]
Un lot de \num{200} pièces contient \num{12} pièces défectueuses. On prélève une pièce au hasard (équiprobabilité).
\begin{itemize}
    \item $\Omega$ = ensemble des \num{200} pièces, $|\Omega|=\num{200}$.
    \item $D$ = « pièce défectueuse », $|D|=\num{12}$.
    \item $P(D) = \frac{12}{200} = \num{0,06} = \num{6}\%$.
    \item $P(\text{conforme}) = 1 - \num{0,06} = \num{0,94}$.
\end{itemize}
\end{exemplebox}

\courssub{Probabilité conditionnelle : l'information change la donne}

Souvent, on apprend qu'un événement $B$ est réalisé (résultat d'un test, information partielle). La probabilité d'un autre événement $A$ doit être \textbf{réévaluée} dans ce nouvel univers réduit à $B$.

\begin{experience}[Découverte : test de dépistage]
Considérons une population de \num{1000} personnes.
\begin{itemize}
    \item \num{100} sont malades ($M$), \num{900} sont saines ($S$).
    \item Un test détecte la maladie : sur les \num{100} malades, \num{90} sont testés positifs ($T^+$) ; sur les \num{900} sains, \num{90} sont testés positifs (faux positifs).
\end{itemize}
\begin{enumerate}
    \item Compléter le tableau de contingence ci-dessous.
    \item Calculer $P(T^+)$ (probabilité brute d'être positif).
    \item Sachant que le test est positif, quelle est la probabilité d'être \textbf{réellement} malade ? Noter cette probabilité $P(M | T^+)$.
    \item Comparer $P(M | T^+)$ et $P(M)$. Que constatez-vous ?
\end{enumerate}

\begin{center}
\begin{tabular}{|c|c|c|c|}\hline
 & Malade ($M$) & Sain ($S$) & Total \\ \hline
Test $+$ ($T^+$) & 90 & 90 & 180 \\
Test $-$ ($T^-$) & 10 & 810 & 820 \\ \hline
Total & 100 & 900 & 1000 \\ \hline
\end{tabular}
\end{center}

\textbf{Analyse :} $P(M) = \num{0,10}$. Mais $P(M | T^+) = \frac{90}{180} = \num{0,50}$. L'information « test positif » multiplie par 5 la probabilité d'être malade. C'est l'essence du conditionnement.
\end{experience}

\begin{definition}[Probabilité conditionnelle]
Soient $A$ et $B$ deux événements de $\Omega$ avec $P(B) > 0$ ($B$ non impossible). La \cle{probabilité conditionnelle de $A$ sachant $B$} est :
\[
\boxed{P(A | B) = \frac{P(A \cap B)}{P(B)}}
\]
\textbf{Interprétation :} On restreint l'univers à $B$ ; les issues favorables à $A$ deviennent celles de $A \cap B$.
\end{definition}

\begin{attention}[Pièges classiques]
\begin{itemize}
    \item \textbf{Confusion $P(A|B)$ et $P(B|A)$ :} En général $P(A|B) \neq P(B|A)$. Dans l'exemple : $P(M|T^+)=\num{0,5}$ mais $P(T^+|M)=\num{0,9}$.
    \item \textbf{Condition impossible :} $P(A|B)$ n'est pas définie si $P(B)=0$.
    \item \textbf{Ordre des événements :} « Sachant $B$ » signifie que $B$ est l'hypothèse (dénominateur), $A$ est la thèse (numérateur).
\end{itemize}
\end{attention}

\begin{propriete}[Formule des probabilités composées (ou multiplication)]
Immédiatement de la définition : $P(A \cap B) = P(B) \times P(A | B) = P(A) \times P(B | A)$.
C'est la base de la construction des \textbf{arbres pondérés} (probabilités sur les branches = conditionnelles).
\end{propriete}

\begin{exemplebox}[Arbre pondéré : Test de dépistage (données réelles)]
Une maladie touche $\num{1}\%$ de la population ($P(M)=\num{0,01}$). Un test a une \textbf{sensibilité} $P(T^+|M)=\num{0,95}$ et une \textbf{spécificité} $P(T^-|S)=\num{0,98}$ (donc $P(T^+|S)=\num{0,02}$).
Construire l'arbre et calculer $P(T^+)$ et $P(M|T^+)$.
\end{exemplebox}

\begin{popfigure}
\begin{tikzpicture}[
    grow=right, level distance=3.5cm, sibling distance=1.5cm,
    edge from parent/.style={draw, thick, -{Stealth[length=3mm]}},
    every node/.style={font=\small},
    branch/.style={edge from parent/.append style={popblue}},
    leaf/.style={edge from parent/.append style={poporange}}
]
\node[draw, rounded corners, fill=popblueL, text width=2.2cm, align=center] {Population\\ $P(M)=\num{0,01}$\\ $P(S)=\num{0,99}$}
    child { node[draw, rounded corners, fill=popgreenL, text width=2.5cm, align=center] {Malade $M$\\ $P(T^+|M)=\num{0,95}$\\ $P(T^-|M)=\num{0,05}$}
        child { node[draw, circle, fill=poporangeL, align=center]{$T^+$\\ $P=\num{0,01}\times\num{0,95}=\num{0,0095}$} edge from parent node[above, sloped, pos=0.6] {\num{0,95}} }
        child { node[draw, circle, fill=poporangeL, align=center]{$T^-$\\ $P=\num{0,01}\times\num{0,05}=\num{0,0005}$} edge from parent node[below, sloped, pos=0.6] {\num{0,05}} }
        edge from parent node[above, sloped, pos=0.6] {\num{0,01}}
    }
    child { node[draw, rounded corners, fill=poppinkL, text width=2.5cm, align=center] {Sain $S$\\ $P(T^+|S)=\num{0,02}$\\ $P(T^-|S)=\num{0,98}$}
        child { node[draw, circle, fill=poporangeL, align=center]{$T^+$\\ $P=\num{0,99}\times\num{0,02}=\num{0,0198}$} edge from parent node[above, sloped, pos=0.6] {\num{0,02}} }
        child { node[draw, circle, fill=poporangeL, align=center]{$T^-$\\ $P=\num{0,99}\times\num{0,98}=\num{0,9702}$} edge from parent node[below, sloped, pos=0.6] {\num{0,98}} }
        edge from parent node[below, sloped, pos=0.6] {\num{0,99}}
    };
\end{tikzpicture}
\legende{Arbre pondéré du test de dépistage. Les probabilités aux feuilles sont les $P(\text{chemin}) = P(\text{racine}) \times P(\text{branche}_1) \times \dots$}
\end{popfigure}

\textbf{Calculs :}
\begin{align*}
P(T^+) &= P(M \cap T^+) + P(S \cap T^+) = \num{0,0095} + \num{0,0198} = \num{0,0293}. \\
P(M | T^+) &= \frac{P(M \cap T^+)}{P(T^+)} = \frac{\num{0,0095}}{\num{0,0293}} \approx \num{0,324} = \num{32,4}\%.
\end{align*}
\textbf{Interprétation technique :} Même avec un test performant (\num{95}\%/\num{98}\%), la rareté de la maladie (\num{1}\%) fait que \textbf{2 tests positifs sur 3 sont des faux positifs}. C'est le \cle{paradoxe du dépistage}, crucial en santé publique et maintenance prédictive.

\courssub{Formule des probabilités totales et Théorème de Bayes}

Ces deux résultats sont les piliers du \textbf{raisonnement inverse} : remonter de l'effet (test positif) à la cause (maladie).

\begin{propriete}[Formule des probabilités totales — Théorème (démonstration exigible au Probatoire)]
Soit $(B_1, B_2, \dots, B_n)$ une \cle{partition} de $\Omega$ (événements deux à deux incompatibles et $\bigcup B_i = \Omega$) avec $P(B_i) > 0$.
Pour tout événement $A$ :
\[
\boxed{P(A) = \sum_{i=1}^n P(B_i) \, P(A | B_i)}
\]
\end{propriete}

\begin{experience}[Démonstration guidée : Formule des probabilités totales]
\begin{enumerate}
    \item Écrire $A$ comme réunion d'événements incompatibles en utilisant la partition : $A = \bigcup_{i=1}^n (A \cap B_i)$.
    \item Justifier que les $(A \cap B_i)$ sont deux à deux incompatibles.
    \item Appliquer la $\sigma$-additivité (somme des probabilités pour événements incompatibles).
    \item Utiliser la formule des probabilités composées sur chaque terme $P(A \cap B_i)$.
    \item Conclure.
\end{enumerate}
\end{experience}

\textbf{Preuve formelle :}
$A = A \cap \Omega = A \cap (\bigcup B_i) = \bigcup (A \cap B_i)$. Les $(A \cap B_i)$ sont incompatibles car les $B_i$ le sont.
$P(A) = \sum P(A \cap B_i) = \sum P(B_i) P(A | B_i)$. \qed

\begin{propriete}[Théorème de Bayes — Théorème (démonstration exigible au Probatoire)]
Mêmes hypothèses. Pour tout $k \in \{1,\dots,n\}$ :
\[
\boxed{P(B_k | A) = \frac{P(B_k) \, P(A | B_k)}{\sum_{i=1}^n P(B_i) \, P(A | B_i)}}
\]
\end{propriete}

\begin{experience}[Démonstration guidée : Théorème de Bayes]
\begin{enumerate}
    \item Écrire la définition de $P(B_k | A)$.
    \item Remplacer le numérateur par la formule des probabilités composées.
    \item Remplacer le dénominateur $P(A)$ par la formule des probabilités totales.
    \item Conclure.
\end{enumerate}
\end{experience}

\textbf{Preuve formelle :}
$P(B_k | A) = \frac{P(B_k \cap A)}{P(A)} = \frac{P(B_k)P(A|B_k)}{\sum P(B_i)P(A|B_i)}$. \qed

\begin{methode}[Résoudre un problème de probabilités conditionnelles / Bayes en 4 étapes]
\begin{enumerate}
    \item \textbf{Identifier la partition (causes) et l'événement observé (effet).} Noter $B_1,\dots,B_n$ les causes (ex: Malade/Sain, Machine 1/2/3) et $A$ l'effet observé (Test +, Pièce défectueuse).
    \item \textbf{Lister les données :} $P(B_i)$ (probabilités \textit{a priori} des causes) et $P(A | B_i)$ (probabilités de l'effet \textit{sachant} chaque cause = fiabilité du test, taux de défaut machine).
    \item \textbf{Calculer $P(A)$ par la formule des probabilités totales.} C'est la probabilité brute d'observer l'effet.
    \item \textbf{Appliquer Bayes pour la cause cherchée :} $P(B_k | A) = \frac{P(B_k)P(A|B_k)}{P(A)}$. Interpréter le résultat en pourcentage.
\end{enumerate}
\end{methode}

\begin{exoresolu}[Contrôle qualité multi-machines]
Un atelier possède 3 machines $M_1, M_2, M_3$ produisant respectivement $\num{50}\%, \num{30}\%, \num{20}\%$ des pièces.
Leurs taux de rebuts (défaut $D$) sont : $\num{2}\%, \num{3}\%, \num{5}\%$.
On choisit une pièce au hasard dans la production globale.
\begin{enumerate}
    \item Quelle est la probabilité qu'elle soit défectueuse ?
    \item Sachant qu'elle est défectueuse, quelle est la probabilité qu'elle vienne de $M_3$ ?
\end{enumerate}
\tcblower
\textbf{Solution détaillée}
\emph{Étape 1 : Partition et événement.}
$B_1=M_1, B_2=M_2, B_3=M_3$ (partition). $A = D$ (défectueuse).
\emph{Étape 2 : Données.}
$P(M_1)=\num{0,50},\; P(M_2)=\num{0,30},\; P(M_3)=\num{0,20}$.
$P(D|M_1)=\num{0,02},\; P(D|M_2)=\num{0,03},\; P(D|M_3)=\num{0,05}$.
\emph{Étape 3 : Probabilités totales.}
\begin{align*}
P(D) &= P(M_1)P(D|M_1) + P(M_2)P(D|M_2) + P(M_3)P(D|M_3) \\
     &= \num{0,50} \times \num{0,02} + \num{0,30} \times \num{0,03} + \num{0,20} \times \num{0,05} \\
     &= \num{0,010} + \num{0,009} + \num{0,010} = \num{0,029}.
\end{align*}
\emph{Étape 4 : Bayes pour $M_3$.}
\[
P(M_3 | D) = \frac{P(M_3)P(D|M_3)}{P(D)} = \frac{\num{0,20} \times \num{0,05}}{\num{0,029}} = \frac{\num{0,010}}{\num{0,029}} \approx \num{0,3448} = \num{34,5}\%.
\]
\textbf{Conclusion :} Bien que $M_3$ ne produise que \num{20}\% du volume, elle génère \textbf{\num{34,5}\% des défauts}. C'est la machine à surveiller en priorité (diagramme de Pareto probabiliste).
\end{exoresolu}



\courssub{Indépendance d'événements : quand l'information ne change rien}

Deux événements sont \cle{indépendants} si la réalisation de l'un ne modifie pas la probabilité de l'autre. C'est une propriété \textbf{mathématique} (pas seulement intuitive), qui se vérifie par un calcul.

\begin{definition}[Indépendance de deux événements]
Les événements $A$ et $B$ (avec $P(A)>0, P(B)>0$) sont \textbf{indépendants} si et seulement si :
\[
\boxed{P(A \cap B) = P(A) \times P(B)}
\]
Équivalent (si $P(B)>0$) : $P(A | B) = P(A)$.
Équivalent (si $P(A)>0$) : $P(B | A) = P(B)$.
\end{definition}

\begin{attention}[Indépendance $\neq$ Incompatibilité — Erreur fatale au Probatoire]
\begin{itemize}
    \item \textbf{Incompatibles :} $A \cap B = \emptyset \Rightarrow P(A \cap B) = 0$. Si $P(A)>0$ et $P(B)>0$, alors $P(A)P(B) > 0$. \textbf{Donc incompatibles $\Rightarrow$ NON indépendants} (savoir que $A$ est réalisé rend $B$ impossible).
    \item \textbf{Indépendants :} $P(A \cap B) = P(A)P(B) > 0$ (si probabilités non nulles). \textbf{Donc indépendants $\Rightarrow$ NON incompatibles} (ils peuvent se réaliser ensemble).
    \item \textbf{Visuel :} Dans un diagramme de Venn, l'indépendance correspond à un « produit des aires » (proportionnalité), pas à une séparation.
\end{itemize}
\end{attention}

\begin{exemplebox}[Test d'indépendance : Dé et Parité]
On lance un dé équilibré ($\Omega=\{1..6\}$).
$A$ = « résultat pair » = $\{2,4,6\}$, $P(A)=\num{1/2}$.
$B$ = « résultat $\le 3$ » = $\{1,2,3\}$, $P(B)=\num{1/2}$.
$A \cap B = \{2\}$, $P(A \cap B) = \num{1/6}$.
$P(A)P(B) = \num{1/4} \neq \num{1/6}$. \textbf{A et B ne sont pas indépendants.}
\emph{Intuition :} Savoir que le résultat est $\le 3$ change la chance qu'il soit pair (\num{1} chance sur \num{3} au lieu de \num{1} sur \num{2}).

$C$ = « résultat $\in \{1,4\}$ », $P(C)=\num{1/3}$.
$A \cap C = \{4\}$, $P(A \cap C)=\num{1/6}$.
$P(A)P(C) = \num{1/2} \times \num{1/3} = \num{1/6}$. \textbf{A et C sont indépendants.}
\end{exemplebox}

\begin{propriete}[Indépendance par paires vs Indépendance mutuelle (rappel)]
Pour trois événements $A,B,C$ :
\begin{itemize}
    \item \textbf{Indépendance par paires :} $P(A\cap B)=P(A)P(B)$, $P(A\cap C)=P(A)P(C)$, $P(B\cap C)=P(B)P(C)$.
    \item \textbf{Indépendance mutuelle :} Indépendance par paires \textbf{ET} $P(A\cap B\cap C) = P(A)P(B)P(C)$.
    \item L'indépendance par paires n'entraîne pas l'indépendance mutuelle (contre-exemple classique : deux lancers de pièce, $A$=premier pile, $B$=deuxième pile, $C$=même face).
\end{itemize}
Au Probatoire, on teste l'indépendance par paires sauf mention contraire.
\end{propriete}

\courssub{Applications intégrées : du diagnostic à la fiabilité}

\begin{exoresolu}[Fiabilité système : redondance active]
Un système critique (ex: freinage avion, serveur bancaire) possède deux composants $C_1, C_2$ en \textbf{parallèle} (redondance active) : le système fonctionne si \textbf{au moins un} composant fonctionne.
$P(C_1 \text{ fonctionne}) = \num{0,99}$, $P(C_2 \text{ fonctionne}) = \num{0,98}$. On suppose l'indépendance.
Calculer la fiabilité du système $P(S)$.
\tcblower
\textbf{Solution}
$S = C_1 \cup C_2$.
$P(S) = P(C_1) + P(C_2) - P(C_1 \cap C_2)$.
Par indépendance : $P(C_1 \cap C_2) = \num{0,99} \times \num{0,98} = \num{0,9702}$.
$P(S) = \num{0,99} + \num{0,98} - \num{0,9702} = \num{0,9998}$.
\textbf{Gain :} La redondance fait passer la probabilité de panne de $\num{1}\%$ (pour $C_1$ seul) à $\num{0,02}\%$ (facteur 50). C'est le principe de la \cle{tolérance aux pannes}.
\end{exoresolu}

\begin{exercice}[Entraînement : Diagnostic automobile]
Un garage reçoit des voitures de deux marques : $\num{60}\%$ Marque A, $\num{40}\%$ Marque B.
Sur la Marque A, $\num{10}\%$ ont un défaut de freinage ($D$). Sur la Marque B, $\num{25}\%$ ont ce défaut.
On choisit une voiture au hasard. On teste les freins : le test révèle un défaut ($T^+$).
Le test a une sensibilité $P(T^+|D)=\num{0,95}$ et une spécificité $P(T^-|\overline{D})=\num{0,90}$.
\begin{enumerate}
    \item Calculer la probabilité qu'une voiture ait un défaut de freinage ($P(D)$).
    \item Calculer la probabilité que le test soit positif ($P(T^+)$).
    \item Sachant que le test est positif, quelle est la probabilité que la voiture ait \textbf{réellement} un défaut ($P(D|T^+)$) ?
    \item La marque de la voiture et le résultat du test sont-ils indépendants ? Justifier.
\end{enumerate}
\end{exercice}

\begin{corrige}[Exercice Diagnostic automobile]
\emph{Notations :} $A$=Marque A, $B$=Marque B, $D$=Défaut, $T^+$=Test positif.
Données : $P(A)=\num{0,6}, P(B)=\num{0,4}$. $P(D|A)=\num{0,1}, P(D|B)=\num{0,25}$.
$P(T^+|D)=\num{0,95}, P(T^+|\overline{D})=\num{0,10}$ (car spécificité $\num{0,90}$).

1. \textbf{Probabilité totale pour $D$ (partition marques) :}
$P(D) = P(A)P(D|A) + P(B)P(D|B) = \num{0,6}\times\num{0,1} + \num{0,4}\times\num{0,25} = \num{0,06} + \num{0,10} = \num{0,16}$.

2. \textbf{Probabilité totale pour $T^+$ (partition $D/\overline{D}$) :}
Il faut $P(\overline{D}) = 1 - \num{0,16} = \num{0,84}$.
$P(T^+) = P(D)P(T^+|D) + P(\overline{D})P(T^+|\overline{D}) = \num{0,16}\times\num{0,95} + \num{0,84}\times\num{0,10} = \num{0,152} + \num{0,084} = \num{0,236}$.

3. \textbf{Bayes pour $P(D|T^+)$ :}
$P(D|T^+) = \frac{P(D)P(T^+|D)}{P(T^+)} = \frac{\num{0,152}}{\num{0,236}} \approx \num{0,644} = \num{64,4}\%$.
Le test positif multiplie par 4 la probabilité de défaut (de \num{16}\% à \num{64}\%), mais \num{35}\% des alertes restent des faux positifs.

4. \textbf{Indépendance Marque / Test ?}
Tester $P(T^+|A)$ vs $P(T^+)$.
$P(T^+|A) = P(D|A)P(T^+|D) + P(\overline{D}|A)P(T^+|\overline{D}) = \num{0,1}\times\num{0,95} + \num{0,9}\times\num{0,10} = \num{0,095} + \num{0,09} = \num{0,185}$.
$P(T^+) = \num{0,236}$.
$\num{0,185} \neq \num{0,236} \Rightarrow$ \textbf{NON indépendants}. La marque influence la prévalence du défaut, donc le taux de positivité du test.
\end{corrige}

\begin{aretenir}[Synthèse : Boîte à outils Probabilités Conditionnelles]
\begin{itemize}
    \item \textbf{Conditionnelle :} $P(A|B) = \frac{P(A\cap B)}{P(B)}$ ($P(B)>0$). « Sachant $B$ » $\to$ dénominateur $P(B)$.
    \item \textbf{Arbre pondéré :} Racine = causes ($P(B_i)$), Branches = effets sachant causes ($P(A|B_i)$), Feuilles = jointes ($P(B_i\cap A)$). Somme des feuilles = 1.
    \item \textbf{Probabilités totales :} $P(A) = \sum P(B_i)P(A|B_i)$. (Partition $B_i$ = causes).
    \item \textbf{Bayes :} $P(B_k|A) = \frac{P(B_k)P(A|B_k)}{\sum P(B_i)P(A|B_i)}$. (Remonter cause $\leftarrow$ effet).
    \item \textbf{Indépendance :} $P(A\cap B) = P(A)P(B)$ $\iff$ $P(A|B)=P(A)$. \textbf{Vérifier par le calcul}, pas par l'intuition.
    \item \textbf{Incompatibles $\neq$ Indépendants.} Si $P(A)>0, P(B)>0$ : Incompatibles $\Rightarrow$ Dépendants.
\end{itemize}
\end{aretenir}

\begin{savaistu}[Le test du VIH et la prévalence : enjeu de santé publique]
Au Cameroun, la prévalence du VIH est d'environ $\num{2,7}\%$ (données récentes). Les tests rapides ont une sensibilité $>\num{99}\%$ et spécificité $>\num{99}\%$.
Appliquons Bayes : $P(\text{VIH}|T^+) = \frac{\num{0,027} \times \num{0,99}}{\num{0,027}\times\num{0,99} + \num{0,973}\times\num{0,01}} \approx \frac{\num{0,0267}}{\num{0,0364}} \approx \num{73}\%$.
Même avec un test excellent, \textbf{1 test positif sur 4 est un faux positif} dans la population générale. C'est pourquoi on \textbf{confirme toujours par un second test} (différent, plus spécifique) avant tout diagnostic définitif. En maintenance industrielle, c'est la même logique : une alarme capteur déclenche une inspection visuelle avant l'arrêt de la ligne.
\end{savaistu}

\coursec{Variables aléatoires et épreuves de Bernoulli}

Dans la section précédente, nous avons appris à calculer la probabilité d'événements, y compris conditionnels. Mais dans la vie pratique — contrôle qualité en usine, jeux de hasard, gestion de stock — on s'intéresse souvent non pas à un événement isolé, mais à une \cle{grandeur numérique} qui dépend du hasard : le nombre de pièces défectueuses dans un lot, le gain d'un joueur, le nombre de pannes d'une machine sur un mois. Cette section introduit l'outil qui permet d'étudier ces grandeurs : la \cle{variable aléatoire}, puis le modèle le plus utile au programme : la \cle{loi binomiale}.

\courssub{Variable aléatoire discrète et loi de probabilité}

\textbf{Intuition.} Un technicien contrôle un lot de 3 pièces sortant d'une machine. Le résultat du contrôle est aléatoire, mais ce qui l'intéresse, c'est le \emph{nombre} de pièces défectueuses : ce nombre peut valoir 0, 1, 2 ou 3. À chaque issue de l'expérience, on associe donc un nombre réel : c'est exactement l'idée d'une variable aléatoire.

\begin{definition}[Variable aléatoire discrète]
Soit $\Omega$ l'univers d'une expérience aléatoire. Une \cle{variable aléatoire discrète} $X$ est une application de $\Omega$ dans $\mathbb{R}$ qui, à chaque issue $\omega$, associe un nombre réel $X(\omega)$, et qui ne prend qu'un nombre fini de valeurs $x_1, x_2, \dots, x_n$.
\end{definition}

\begin{definition}[Loi de probabilité]
La \cle{loi de probabilité} d'une variable aléatoire discrète $X$ est la donnée de toutes les valeurs $x_i$ prises par $X$ et des probabilités associées $p_i = P(X = x_i)$. On la présente dans un \cle{tableau de répartition} :
\begin{center}
\begin{tabular}{|c|ccccc|}\hline
$x_i$ & $x_1$ & $x_2$ & $\dots$ & $x_n$ & Total \\ \hline
$P(X=x_i)$ & $p_1$ & $p_2$ & $\dots$ & $p_n$ & $1$ \\ \hline
\end{tabular}
\end{center}
On a toujours $p_i \geq 0$ pour tout $i$, et $p_1 + p_2 + \dots + p_n = 1$.
\end{definition}

\begin{attention}[Vérification obligatoire]
Avant tout calcul, vérifie que la somme des probabilités vaut $1$. Si $\sum p_i \neq 1$, le tableau n'est pas une loi de probabilité : il y a une erreur dans l'énoncé ou dans tes calculs. Attention aussi à ne pas oublier une valeur possible de $X$ (par exemple la valeur $0$).
\end{attention}

\begin{exoresolu}[Loi de probabilité d'un contrôle qualité]
Un atelier contrôle 2 pièces tirées au hasard dans une production contenant $10\,\%$ de pièces défectueuses (on suppose les tirages indépendants). On note $X$ le nombre de pièces défectueuses parmi les 2. Déterminer la loi de probabilité de $X$.
\tcblower
$X$ prend les valeurs $0$, $1$ et $2$. Avec $p = 0{,}1$ (probabilité qu'une pièce soit défectueuse) et $q = 1-p = 0{,}9$ :
\begin{align*}
P(X=0) &= 0{,}9 \times 0{,}9 = 0{,}81,\\
P(X=1) &= 2 \times 0{,}1 \times 0{,}9 = 0{,}18,\\
P(X=2) &= 0{,}1 \times 0{,}1 = 0{,}01.
\end{align*}
(Le facteur $2$ dans $P(X=1)$ vient des deux cas possibles : défectueuse puis bonne, ou bonne puis défectueuse.) Vérification : $0{,}81 + 0{,}18 + 0{,}01 = 1$. La loi de $X$ est :
\begin{center}
\begin{tabular}{|c|ccc|c|}\hline
$x_i$ & $0$ & $1$ & $2$ & Total \\ \hline
$P(X=x_i)$ & $0{,}81$ & $0{,}18$ & $0{,}01$ & $1$ \\ \hline
\end{tabular}
\end{center}
\end{exoresolu}

\courssub{Espérance, variance et écart-type}

\textbf{Intuition.} Si un joueur répète un jeu de hasard un très grand nombre de fois, quel sera son gain \emph{moyen} par partie ? C'est l'\cle{espérance}. Et pour savoir si les gains sont stables ou très dispersés autour de cette moyenne, on utilise la \cle{variance} et l'\cle{écart-type}.

\begin{definition}[Espérance, variance, écart-type]
Soit $X$ une variable aléatoire discrète prenant les valeurs $x_1, \dots, x_n$ avec les probabilités $p_1, \dots, p_n$.
\begin{itemize}
\item L'\cle{espérance} de $X$ est : $\displaystyle E(X) = \sum_{i=1}^{n} p_i\, x_i = p_1 x_1 + p_2 x_2 + \dots + p_n x_n$.
\item La \cle{variance} de $X$ est : $\displaystyle V(X) = \sum_{i=1}^{n} p_i \bigl(x_i - E(X)\bigr)^2 = E(X^2) - \bigl[E(X)\bigr]^2$.
\item L'\cle{écart-type} est : $\sigma(X) = \sqrt{V(X)}$.
\end{itemize}
\end{definition}

\begin{aretenir}[Interprétation]
\begin{itemize}
\item $E(X)$ est la \textbf{valeur moyenne} de $X$ sur un grand nombre de répétitions de l'expérience.
\item $\sigma(X)$ mesure la \textbf{dispersion} des valeurs autour de la moyenne : plus $\sigma$ est grand, plus les résultats sont imprévisibles.
\end{itemize}
\end{aretenir}

\begin{exemplebox}[Calcul complet sur l'exemple de l'atelier]
Reprenons la loi de $X$ ci-dessus : $E(X) = 0{,}81 \times 0 + 0{,}18 \times 1 + 0{,}01 \times 2 = 0{,}20$. En moyenne, on trouve donc $0{,}2$ pièce défectueuse par contrôle de 2 pièces (cohérent : $10\,\%$ de 2). Pour la variance, on calcule d'abord $E(X^2) = 0{,}81 \times 0^2 + 0{,}18 \times 1^2 + 0{,}01 \times 2^2 = 0{,}22$, d'où $V(X) = E(X^2) - [E(X)]^2 = 0{,}22 - 0{,}20^2 = 0{,}18$ et $\sigma(X) = \sqrt{0{,}18} \approx 0{,}42$.
\end{exemplebox}

\begin{attention}[Pièges fréquents]
\begin{itemize}
\item Ne confonds pas $E(X^2)$ et $[E(X)]^2$ : ce sont deux quantités différentes, et la formule de la variance utilise les deux.
\item L'écart-type est la \textbf{racine carrée} de la variance : n'oublie pas cette dernière étape quand l'énoncé demande $\sigma$.
\item Une variance est toujours \textbf{positive ou nulle}. Un résultat négatif signale une erreur de calcul.
\end{itemize}
\end{attention}

\courssub{Épreuve de Bernoulli}

\begin{definition}[Épreuve de Bernoulli]
Une \cle{épreuve de Bernoulli} est une expérience aléatoire qui n'a que \textbf{deux issues} :
\begin{itemize}
\item le \cle{succès} $S$, de probabilité $p$ ;
\item l'\cle{échec} $\overline{S}$, de probabilité $q = 1 - p$.
\end{itemize}
Le nombre $p$ (avec $0 < p < 1$) est appelé le \cle{paramètre} de l'épreuve. La variable aléatoire $X$ qui vaut $1$ en cas de succès et $0$ en cas d'échec est appelée \cle{variable de Bernoulli} ; on a $E(X) = p$ et $V(X) = p(1-p)$.
\end{definition}

\begin{exemplebox}[Exemples d'épreuves de Bernoulli]
\begin{itemize}
\item Contrôler une pièce : défectueuse ($p = 0{,}02$) ou conforme ($q = 0{,}98$).
\item Lancer une pièce de monnaie : pile ($p = 0{,}5$) ou face.
\item Un client entre dans un magasin : il achète ($p$) ou il n'achète pas.
\end{itemize}
\end{exemplebox}

La fonction de répartition d'une variable de Bernoulli $X$ est définie par $F(x) = P(X \leq x)$. Elle vaut $0$ avant $0$, puis $1-p$ entre $0$ (compris) et $1$ (non compris), puis $1$ à partir de $1$ : c'est une fonction « en escalier ».

\begin{popfigure}
\begin{center}
\begin{tikzpicture}[scale=1.1]
\draw[->, thick, popdark] (-1.2,0) -- (2.6,0) node[below left] {$x$};
\draw[->, thick, popdark] (0,-0.3) -- (0,1.5) node[left] {$F(x)$};
\draw[dashed, popmuted] (0,1) node[left] {$1$} -- (2.4,1);
\draw[dashed, popmuted] (0,0.6) node[left] {$1-p$} -- (0,0.6);
\draw[very thick, popblue] (-1.1,0) -- (0,0);
\draw[very thick, popblue] (0,0.6) -- (1,0.6);
\draw[very thick, popblue] (1,1) -- (2.4,1);
\fill[popblue] (0,0.6) circle (1.6pt);
\fill[white, draw=popblue, thick] (0,0) circle (1.6pt);
\fill[popblue] (1,1) circle (1.6pt);
\fill[white, draw=popblue, thick] (1,0.6) circle (1.6pt);
\draw[dashed, popmuted] (1,0) node[below] {$1$} -- (1,1);
\node[below left] at (0,0) {$0$};
\end{tikzpicture}
\end{center}
\legende{Fonction de répartition d'une variable de Bernoulli : paliers en $0$ et en $1$ (points pleins = valeurs prises, points creux = valeurs exclues).}
\end{popfigure}

\courssub{Schéma de Bernoulli et loi binomiale}

\textbf{Intuition.} On ne contrôle pas une seule pièce, mais un lot de $n$ pièces, dans les mêmes conditions et indépendamment les unes des autres. On compte alors le nombre total de succès : c'est le \cle{schéma de Bernoulli}.

\begin{definition}[Schéma de Bernoulli, loi binomiale]
Un \cle{schéma de Bernoulli} est la répétition de $n$ épreuves de Bernoulli \textbf{identiques et indépendantes}, de paramètre $p$. La variable aléatoire $X$ qui compte le nombre de succès suit la \cle{loi binomiale} de paramètres $n$ et $p$, notée $\mathcal{B}(n,p)$. Pour tout entier $k$ compris entre $0$ et $n$ :
\[ P(X = k) = \binom{n}{k} p^k (1-p)^{n-k}, \qquad \text{où } \binom{n}{k} = \frac{n!}{k!(n-k)!}. \]
Le coefficient $\binom{n}{k}$ (lu « $k$ parmi $n$ ») compte le nombre de façons de placer les $k$ succès parmi les $n$ épreuves.
\end{definition}

\begin{attention}[Épreuve ou schéma ?]
Ne confonds pas :
\begin{itemize}
\item l'\textbf{épreuve de Bernoulli} : une seule expérience ($n = 1$), $X$ vaut $0$ ou $1$ ;
\item le \textbf{schéma de Bernoulli} : $n > 1$ répétitions, $X$ peut valoir $0, 1, \dots, n$.
\end{itemize}
Avant d'appliquer la formule binomiale, vérifie toujours les trois conditions : \textbf{deux issues}, épreuves \textbf{identiques} (même $p$) et \textbf{indépendantes}. Un tirage \emph{sans remise} dans un petit lot n'est pas un schéma de Bernoulli ; en revanche, si le lot est très grand, on assimile le tirage à un tirage avec remise et le modèle binomial s'applique.
\end{attention}

\begin{propriete}[Espérance et variance de la loi binomiale]
Si $X \sim \mathcal{B}(n,p)$, alors :
\[ E(X) = np \qquad \text{et} \qquad V(X) = np(1-p), \quad \sigma(X) = \sqrt{np(1-p)}. \]
Ces formules ne se démontrent pas au Probatoire ; elles s'appliquent directement.
\end{propriete}

\begin{experience}[D'où vient la formule $E(X) = np$ ?]
Une épreuve de Bernoulli a pour espérance $p$. Le schéma de Bernoulli répète $n$ fois cette épreuve, et l'espérance du nombre total de succès est la somme des espérances : $E(X) = \underbrace{p + p + \dots + p}_{n \text{ fois}} = np$. Retiens l'idée : sur $n$ essais avec une proportion $p$ de succès, on attend en moyenne $np$ succès.
\end{experience}

\begin{methode}[Calculer $P(X = k)$, $P(X \leq k)$ ou $P(X \geq k)$]
\begin{enumerate}
\item Identifier $n$, $p$ et vérifier les conditions du schéma de Bernoulli.
\item Pour $P(X = k)$ : utiliser directement la formule $\binom{n}{k} p^k (1-p)^{n-k}$ (touche \texttt{nCr} de la calculatrice pour les coefficients).
\item Pour $P(X \leq k)$ : additionner $P(X=0) + P(X=1) + \dots + P(X=k)$, ou utiliser la fonction \texttt{binomFRép} / \texttt{BinomCD} de la calculatrice.
\item Pour $P(X \geq k)$ : passer par l'événement contraire : $P(X \geq k) = 1 - P(X \leq k-1)$.
\item Arrondir au degré de précision demandé (souvent $10^{-3}$ ou $10^{-4}$).
\end{enumerate}
\end{methode}

\begin{exoresolu}[Contrôle qualité dans une usine]
Une usine de Douala produit $2\,\%$ de pièces défectueuses. On prélève au hasard un lot de $100$ pièces dans la production (la production est suffisamment grande pour assimiler le prélèvement à des tirages indépendants). On note $X$ le nombre de pièces défectueuses dans le lot.
\begin{enumerate}
\item Justifier que $X$ suit une loi binomiale et donner ses paramètres.
\item Calculer $E(X)$ et interpréter.
\item Calculer la probabilité d'avoir au plus $3$ pièces défectueuses dans le lot (arrondir à $10^{-3}$).
\end{enumerate}
\tcblower
\textbf{1.} Chaque pièce contrôlée est une épreuve de Bernoulli : défectueuse (succès, $p = 0{,}02$) ou conforme. Les $100$ contrôles sont identiques et indépendants. Donc $X \sim \mathcal{B}(100\,;\,0{,}02)$.

\textbf{2.} $E(X) = np = 100 \times 0{,}02 = 2$. En moyenne, un lot de $100$ pièces contient $2$ pièces défectueuses.

\textbf{3.} « Au plus $3$ » signifie $X \leq 3$ :
\begin{align*}
P(X \leq 3) &= P(X=0) + P(X=1) + P(X=2) + P(X=3)\\
&= \binom{100}{0} 0{,}02^0\, 0{,}98^{100} + \binom{100}{1} 0{,}02^1\, 0{,}98^{99} \\
&\quad + \binom{100}{2} 0{,}02^2\, 0{,}98^{98} + \binom{100}{3} 0{,}02^3\, 0{,}98^{97}\\
&\approx 0{,}1326 + 0{,}2707 + 0{,}2734 + 0{,}1823 \approx 0{,}859.
\end{align*}
La probabilité d'avoir au plus $3$ pièces défectueuses est d'environ $0{,}859$, soit environ $86\,\%$.
\end{exoresolu}

\begin{popfigure}
\begin{center}
\begin{tikzpicture}
\begin{axis}[
    width=0.85\linewidth, height=5.5cm,
    ybar, bar width=8pt,
    xlabel={$k$}, ylabel={$P(X=k)$},
    xmin=-0.6, xmax=10.6, ymin=0, ymax=0.28,
    xtick={0,1,2,3,4,5,6,7,8,9,10},
    axis lines=left,
    tick label style={font=\small},
    label style={font=\small}
]
\addplot[fill=popblue, draw=popdark] coordinates {
(0,0.0282) (1,0.1211) (2,0.2335) (3,0.2668) (4,0.2001)
(5,0.1029) (6,0.0368) (7,0.0090) (8,0.0014) (9,0.0001) (10,0.0000)
};
\end{axis}
\end{tikzpicture}
\end{center}
\legende{Histogramme de la loi binomiale $\mathcal{B}(10\,;\,0{,}3)$ : le maximum est atteint en $k = 3$, proche de l'espérance $np = 3$.}
\end{popfigure}

\begin{attention}[Erreurs classiques sur la loi binomiale]
\begin{itemize}
\item Oublier le coefficient $\binom{n}{k}$ dans $P(X=k)$ : sans lui, on ne compte qu'un seul arrangement des succès.
\item Confondre $P(X \leq k)$ et $P(X < k)$ : le second exclut la valeur $k$.
\item Écrire $P(X \geq k) = 1 - P(X \leq k)$ : c'est faux ! La bonne formule est $P(X \geq k) = 1 - P(X \leq k-1)$.
\end{itemize}
\end{attention}

\courssub{Vers la loi normale : le théorème de Moivre-Laplace (évocation)}

Observe l'histogramme de $\mathcal{B}(10\,;\,0{,}3)$ ci-dessus : il a une forme « en cloche ». Lorsque $n$ devient grand, cette forme en cloche se rapproche d'une courbe célèbre : la \cle{courbe en cloche} de la \cle{loi normale}. C'est le contenu du théorème de \cle{Moivre-Laplace} : pour $n$ grand, la loi binomiale $\mathcal{B}(n,p)$ peut être approchée par une loi normale de moyenne $np$ et d'écart-type $\sqrt{np(1-p)}$. Cette approximation, très utilisée en contrôle qualité et dans les sondages, sera étudiée plus en détail dans la suite du programme ; ici, il suffit de retenir l'idée : \emph{la répétition d'un grand nombre d'épreuves identiques produit une répartition en cloche autour de la moyenne $np$}.

\begin{savaistu}[La loi binomiale, un modèle universel]
La loi binomiale modélise le nombre de succès dans $n$ épreuves indépendantes et identiques : contrôle qualité dans les usines, sondages électoraux, tests médicaux, tirs au but... Elle porte le nom de Jacques Bernoulli (1654--1705), mathématicien suisse dont l'ouvrage \emph{Ars Conjectandi} (1713) jette les bases du calcul des probabilités. Au Cameroun, les bureaux d'études et les services de contrôle qualité des entreprises agroalimentaires utilisent quotidiennement ce modèle pour décider si un lot de produits peut être mis sur le marché.
\end{savaistu}

\begin{exercice}[À toi de jouer]
Un candidat répond au hasard à un QCM de $8$ questions ; chaque question a $4$ propositions dont une seule est correcte. On note $X$ le nombre de bonnes réponses.
\begin{enumerate}
\item Justifier que $X$ suit une loi binomiale et préciser ses paramètres.
\item Calculer $E(X)$ et $\sigma(X)$.
\item Calculer la probabilité que le candidat ait au moins une bonne réponse (arrondir à $10^{-3}$).
\end{enumerate}
\end{exercice}

\begin{corrige}[Exercice 3.5]
\textbf{1.} Chaque question est une épreuve de Bernoulli : bonne réponse (succès) avec $p = \dfrac{1}{4} = 0{,}25$. Les $8$ réponses sont indépendantes et identiques, donc $X \sim \mathcal{B}(8\,;\,0{,}25)$.

\textbf{2.} $E(X) = np = 8 \times 0{,}25 = 2$ ; $V(X) = np(1-p) = 8 \times 0{,}25 \times 0{,}75 = 1{,}5$, d'où $\sigma(X) = \sqrt{1{,}5} \approx 1{,}22$.

\textbf{3.} $P(X \geq 1) = 1 - P(X = 0) = 1 - 0{,}75^8 \approx 1 - 0{,}1001 \approx 0{,}900$. Le candidat a environ $90\,\%$ de chances d'avoir au moins une bonne réponse en répondant au hasard.
\end{corrige}

\begin{aretenir}[L'essentiel de la section]
\begin{itemize}
\item Une \textbf{variable aléatoire discrète} associe un nombre à chaque issue ; sa \textbf{loi de probabilité} est un tableau dont la somme vaut $1$.
\item $E(X) = \sum p_i x_i$ (moyenne), $V(X) = E(X^2) - [E(X)]^2$ (dispersion), $\sigma = \sqrt{V}$.
\item \textbf{Épreuve de Bernoulli} : deux issues, succès de probabilité $p$.
\item \textbf{Loi binomiale} $\mathcal{B}(n,p)$ : $P(X=k) = \binom{n}{k} p^k (1-p)^{n-k}$, avec $E(X) = np$ et $V(X) = np(1-p)$.
\item Pour $n$ grand, la loi binomiale s'approche d'une courbe en cloche (Moivre-Laplace).
\end{itemize}
\end{aretenir}

\coursec{Présentation générale des coniques et propriétés}

Après avoir exploré le hasard et la modélisation des phénomènes aléatoires avec les variables aléatoires et les épreuves de Bernoulli, nous changeons radicalement de registre pour entrer dans la géométrie analytique. Les \textbf{coniques} sont des courbes d'une élégance remarquable, nées de l'intersection d'un plan et d'un cône. Elles ne sont pas de simples objets théoriques : elles décrivent les orbites des planètes (ellipses), la trajectoire d'un projectile ou le faisceau d'un phare (paraboles), et la forme de certaines antennes ou ombres portées (hyperboles). Maîtriser leurs équations réduites et leurs propriétés géométriques est indispensable pour tout technicien confronté à l'optique, la mécanique céleste, le calcul de structures ou le traitement du signal.

\courssub{Définition géométrique : sections d'un cône de révolution}

Imaginons un \textbf{cône de révolution} à deux nappes (deux cônes accolés par leur sommet $S$), engendré par une droite génératrice tournant autour d'un axe $\Delta$. Coupons ce cône par un plan $\mathcal{P}$ qui ne passe pas par le sommet $S$. Selon l'inclinaison de $\mathcal{P}$ par rapport à l'axe $\Delta$, nous obtenons trois types de courbes fermées ou ouvertes : les \textbf{coniques}.

\begin{popfigure}
\begin{tikzpicture}[scale=0.9, >=Stealth]
% Cone double nappe
\def\coneangle{30}
\def\coneheight{4}
\def\radius{2.309} % tan(30)*4
% Back lines of cone
\draw[popmuted, thick] (0,0) -- (-\radius,-\coneheight);
\draw[popmuted, thick] (0,0) -- (\radius,-\coneheight);
\draw[popmuted, thick] (0,0) -- (-\radius,\coneheight);
\draw[popmuted, thick] (0,0) -- (\radius,\coneheight);
% Ellipses for base/top
\draw[popmuted, thick] (0,-\coneheight) ellipse ({\radius} and {0.3});
\draw[popmuted, thick] (0,\coneheight) ellipse ({\radius} and {0.3});
% Axis
\draw[popmuted, dashed] (0,-4.5) -- (0,4.5) node[above] {$\Delta$};
% Summit
\fill[popdark] (0,0) circle (2pt) node[right] {$S$};
% Plane for Ellipse (horizontal-ish)
\begin{scope}[yshift=-1.5cm]
\fill[popblue!20, opacity=0.6] (-2.5,0) -- (2.5,0) -- (2.5,0.4) -- (-2.5,0.4) -- cycle;
\draw[popblue, thick] (-2.5,0.2) -- (2.5,0.2) node[right] {$\mathcal{P}_1$ (Ellipse)};
\draw[popblue, thick] (0,0.2) ellipse ({1.2} and {0.3});
\end{scope}
% Plane for Parabola (parallel to generator)
\begin{scope}[yshift=1.5cm, rotate=-30]
\fill[poporange!20, opacity=0.6] (-2.5,-0.2) -- (2.5,-0.2) -- (2.5,0.2) -- (-2.5,0.2) -- cycle;
\draw[poporange, thick] (-2.5,0) -- (2.5,0) node[right] {$\mathcal{P}_2$ (Parabole)};
% Parabola curve approx
\draw[poporange, thick, domain=-1.5:1.5, smooth, variable=\x] plot ({\x}, {0.3*\x*\x});
\end{scope}
% Plane for Hyperbola (vertical-ish)
\begin{scope}[yshift=3.5cm]
\fill[popgreen!20, opacity=0.6] (-2.5,-0.2) -- (2.5,-0.2) -- (2.5,0.2) -- (-2.5,0.2) -- cycle;
\draw[popgreen, thick] (-2.5,0) -- (2.5,0) node[right] {$\mathcal{P}_3$ (Hyperbole)};
% Hyperbola branches
\draw[popgreen, thick, domain=-2:-0.5, smooth, variable=\x] plot ({\x}, {0.8/\x});
\draw[popgreen, thick, domain=0.5:2, smooth, variable=\x] plot ({\x}, {0.8/\x});
\draw[popgreen, thick, domain=-2:-0.5, smooth, variable=\x] plot ({\x}, {-0.8/\x});
\draw[popgreen, thick, domain=0.5:2, smooth, variable=\x] plot ({\x}, {-0.8/\x});
\end{scope}
% Generator line angle indicator
\draw[popmuted, dashed] (0,0) -- (4,0) node[midway, above] {$\alpha$};
\draw[popmuted, dashed] (0,0) -- (3.464, -2);
\draw (0.5,0) arc (0:-30:0.5);
\node at (0.7,-0.2) {$\alpha$};
\end{tikzpicture}
\legende{Les trois types de coniques obtenus par section d'un cône de révolution par un plan $\mathcal{P}$ : ellipse (plan peu incliné), parabole (plan parallèle à une génératrice), hyperbole (plan très incliné coupant les deux nappes).}
\end{popfigure}

\begin{definition}[Conique — Définition géométrique]
Une \textbf{conique} est la courbe d'intersection d'un cône de révolution à deux nappes par un plan ne passant pas par le sommet du cône.
\begin{itemize}
    \item Si le plan coupe une seule nappe et n'est pas parallèle à une génératrice $\rightarrow$ \textbf{Ellipse} (cas particulier : \textbf{Cercle} si le plan est perpendiculaire à l'axe).
    \item Si le plan est parallèle à une génératrice du cône $\rightarrow$ \textbf{Parabole}.
    \item Si le plan coupe les deux nappes $\rightarrow$ \textbf{Hyperbole}.
\end{itemize}
\end{definition}

\courssub{Définition unifiée par foyer, directrice et excentricité}

Au-delà de la construction solide, les coniques possèdent une définition plane, métrique, qui permet de les étudier dans un repère orthonormé. C'est la définition par \textbf{foyer}, \textbf{directrice} et \textbf{excentricité}.

\begin{definition}[Foyer, Directrice, Excentricité]
Soit $F$ un point (le \cle{foyer}), $\mathcal{D}$ une droite ne passant pas par $F$ (la \cle{directrice}) et $e$ un nombre réel strictement positif (l'\cle{excentricité}).
La conique de foyer $F$, de directrice $\mathcal{D}$ et d'excentricité $e$ est l'ensemble des points $M$ du plan tels que :
\[
MF = e \cdot d(M, \mathcal{D})
\]
où $d(M, \mathcal{D})$ est la distance du point $M$ à la droite $\mathcal{D}$.
\end{definition}

La valeur de $e$ détermine entièrement la nature de la courbe :

\begin{propriete}[Classification par l'excentricité $e$]
\begin{center}
\begin{tabularx}{\linewidth}{|c|X|X|X|}\hline
\textbf{Type} & \textbf{Condition sur $e$} & \textbf{Forme} & \textbf{Position foyer/directrice} \\ \hline
Ellipse & $0 < e < 1$ & Courbe fermée, bornée & Le foyer est \emph{entre} le centre et la directrice \\ \hline
Parabole & $e = 1$ & Courbe ouverte, une branche & Le foyer est à \emph{mi-distance} du sommet et de la directrice \\ \hline
Hyperbole & $e > 1$ & Courbe ouverte, deux branches & Le foyer est \emph{au-delà} du sommet par rapport à la directrice \\ \hline
\end{tabularx}
\end{center}
\end{propriete}

\begin{attention}[Piège : Cercle et Excentricité]
Le cercle est un cas particulier d'ellipse où $e = 0$. Dans la définition foyer-directrice, cela correspond à une directrice « à l'infini ». On ne traite pas le cercle par la définition foyer-directrice standard ($e>0$), mais par sa définition usuelle (ensemble des points équidistants d'un centre).
\end{attention}



\courssub{Équations réduites dans un repère orthonormé}

Pour exploiter algébriquement les coniques, on les place dans un repère orthonormé $(O; \vec{\imath}, \vec{\jmath})$ de manière à ce que leurs axes de symétries coïncident avec les axes du repère. On obtient alors les \textbf{équations réduites} (ou canoniques), qu'il faut connaître par cœur. \emph{Le passage de la définition foyer-directrice à ces équations est une démonstration classique exigible au Probatoire.}

\begin{propriete}[Équations réduites des coniques — À connaître par cœur]
Dans un repère orthonormé bien choisi (centre au centre de symétrie pour ellipse/hyperbole, sommet au sommet pour parabole) :
\begin{enumerate}
    \item \textbf{Ellipse} (centre $O$, axes sur les axes du repère, $a \ge b > 0$) :
    \[
    \frac{x^2}{a^2} + \frac{y^2}{b^2} = 1
    \]
    \item \textbf{Parabole} (sommet $O$, axe focal sur $Ox$, paramètre $p > 0$) :
    \[
    y^2 = 2px \quad \text{(branche vers la droite)}
    \]
    \item \textbf{Hyperbole} (centre $O$, axe focal sur $Ox$, $a>0, b>0$) :
    \[
    \frac{x^2}{a^2} - \frac{y^2}{b^2} = 1
    \]
\end{enumerate}
\end{propriete}

\begin{aretenir}[Paramètres caractéristiques à retenir]
Pour chaque conique, identifiez systématiquement :
\begin{itemize}
    \item \textbf{Ellipse} : $a$ (grand demi-axe), $b$ (petit demi-axe), $c = \sqrt{a^2 - b^2}$ (demi-distance focale), $e = c/a$ (excentricité). Foyers $F(\pm c, 0)$, Sommets $S(\pm a, 0)$, Directrices $x = \pm a/e$.
    \item \textbf{Parabole} : $p$ (paramètre, distance foyer-directrice = $p$, distance sommet-foyer = $p/2$). Foyer $F(p/2, 0)$, Sommet $S(0,0)$, Directrice $x = -p/2$.
    \item \textbf{Hyperbole} : $a$ (demi-axe transverse), $b$ (demi-axe conjugate), $c = \sqrt{a^2 + b^2}$ (demi-distance focale), $e = c/a > 1$. Foyers $F(\pm c, 0)$, Sommets $S(\pm a, 0)$, Directrices $x = \pm a/e$, \textbf{Asymptotes} : $y = \pm \frac{b}{a}x$.
\end{itemize}
\end{aretenir}

\begin{attention}[Ordre des axes : $a$ et $b$]
Pour l'ellipse $\frac{x^2}{a^2}+\frac{y^2}{b^2}=1$, la convention standard impose $a \ge b$. L'axe focal est $Ox$ (le grand axe).
Si l'équation est $\frac{x^2}{4}+\frac{y^2}{9}=1$, alors $a^2=9, b^2=4 \Rightarrow a=3, b=2$. L'axe focal est $Oy$ ! Ne confondez pas la position de $a^2$ (sous la variable de l'axe focal) et la valeur numérique.
\end{attention}

\courssub{Éléments caractéristiques : Foyers, Directrices, Asymptotes}

\courssub{L'Ellipse $\frac{x^2}{a^2}+\frac{y^2}{b^2}=1$ ($a > b > 0$)}

\begin{popfigure}
\begin{tikzpicture}[scale=1.2, >=Stealth, domain=-3:3]
% Axes
\draw[popline] (-3.5,0) -- (3.5,0) node[right] {$x$};
\draw[popline] (0,-2.2) -- (0,2.2) node[above] {$y$};
% Ellipse a=3, b=1.5 -> c=sqrt(9-2.25)=sqrt(6.75)=2.598
\def\a{3} \def\b{1.5} \def\c{2.598}
\draw[popblue, thick] plot[domain=0:360, samples=100] ({\a*cos(\x)}, {\b*sin(\x)});
% Foci
\fill[popdark] (-\c,0) circle (2pt) node[below left] {$F'$};
\fill[popdark] (\c,0) circle (2pt) node[below right] {$F$};
% Vertices
\fill[popblue] (-\a,0) circle (1.5pt) node[below] {$A'$};
\fill[popblue] (\a,0) circle (1.5pt) node[below] {$A$};
\fill[popblue] (0,-\b) circle (1.5pt) node[left] {$B'$};
\fill[popblue] (0,\b) circle (1.5pt) node[left] {$B$};
% Directrices x = +/- a/e = +/- a^2/c
\pgfmathsetmacro{\dd}{\a*\a/\c}
\draw[popmuted, dashed] (\dd,-2.2) -- (\dd,2.2) node[above] {$\mathcal{D}$};
\draw[popmuted, dashed] (-\dd,-2.2) -- (-\dd,2.2) node[above] {$\mathcal{D}'$};
% Labels
\draw[<->] (0,-1.8) -- (\a,-1.8) node[midway, below] {$a$};
\draw[<->] (0,-1.5) -- (\c,-1.5) node[midway, below] {$c$};
\draw[<->, popmuted] (\c,0.2) -- (\dd,0.2) node[midway, above] {$a/e - c$};
\end{tikzpicture}
\legende{Ellipse : Foyers $F, F'$, Sommets $A, A', B, B'$, Directrices $\mathcal{D}, \mathcal{D}'$. Relation fondamentale : $a^2 = b^2 + c^2$ et $e = c/a$.}
\end{popfigure}

\courssub{La Parabole $y^2 = 2px$ ($p > 0$)}

\begin{popfigure}
\begin{tikzpicture}[scale=1.2, >=Stealth]
% Axes
\draw[popline] (-1,0) -- (4,0) node[right] {$x$};
\draw[popline] (0,-3) -- (0,3) node[above] {$y$};
% Parabola y^2 = 2px, p=2 -> x = y^2/4
\def\p{2}
\draw[poporange, thick, domain=-2.8:2.8, smooth, variable=\y] plot ({\y*\y/(2*\p)}, \y);
% Focus
\fill[popdark] (\p/2, 0) circle (2pt) node[below right] {$F(p/2, 0)$};
% Vertex
\fill[poporange] (0,0) circle (1.5pt) node[below left] {$S(0,0)$};
% Directrix x = -p/2
\draw[popmuted, dashed] (-\p/2,-3) -- (-\p/2,3) node[above] {$\mathcal{D}: x=-p/2$};
% Latus rectum (focal chord)
\draw[popmuted, thin] (\p/2, -\p) -- (\p/2, \p);
\draw[<->, popmuted] (0, -1.5) -- (\p/2, -1.5) node[midway, below] {$p/2$};
\draw[<->, popmuted] (-\p/2, -1.5) -- (0, -1.5) node[midway, below] {$p/2$};
\end{tikzpicture}
\legende{Parabole : Foyer $F$, Sommet $S$, Directrice $\mathcal{D}$. Le paramètre $p$ est la distance foyer-directrice. La corde focale (latus rectum) a pour longueur $2p$.}
\end{popfigure}

\courssub{L'Hyperbole $\frac{x^2}{a^2}-\frac{y^2}{b^2}=1$ ($a, b > 0$)}

\begin{popfigure}
\begin{tikzpicture}[scale=1.2, >=Stealth]
% Axes
\draw[popline] (-4.5,0) -- (4.5,0) node[right] {$x$};
\draw[popline] (0,-3) -- (0,3) node[above] {$y$};
% Hyperbola a=2, b=1.5 -> c=2.5
\def\a{2} \def\b{1.5} \def\c{2.5}
% Asymptotes
\draw[popmuted, dashed] (-4.5, -\b/\a*4.5) -- (4.5, \b/\a*4.5);
\draw[popmuted, dashed] (-4.5, \b/\a*4.5) -- (4.5, -\b/\a*4.5);
\node[popmuted] at (4.2, 3.2) {$y=\frac{b}{a}x$};
\node[popmuted] at (4.2, -3.2) {$y=-\frac{b}{a}x$};
% Branches
\draw[popgreen, thick, domain=-4.5:-\a, samples=100] plot (\x, {\b*sqrt(max(0,\x*\x/\a/\a - 1))});
\draw[popgreen, thick, domain=-4.5:-\a, samples=100] plot (\x, {-\b*sqrt(max(0,\x*\x/\a/\a - 1))});
\draw[popgreen, thick, domain=\a:4.5, samples=100] plot (\x, {\b*sqrt(max(0,\x*\x/\a/\a - 1))});
\draw[popgreen, thick, domain=\a:4.5, samples=100] plot (\x, {-\b*sqrt(max(0,\x*\x/\a/\a - 1))});
% Foci
\fill[popdark] (-\c,0) circle (2pt) node[below left] {$F'$};
\fill[popdark] (\c,0) circle (2pt) node[below right] {$F$};
% Vertices
\fill[popgreen] (-\a,0) circle (1.5pt) node[below] {$A'$};
\fill[popgreen] (\a,0) circle (1.5pt) node[below] {$A$};
% Directrices
\pgfmathsetmacro{\dd}{\a*\a/\c}
\draw[popmuted, dashed] (\dd,-3) -- (\dd,3) node[above] {$\mathcal{D}$};
\draw[popmuted, dashed] (-\dd,-3) -- (-\dd,3) node[above] {$\mathcal{D}'$};
% Rectangle asymptote guide
\draw[popmuted, very thin] (-\a,-\b) rectangle (\a,\b);
\draw[<->] (0,-2.5) -- (\a,-2.5) node[midway, below] {$a$};
\draw[<->] (0,-2.2) -- (\c,-2.2) node[midway, below] {$c$};
\end{tikzpicture}
\legende{Hyperbole : Foyers $F, F'$, Sommets $A, A'$, Directrices $\mathcal{D}, \mathcal{D}'$, Asymptotes (diagonales du rectangle $[-a,a]\times[-b,b]$). Relation : $c^2 = a^2 + b^2$, $e = c/a > 1$.}
\end{popfigure}

\courssub{Propriétés optiques (réflexion) — Énoncés}

Ces propriétés expliquent l'usage industriel des coniques (antennes, phares, fours solaires, billard). Elles sont admises au Probatoire (pas de démonstration exigible), mais il faut savoir les énoncer et les utiliser.

\begin{propriete}[Propriétés optiques des coniques]
\begin{enumerate}
    \item \textbf{Ellipse :} Tout rayon lumineux émis par un foyer $F$ et réfléchi par la paroi de l'ellipse passe par l'autre foyer $F'$. (Angles d'incidence et de réflexion égaux par rapport à la normale, qui bissecte l'angle $\widehat{FMF'}$).
    \item \textbf{Parabole :} Tout rayon lumineux émis par le foyer $F$ et réfléchi par la parabole en ressort \textbf{parallèle à l'axe} de la parabole. Réciproquement, des rayons parallèles à l'axe incidents sur la parabole convergent vers le foyer $F$.
    \item \textbf{Hyperbole :} Tout rayon lumineux émis par un foyer $F$ et réfléchi par une branche de l'hyperbole semble provenir de l'autre foyer $F'$ (réflexion vers l'extérieur de la branche). Un rayon dirigé vers un foyer $F'$ et incident sur la branche associée à $F$ est réfléchi vers $F$.
\end{enumerate}
\end{propriete}

\begin{savaistu}[Kepler et les orbites planétaires]
Au début du XVII\textsuperscript{e} siècle, Johannes Kepler découvre que les planètes ne tournent pas en cercles parfaits autour du Soleil, mais sur des \textbf{ellipses} dont le Soleil occupe un foyer (1\textsuperscript{re} loi de Kepler, 1609). C'est la première application physique majeure des coniques. Plus tard, Newton montrera que la gravitation universelle entraîne nécessairement des trajectoires coniques (ellipses pour les planètes, paraboles/hyperboles pour certaines comètes ou sondes spatiales en survol). Au Cameroun, les antennes paraboliques (réception satellite) et les réflecteurs de phares de voiture sont des applications quotidiennes de la propriété optique de la parabole.
\end{savaistu}

\courssub{Identifier une conique à partir de son équation cartésienne générale}

En technologie, on rencontre souvent l'équation générale du second degré :
\[
ax^2 + bxy + cy^2 + dx + ey + f = 0 \quad (a,b,c \text{ non tous nuls})
\]
Sans la réduire complètement, le \textbf{discriminant} $\Delta = b^2 - 4ac$ permet de connaître la nature de la conique (si elle n'est pas dégénérée : vide, point, droites).

\begin{methode}[Identifier le type de conique par le discriminant $\Delta$]
Donnée l'équation $ax^2 + bxy + cy^2 + dx + ey + f = 0$ :
\begin{enumerate}
    \item Calculer $\Delta = b^2 - 4ac$.
    \item Conclure :
    \begin{itemize}
        \item $\Delta < 0$ $\rightarrow$ \textbf{Ellipse} (ou cercle si $a=c$ et $b=0$, ou ensemble vide/point).
        \item $\Delta = 0$ $\rightarrow$ \textbf{Parabole} (ou ensemble vide, droite, deux droites parallèles).
        \item $\Delta > 0$ $\rightarrow$ \textbf{Hyperbole} (ou deux droites sécantes).
    \end{itemize}
    \item (Optionnel) Vérifier la non-dégénérescence par le déterminant de la matrice $3\times3$ associée si le problème l'exige.
\end{enumerate}
\end{methode}

\begin{exemplebox}[Application du discriminant]
Déterminer la nature des coniques :
\begin{enumerate}
    \item $3x^2 + 2xy + 3y^2 - 4 = 0$ : $a=3, b=2, c=3 \Rightarrow \Delta = 4 - 36 = -32 < 0$ $\rightarrow$ \textbf{Ellipse}.
    \item $x^2 - 2xy + y^2 + x - y = 0$ : $a=1, b=-2, c=1 \Rightarrow \Delta = 4 - 4 = 0$ $\rightarrow$ \textbf{Parabole} (ici dégénérée en deux droites parallèles : $(x-y)^2+(x-y)=0 \Leftrightarrow (x-y)(x-y+1)=0$).
    \item $2x^2 - 5xy - 3y^2 + 4 = 0$ : $a=2, b=-5, c=-3 \Rightarrow \Delta = 25 + 24 = 49 > 0$ $\rightarrow$ \textbf{Hyperbole}.
\end{enumerate}
\end{exemplebox}

\courssub{Exemple résolu complet : Ellipse $9x^2+25y^2=225$}

\begin{exoresolu}[Analyse complète d'une ellipse]
{Donner l'équation réduite, les coordonnées des foyers, les équations des directrices et l'excentricité de l'ellipse d'équation $9x^2 + 25y^2 = 225$.}

\tcblower
\textbf{1. Mise sous forme réduite.}
On divise par 225 :
\[
\frac{9x^2}{225} + \frac{25y^2}{225} = 1 \iff \frac{x^2}{25} + \frac{y^2}{9} = 1
\]
On identifie $a^2 = 25$ et $b^2 = 9$. Comme $a^2 > b^2$, l'axe focal est $Ox$.
\[
a = 5, \quad b = 3.
\]

\textbf{2. Calcul de $c$ (demi-distance focale).}
Pour l'ellipse : $c^2 = a^2 - b^2 = 25 - 9 = 16 \Rightarrow c = 4$.

\textbf{3. Coordonnées des foyers.}
L'axe focal est $Ox$, les foyers sont sur l'axe des abscisses :
\[
F(-c, 0) = F(-4, 0) \quad \text{et} \quad F'(c, 0) = F'(4, 0).
\]

\textbf{4. Excentricité $e$.}
\[
e = \frac{c}{a} = \frac{4}{5} = 0,8.
\]
Vérification : $0 < e < 1$, c'est bien une ellipse.

\textbf{5. Équations des directrices.}
Les directrices sont perpendiculaires à l'axe focal ($Ox$), d'équations $x = \pm \frac{a}{e} = \pm \frac{a^2}{c}$.
\[
\frac{a}{e} = \frac{5}{0,8} = 6,25 \quad \text{ou} \quad \frac{a^2}{c} = \frac{25}{4} = 6,25.
\]
Directrices : $\mathcal{D} : x = 6,25$ et $\mathcal{D}' : x = -6,25$.

\textbf{Réponse :} Équation réduite $\frac{x^2}{25}+\frac{y^2}{9}=1$. Foyers $F(-4,0), F'(4,0)$. Excentricité $e=0,8$. Directrices $x=\pm 6,25$.
\end{exoresolu}

\begin{attention}[Erreurs fréquentes à ne pas commettre]
\begin{itemize}
    \item \textbf{Inverser $a$ et $b$ :} Toujours prendre $a^2$ comme le plus grand dénominateur pour l'ellipse. Si l'équation est $\frac{x^2}{9}+\frac{y^2}{25}=1$, alors $a=5$ (sur $y$), l'axe focal est $Oy$, les foyers sont $(0, \pm 4)$.
    \item \textbf{Confondre $c$ ellipse et $c$ hyperbole :} Ellipse : $c^2 = a^2 - b^2$. Hyperbole : $c^2 = a^2 + b^2$.
    \item \textbf{Oublier le $2$ dans la parabole :} $y^2 = 2px$, pas $y^2 = px$. Le foyer est à $(\frac{p}{2}, 0)$, pas $(p, 0)$.
    \item \textbf{Asymptotes de l'hyperbole :} Équations $y = \pm \frac{b}{a}x$ (si axe focal $Ox$). Ne pas écrire $y = \pm \frac{a}{b}x$.
\end{itemize}
\end{attention}

\courssub{Exercices d'application}

\begin{exercice}[Exercice 1 : Identification et paramètres]
On considère la conique $\mathcal{C}$ d'équation $4x^2 + 9y^2 - 16x + 18y - 11 = 0$.
\begin{enumerate}
    \item Déterminer la nature de $\mathcal{C}$ par le discriminant.
    \item Réduire l'équation pour trouver le centre, les axes, $a$, $b$, $c$, $e$.
    \item Donner les coordonnées des foyers et les équations des directrices.
    \item Tracer $\mathcal{C}$ dans un repère (indiquer les éléments caractéristiques).
\end{enumerate}
\end{exercice}

\begin{exercice}[Exercice 2 : Parabole et propriété optique]
Une antenne parabolique a une ouverture de $80$ cm de diamètre et une profondeur de $10$ cm au centre. On place son sommet à l'origine et son axe sur $Ox$ ($x \ge 0$).
\begin{enumerate}
    \item Déterminer l'équation réduite de la parabole.
    \item Calculer la position du foyer $F$ (emplacement de la tête de réception).
    \item Vérifier que la propriété optique garantit que tous les signaux incidents parallèles à l'axe convergent vers $F$.
\end{enumerate}
\end{exercice}

\begin{exercice}[Exercice 3 : Hyperbole et asymptotes]
Soit l'hyperbole $\mathcal{H}$ d'équation $4x^2 - y^2 = 16$.
\begin{enumerate}
    \item Mettre sous forme réduite. Préciser $a, b, c, e$.
    \item Déterminer les équations des asymptotes.
    \item On considère le point $M(4, 4\sqrt{3})$ sur $\mathcal{H}$. Calculer les distances $MF$ et $MF'$ ($F, F'$ foyers). Vérifier la relation $|MF' - MF| = 2a$.
\end{enumerate}
\end{exercice}

\begin{corrige}[Exercice 1]
\begin{enumerate}
    \item $a=4, b=0, c=9 \Rightarrow \Delta = 0 - 144 = -144 < 0 \rightarrow$ \textbf{Ellipse}.
    \item $4(x^2 - 4x) + 9(y^2 + 2y) = 11 \iff 4[(x-2)^2 - 4] + 9[(y+1)^2 - 1] = 11 \iff 4(x-2)^2 + 9(y+1)^2 = 36$.
    Équation réduite : $\frac{(x-2)^2}{9} + \frac{(y+1)^2}{4} = 1$.
    Centre $\Omega(2, -1)$. $a=3$ (axe $Ox$), $b=2$. $c = \sqrt{9-4} = \sqrt{5}$. $e = \frac{\sqrt{5}}{3}$.
    \item Foyers : $F(2-\sqrt{5}, -1)$ et $F'(2+\sqrt{5}, -1)$.
    Directrices : $x = 2 \pm \frac{a}{e} = 2 \pm \frac{9}{\sqrt{5}}$.
    \item Tracé : Repère centré en $\Omega$, rectangle $[-3,3]\times[-2,2]$, ellipse inscrite.
\end{enumerate}
\end{corrige}

\begin{corrige}[Exercice 2]
\begin{enumerate}
    \item Diamètre 80 cm $\rightarrow$ rayon 40 cm. Point $M(10, 40)$ sur la parabole $y^2 = 2px$.
    $40^2 = 2p \times 10 \Rightarrow 1600 = 20p \Rightarrow p = 80$ cm.
    Équation : $y^2 = 160x$ ($x$ en cm).
    \item Foyer $F(\frac{p}{2}, 0) = F(40, 0)$. La tête de réception est à 40 cm du sommet (au fond de l'antenne).
    \item Propriété optique : Tout rayon incident parallèle à l'axe ($Ox$) est réfléchi vers le foyer $F$. Les ondes satellites (quasi parallèles) sont donc concentrées en $F$.
\end{enumerate}
\end{corrige}

\begin{corrige}[Exercice 3]
\begin{enumerate}
    \item $\frac{x^2}{4} - \frac{y^2}{16} = 1$. $a=2, b=4$. $c = \sqrt{4+16} = \sqrt{20} = 2\sqrt{5}$. $e = \frac{c}{a} = \sqrt{5} > 1$.
    \item Asymptotes : $y = \pm \frac{b}{a}x = \pm 2x$.
    \item Foyers : $F(-2\sqrt{5}, 0)$ et $F'(2\sqrt{5}, 0)$.
    Point $M(4, 4\sqrt{3})$ : vérification $\frac{16}{4} - \frac{48}{16} = 4 - 3 = 1$. $M \in \mathcal{H}$ (branche droite).
    \begin{align*}
    MF &= \sqrt{(4 + 2\sqrt{5})^2 + (4\sqrt{3})^2} = \sqrt{16 + 20 + 16\sqrt{5} + 48} = \sqrt{84 + 16\sqrt{5}} \\
    MF' &= \sqrt{(4 - 2\sqrt{5})^2 + (4\sqrt{3})^2} = \sqrt{16 + 20 - 16\sqrt{5} + 48} = \sqrt{84 - 16\sqrt{5}}
    \end{align*}
    Sur la branche droite, $MF' > MF$. Calculons $(MF' - MF)^2$ :
    \begin{align*}
    (MF' - MF)^2 &= MF'^2 + MF^2 - 2\sqrt{MF'^2 \cdot MF^2} \\
    &= (84 - 16\sqrt{5}) + (84 + 16\sqrt{5}) - 2\sqrt{84^2 - (16\sqrt{5})^2} \\
    &= 168 - 2\sqrt{7056 - 1280} \\
    &= 168 - 2\sqrt{5776} \\
    &= 168 - 2 \times 76 = 168 - 152 = 16.
    \end{align*}
    Donc $MF' - MF = 4 = 2a$. \textbf{Propriété vérifiée.}
\end{enumerate}
\end{corrige}

\coursec{Tangente en un point à une conique}

Dans la section précédente, nous avons présenté les trois familles de coniques — ellipse, parabole et hyperbole — à travers leurs équations réduites, leurs foyers et leurs propriétés géométriques. Nous savons maintenant reconnaître une conique et tracer son allure. Mais dans la pratique des métiers techniques, une question revient sans cesse : comment une courbe se comporte-t-elle \emph{localement}, au voisinage d'un point précis ? Un opticien qui taille une lentille, un technicien qui installe une antenne parabolique, un mécanicien qui profile une pièce : tous ont besoin de connaître la \cle{tangente} à la courbe en un point donné. C'est l'objet de cette section.

\courssub{Intuition : qu'est-ce qu'une tangente à une conique ?}

\begin{definition}[Tangente à une conique]
Soit $\mathcal{C}$ une conique (ellipse, parabole ou hyperbole) et $M_0$ un point de $\mathcal{C}$. La \cle{tangente} à $\mathcal{C}$ en $M_0$ est la droite passant par $M_0$ qui « épouse » la courbe au voisinage de ce point : elle touche la conique en $M_0$ sans la traverser localement. C'est la position limite des sécantes $(M_0M)$ lorsque le point $M$ de la conique se rapproche de $M_0$.
\end{definition}

\begin{attention}[Un point capital]
Dans toute cette section, le point $M_0(x_0\,;\,y_0)$ où l'on cherche la tangente doit \textbf{appartenir à la conique}. Les formules que nous allons établir ne sont valables que sous cette condition. Avant d'appliquer une formule, il faut \textbf{toujours vérifier} que les coordonnées de $M_0$ satisfont l'équation de la conique.
\end{attention}

\courssub{Équations des tangentes : les trois théorèmes}

\begin{propriete}[Tangente à l'ellipse]
Soit $\mathcal{E}$ l'ellipse d'équation réduite $\dfrac{x^2}{a^2}+\dfrac{y^2}{b^2}=1$ et $M_0(x_0\,;\,y_0)$ un point de $\mathcal{E}$. La tangente à $\mathcal{E}$ en $M_0$ a pour équation :
\[ \frac{x\,x_0}{a^2}+\frac{y\,y_0}{b^2}=1. \]
\emph{(Démonstration exigible au Probatoire, voir activité guidée ci-dessous.)}
\end{propriete}

\begin{propriete}[Tangente à la parabole]
Soit $\mathcal{P}$ la parabole d'équation réduite $y^2=2px$ et $M_0(x_0\,;\,y_0)$ un point de $\mathcal{P}$. La tangente à $\mathcal{P}$ en $M_0$ a pour équation :
\[ y\,y_0 = p\,(x+x_0). \]
\end{propriete}

\begin{propriete}[Tangente à l'hyperbole]
Soit $\mathcal{H}$ l'hyperbole d'équation réduite $\dfrac{x^2}{a^2}-\dfrac{y^2}{b^2}=1$ et $M_0(x_0\,;\,y_0)$ un point de $\mathcal{H}$. La tangente à $\mathcal{H}$ en $M_0$ a pour équation :
\[ \frac{x\,x_0}{a^2}-\frac{y\,y_0}{b^2}=1. \]
\end{propriete}

\begin{aretenir}[Règle de dédoublement — moyen mnémotechnique]
On obtient l'équation de la tangente en « dédoublant » l'équation de la conique :
\begin{itemize}
\item $x^2$ devient $x\,x_0$ et $y^2$ devient $y\,y_0$ ;
\item $2x$ devient $x+x_0$ et $2y$ devient $y+y_0$ (le coefficient $p$ reste inchangé).
\end{itemize}
Ainsi $\frac{x^2}{a^2}+\frac{y^2}{b^2}=1$ donne $\frac{x x_0}{a^2}+\frac{y y_0}{b^2}=1$, et $y^2=2px$ donne $y y_0 = p(x+x_0)$.
\end{aretenir}

\courssub{Démonstration guidée : le cas de l'ellipse}

La démonstration qui suit utilise la \cle{dérivation implicite} : on considère $y$ comme une fonction de $x$ et on dérive les deux membres de l'équation de la conique. Cette méthode est transposable à la parabole et à l'hyperbole.

\begin{experience}[Démonstration : tangente à l'ellipse]
On considère l'ellipse $\mathcal{E}$ d'équation $\dfrac{x^2}{a^2}+\dfrac{y^2}{b^2}=1$ et un point $M_0(x_0\,;\,y_0)$ de $\mathcal{E}$, avec $y_0\neq 0$.

\textbf{Étape 1 : dérivation implicite.} Au voisinage de $M_0$, l'ellipse est le graphe d'une fonction $y$ de $x$. En dérivant les deux membres de $\dfrac{x^2}{a^2}+\dfrac{y^2}{b^2}=1$ par rapport à $x$ :
\[ \frac{2x}{a^2}+\frac{2y\,y'}{b^2}=0. \]

\textbf{Étape 2 : coefficient directeur de la tangente.} On isole $y'$ :
\[ y'=-\frac{b^2\,x}{a^2\,y}. \]
Au point $M_0(x_0\,;\,y_0)$, le coefficient directeur de la tangente vaut donc $m=-\dfrac{b^2 x_0}{a^2 y_0}$.

\textbf{Étape 3 : équation de la tangente.} La tangente passe par $M_0$ avec le coefficient directeur $m$ :
\[ y-y_0=-\frac{b^2 x_0}{a^2 y_0}\,(x-x_0). \]
En multipliant les deux membres par $\dfrac{y_0}{b^2}$ :
\[ \frac{y\,y_0}{b^2}-\frac{y_0^2}{b^2}=-\frac{x\,x_0}{a^2}+\frac{x_0^2}{a^2}, \]
soit
\[ \frac{x\,x_0}{a^2}+\frac{y\,y_0}{b^2}=\frac{x_0^2}{a^2}+\frac{y_0^2}{b^2}. \]

\textbf{Étape 4 : utilisation de l'appartenance.} Comme $M_0\in\mathcal{E}$, on a $\dfrac{x_0^2}{a^2}+\dfrac{y_0^2}{b^2}=1$. On conclut :
\[ \frac{x\,x_0}{a^2}+\frac{y\,y_0}{b^2}=1. \]
Le cas $y_0=0$ (sommets $(\pm a\,;\,0)$) donne les tangentes verticales $x=\pm a$, que la formule fournit aussi directement. $\blacksquare$
\end{experience}

\courssub{Méthode pratique et exemples résolus}

\begin{methode}[Déterminer l'équation d'une tangente en 3 étapes]
\begin{enumerate}
\item \textbf{Vérifier} que le point $M_0(x_0\,;\,y_0)$ appartient à la conique, en remplaçant ses coordonnées dans l'équation.
\item \textbf{Écrire} la formule de dédoublement adaptée à la conique (ellipse, parabole ou hyperbole).
\item \textbf{Simplifier} l'équation obtenue pour la mettre sous une forme réduite $y=mx+p$ si demandé.
\end{enumerate}
\end{methode}

\begin{exoresolu}[Tangente à une hyperbole]
Déterminer l'équation de la tangente à l'hyperbole $\mathcal{H}$ d'équation $\dfrac{x^2}{4}-\dfrac{y^2}{9}=1$ au point $M_0\left(4\,;\,3\sqrt{3}\right)$.
\tcblower
\textbf{Étape 1 : appartenance.} On vérifie :
\[ \frac{4^2}{4}-\frac{(3\sqrt{3})^2}{9}=\frac{16}{4}-\frac{27}{9}=4-3=1. \]
Donc $M_0\in\mathcal{H}$ : la formule s'applique.

\textbf{Étape 2 : formule de dédoublement.} Ici $a^2=4$, $b^2=9$, $x_0=4$, $y_0=3\sqrt{3}$ :
\[ \frac{x\times 4}{4}-\frac{y\times 3\sqrt{3}}{9}=1
\quad\Longleftrightarrow\quad
x-\frac{\sqrt{3}}{3}\,y=1. \]

\textbf{Étape 3 : forme réduite.} On isole $y$ :
\[ y=\sqrt{3}\,(x-1)=\sqrt{3}\,x-\sqrt{3}. \]
La tangente cherchée est la droite d'équation $y=\sqrt{3}\,x-\sqrt{3}$, de coefficient directeur $\sqrt{3}$.
\end{exoresolu}

\begin{exoresolu}[Tangente à une parabole]
Déterminer la tangente à la parabole $y^2=8x$ au point $M_0(2\,;\,4)$.
\tcblower
\textbf{Étape 1 :} $4^2=16$ et $8\times 2=16$, donc $M_0$ appartient bien à la parabole.

\textbf{Étape 2 :} l'équation est de la forme $y^2=2px$ avec $2p=8$, donc $p=4$. La formule $y\,y_0=p(x+x_0)$ donne :
\[ 4y=4(x+2). \]

\textbf{Étape 3 :} en simplifiant par $4$ : $y=x+2$. La tangente est la droite d'équation $y=x+2$.
\end{exoresolu}

\begin{attention}[Erreurs fréquentes]
\begin{itemize}
\item \textbf{Appliquer la formule à un point extérieur.} Si $M_0\notin\mathcal{C}$, la droite obtenue par dédoublement n'est \emph{pas} une tangente : vérifier l'appartenance est obligatoire.
\item \textbf{Oublier le signe moins} pour l'hyperbole : $\frac{x x_0}{a^2}-\frac{y y_0}{b^2}=1$.
\item \textbf{Mal identifier $p$} pour la parabole : dans $y^2=2px$, c'est bien $2p$ qui est le coefficient de $x$. Pour $y^2=8x$, on a $p=4$ et non $p=8$.
\item \textbf{Confondre $x_0$ et $x$} : dans la formule, $x_0$ et $y_0$ sont des constantes (coordonnées du point de contact), $x$ et $y$ sont les variables.
\end{itemize}
\end{attention}

\courssub{Propriétés de réflexion des coniques}

Les tangentes aux coniques possèdent des propriétés géométriques remarquables, à la base de nombreuses applications technologiques.

\begin{propriete}[Propriété de réflexion de l'ellipse]
Soit $\mathcal{E}$ une ellipse de foyers $F$ et $F'$, et $M_0$ un point de $\mathcal{E}$. La tangente en $M_0$ est la bissectrice extérieure de l'angle $\widehat{FM_0F'}$ : un rayon issu du foyer $F$ qui se réfléchit sur l'ellipse en $M_0$ passe ensuite par l'autre foyer $F'$.
\end{propriete}

\begin{propriete}[Propriété de réflexion de la parabole]
Soit $\mathcal{P}$ une parabole de foyer $F$ et d'axe $(\Delta)$. Tout rayon issu du foyer $F$ se réfléchit sur la parabole en un rayon \textbf{parallèle à l'axe}. Réciproquement, tout rayon parallèle à l'axe se réfléchit en passant par le foyer.
\end{propriete}

\begin{popfigure}
\begin{center}
\begin{tikzpicture}[scale=1.1]
% Ellipse a=3, b=2 ; foyers en (±sqrt(5),0)
\draw[popline,->] (-4,0) -- (4,0) node[below left]{$x$};
\draw[popline,->] (0,-2.8) -- (0,2.8) node[below left]{$y$};
\draw[popblue,thick] (0,0) ellipse (3 and 2);
% Point M0 sur l'ellipse
\coordinate (M) at (1.5,1.732);
% Tangente : x*1.5/9 + y*1.732/4 = 1
\draw[poporange,thick] (-1.2,3.3) -- (4.2,-0.35) node[right,popdark]{tangente};
% Foyers
\fill[popink] (2.236,0) circle (1.5pt) node[below right]{$F$};
\fill[popink] (-2.236,0) circle (1.5pt) node[below left]{$F'$};
\fill[poporange] (M) circle (2pt) node[above right]{$M_0$};
% Rayons réfléchis
\draw[popgreen,thick,->] (-2.236,0) -- (M);
\draw[popgreen,thick,->] (M) -- (2.236,0);
\end{tikzpicture}
\end{center}
\legende{Ellipse $\frac{x^2}{9}+\frac{y^2}{4}=1$, tangente en $M_0$ et propriété de réflexion : le rayon issu de $F'$ se réfléchit vers $F$.}
\end{popfigure}

\begin{popfigure}
\begin{center}
\begin{tikzpicture}[scale=0.9]
% Parabole y^2 = 2px avec p=1 : x = y^2/2
\draw[popline,->] (-0.5,0) -- (5.5,0) node[below left]{$x$ (axe)};
\draw[popline,->] (0,-3) -- (0,3) node[below left]{$y$};
\draw[popblue,thick,domain=-2.6:2.6,samples=100] plot ({0.5*\x*\x},{\x});
% Foyer F(p/2,0) = (0.5,0)
\fill[popink] (0.5,0) circle (2pt) node[below right]{$F$};
% Points de réflexion
\coordinate (A) at (2,2);
\coordinate (B) at (1.125,-1.5);
\coordinate (C) at (3.125,2.5);
% Rayons issus du foyer
\draw[popgreen,thick,->] (0.5,0) -- (A);
\draw[popgreen,thick,->] (0.5,0) -- (B);
\draw[popgreen,thick,->] (0.5,0) -- (C);
% Rayons réfléchis parallèles à l'axe
\draw[poporange,thick,->] (A) -- (5.2,2);
\draw[poporange,thick,->] (B) -- (5.2,-1.5);
\draw[poporange,thick,->] (C) -- (5.2,2.5);
% Tangentes (pointillés)
\draw[popmuted,dashed] (0.5,3) -- (3.5,1);
\draw[popmuted,dashed] (0,-2.25) -- (2.25,-0.75);
\end{tikzpicture}
\end{center}
\legende{Propriété optique de la parabole : tout rayon issu du foyer $F$ se réfléchit parallèlement à l'axe (réflecteur parabolique).}
\end{popfigure}

\begin{savaistu}[Les antennes paraboliques]
Les antennes paraboliques que l'on voit sur les toits de Douala ou de Yaoundé exploitent exactement la propriété de réflexion de la parabole : les ondes émises par le satellite arrivent parallèles à l'axe de la parabole, se réfléchissent sur la surface et convergent toutes au foyer, où est placé le récepteur (la tête LNB). À l'inverse, les phares de voiture et les lampes torches placent la source lumineuse au foyer : les rayons ressortent alors en un faisceau parallèle, puissant et dirigé. Les fours solaires et les radiotélescopes fonctionnent sur le même principe !
\end{savaistu}

\begin{exercice}[À vous de jouer]
\begin{enumerate}
\item Déterminer l'équation de la tangente à l'ellipse $\dfrac{x^2}{25}+\dfrac{y^2}{9}=1$ au point $M_0\left(3\,;\,\dfrac{12}{5}\right)$.
\item Déterminer l'équation de la tangente à la parabole $y^2=4x$ au point d'abscisse $x_0=1$ et d'ordonnée positive.
\end{enumerate}
\end{exercice}

\begin{corrige}[Corrigé]
\textbf{1.} Vérification : $\dfrac{9}{25}+\dfrac{144/25}{9}=\dfrac{9}{25}+\dfrac{16}{25}=1$, donc $M_0$ appartient à l'ellipse. La formule donne $\dfrac{3x}{25}+\dfrac{12y/5}{9}=1$, soit $\dfrac{3x}{25}+\dfrac{4y}{15}=1$. En multipliant par $75$ : $9x+20y=75$.

\textbf{2.} Pour $x_0=1$ : $y_0^2=4$, donc $y_0=2$ (ordonnée positive). Ici $2p=4$ donc $p=2$. La tangente a pour équation $2y=2(x+1)$, c'est-à-dire $y=x+1$.
\end{corrige}

\begin{aretenir}[L'essentiel de la section]
\begin{itemize}
\item Tangente à l'ellipse $\frac{x^2}{a^2}+\frac{y^2}{b^2}=1$ en $M_0$ : $\dfrac{x x_0}{a^2}+\dfrac{y y_0}{b^2}=1$.
\item Tangente à la parabole $y^2=2px$ en $M_0$ : $y y_0=p(x+x_0)$.
\item Tangente à l'hyperbole $\frac{x^2}{a^2}-\frac{y^2}{b^2}=1$ en $M_0$ : $\dfrac{x x_0}{a^2}-\dfrac{y y_0}{b^2}=1$.
\item Condition d'application : $M_0$ doit appartenir à la conique.
\item Propriété de réflexion : un rayon issu d'un foyer d'une ellipse ressort vers l'autre foyer ; un rayon issu du foyer d'une parabole ressort parallèle à l'axe.
\end{itemize}
\end{aretenir}

\coursec{Construction d'une conique et applications intégrées}

Dans les sections précédentes, nous avons appris à déterminer l'équation de la tangente en un point d'une conique à partir de son équation réduite. Nous savons donc maintenant \emph{calculer} sur les coniques. Mais une question naturelle se pose dans l'atelier ou sur le chantier : comment \emph{tracer réellement} une ellipse, une parabole ou une hyperbole ? Un menuisier qui doit découper une table ovale, un technicien qui fabrique un réflecteur parabolique, un ingénieur qui modélise une trajectoire : tous ont besoin de méthodes de \cle{construction} pratiques. C'est l'objet de cette dernière section, qui se conclura par un beau pont entre géométrie et probabilités : les \cle{probabilités géométriques}.

\courssub{Construction de l'ellipse : la méthode du jardinier}

\textbf{Intuition.} Rappelons la définition bifocale de l'ellipse : étant donnés deux points fixes $F_1$ et $F_2$ (les \cle{foyers}) et un nombre réel $2a > F_1F_2$, l'ellipse est l'ensemble des points $M$ tels que
\[ MF_1 + MF_2 = 2a. \]
Autrement dit, la \emph{somme} des distances d'un point de l'ellipse aux deux foyers est constante. Cette propriété fournit immédiatement un procédé mécanique de tracé : il suffit d'une ficelle de longueur constante dont les extrémités sont fixées aux foyers.

\begin{methode}[Construire une ellipse connaissant ses foyers et la longueur $2a$ (méthode du jardinier)]
\begin{enumerate}
\item Fixer deux punaises (ou deux piquets) aux points $F_1$ et $F_2$, distants de $2c$, avec $2c < 2a$.
\item Prendre une ficelle de longueur exactement égale à $2a$ et attacher ses deux extrémités aux punaises.
\item Tendre la ficelle avec la pointe d'un crayon (ou d'un piquet) placé en un point $M$.
\item Faire glisser le crayon en gardant la ficelle toujours tendue : la pointe décrit l'ellipse de foyers $F_1$, $F_2$ et de grand axe $2a$.
\end{enumerate}
\textbf{Justification :} à chaque position, la ficelle est pliée en deux segments $MF_1$ et $MF_2$ dont la somme des longueurs vaut $2a$, donc $M$ vérifie $MF_1+MF_2=2a$ : c'est exactement la définition de l'ellipse.
\end{methode}

\begin{popfigure}
\begin{center}
\begin{tikzpicture}[scale=1.1]
% Ellipse a=3, c=1.8, b=sqrt(9-3.24)=2.4
\draw[popblue, very thick] (0,0) ellipse (3 and 2.4);
\coordinate (F1) at (-1.8,0);
\coordinate (F2) at (1.8,0);
\coordinate (M) at (1.8,1.92);
% ficelle
\draw[poporange, very thick] (F1) -- (M) -- (F2);
% axes
\draw[popmuted, dashed] (-3.6,0) -- (3.6,0);
\draw[popmuted, dashed] (0,-2.9) -- (0,2.9);
% points
\fill[popink] (F1) circle (2pt) node[below left=1pt] {$F_1$};
\fill[popink] (F2) circle (2pt) node[below right=1pt] {$F_2$};
\fill[popdark] (M) circle (2.5pt) node[above right=1pt] {$M$ (crayon)};
\fill[popink] (0,0) circle (1.5pt) node[below left=1pt] {$O$};
% annotations
\node[poporange] at (-1.6,1.5) {ficelle};
\node[popblue] at (3.1,1.3) {ellipse};
\draw[<->, popmuted] (-3,-2.7) -- (3,-2.7) node[midway, below] {$2a$};
\end{tikzpicture}
\end{center}
\legende{Méthode du jardinier : la ficelle $F_1MF_2$ a une longueur constante $2a$ ; en la tendant avec un crayon, on trace l'ellipse.}
\end{popfigure}

\begin{exemplebox}[Tracer une ellipse de table]
Un ébéniste veut tracer un plateau ovale de \SI{1,20}{m} de long et \SI{0,80}{m} de large. On a $2a = 1{,}20$~m et $2b = 0{,}80$~m, donc $a = 0{,}60$~m et $b = 0{,}40$~m. La distance focale est
\[ c = \sqrt{a^2-b^2} = \sqrt{0{,}36-0{,}16} = \sqrt{0{,}20} \approx 0{,}447~\text{m}. \]
Il plante deux clous distants de $2c \approx \SI{0,89}{m}$ sur l'axe de longueur, attache une ficelle de \SI{1,20}{m} et trace : le plateau est parfaitement elliptique.
\end{exemplebox}

\begin{attention}[Conditions d'application de la méthode du jardinier]
\begin{itemize}
\item Il faut impérativement $2a > F_1F_2$ (c'est-à-dire $a > c$). Si la ficelle est trop courte (longueur égale à $F_1F_2$), le crayon ne peut glisser que sur le segment $[F_1F_2]$ : l'ellipse dégénère en un segment.
\item La ficelle ne doit pas être élastique, sinon la somme $MF_1+MF_2$ varie et la courbe tracée n'est plus une ellipse.
\item Plus les punaises sont proches ($c$ petit), plus l'ellipse ressemble à un cercle ; le cas $c=0$ (deux punaises confondues) donne un cercle de rayon $a$.
\end{itemize}
\end{attention}

\courssub{Construction de la parabole : foyer et directrice}

\textbf{Intuition.} La parabole est l'ensemble des points $M$ \cle{équidistants} d'un point fixe $F$ (le \cle{foyer}) et d'une droite fixe $(d)$ (la \cle{directrice}) :
\[ MF = \operatorname{dist}(M,(d)). \]
Cette définition se traduit en un tracé point par point à la règle et à l'équerre.

\begin{methode}[Tracer une parabole point par point (règle et équerre)]
On se donne le foyer $F$ et la directrice $(d)$. Pour obtenir des points de la parabole :
\begin{enumerate}
\item Tracer une droite $(\Delta)$ parallèle à la directrice $(d)$, à une distance $h$ de celle-ci (du côté du foyer).
\item Tout point $M$ de la parabole situé sur $(\Delta)$ vérifie $\operatorname{dist}(M,(d)) = h$, donc doit vérifier $MF = h$.
\item Tracer le cercle de centre $F$ et de rayon $h$ : ses intersections avec $(\Delta)$ donnent deux points de la parabole (s'ils existent).
\item Recommencer avec plusieurs valeurs de $h$, puis relier les points obtenus à main levée par une courbe régulière.
\end{enumerate}
Le point situé à égale distance de $F$ et de $(d)$ sur la perpendiculaire commune est le \cle{sommet} $S$ de la parabole, milieu du segment joignant $F$ à sa projection sur $(d)$.
\end{methode}

\begin{exemplebox}[Parabole de paramètre $p$]
Soit la parabole d'équation réduite $y^2 = 4x$. On a $2p = 4$, donc $p = 2$.
\begin{itemize}
\item Foyer $F(\frac{p}{2}\,;\,0) = F(1\,;\,0)$.
\item Directrice $(d) : x = -\frac{p}{2} = -1$.
\item Sommet $S(0\,;\,0)$.
\end{itemize}
\textbf{Test d'un point n'appartenant pas à la parabole :} $M(2\,;\,2)$.
$MF = \sqrt{(2-1)^2 + 2^2} = \sqrt{5}$.
$\operatorname{dist}(M,(d)) = |2 - (-1)| = 3$.
$\sqrt{5} \neq 3$, donc $M$ n'est pas sur la parabole (et effectivement $2^2 = 4 \neq 4 \times 2 = 8$).

\textbf{Test d'un point appartenant à la parabole :} $M(1\,;\,2)$ (car $2^2 = 4 \times 1$).
$MF = \sqrt{(1-1)^2 + 2^2} = 2$.
$\operatorname{dist}(M,(d)) = |1 - (-1)| = 2$.
L'équidistance est vérifiée : $M$ appartient bien à la parabole.
\end{exemplebox}

\begin{attention}[Piège fréquent]
La distance d'un point $M$ à la directrice est la distance \emph{perpendiculaire} : c'est la longueur $MH$ où $H$ est le projeté orthogonal de $M$ sur $(d)$, et non la distance de $M$ à un point quelconque de la droite. Dans les calculs, pour $(d) : x = -p/2$ et $M(x\,;\,y)$, on a $\operatorname{dist}(M,(d)) = \left|x+\dfrac{p}{2}\right|$.
\end{attention}

\courssub{Construction de l'hyperbole : différence constante des distances}

\textbf{Intuition.} L'hyperbole est l'ensemble des points $M$ tels que la \emph{différence} des distances aux deux foyers soit constante en valeur absolue :
\[ |MF_1 - MF_2| = 2a, \qquad \text{avec } 2a < F_1F_2 = 2c. \]
La condition $2a < 2c$ est ici inversée par rapport à l'ellipse : c'est elle qui garantit l'existence de points et la forme en deux branches.

\begin{methode}[Tracer une hyperbole point par point (règle et compas)]
\begin{enumerate}
\item Placer les foyers $F_1$ et $F_2$ avec $F_1F_2 = 2c$, et choisir $2a < 2c$.
\item Choisir un rayon $r > c - a$. Tracer le cercle de centre $F_1$ et de rayon $r$, puis le cercle de centre $F_2$ et de rayon $r + 2a$.
\item Les points d'intersection $M$ de ces deux cercles vérifient $MF_2 - MF_1 = 2a$ : ce sont des points de la branche la plus proche de $F_1$.
\item En échangeant les rôles de $F_1$ et $F_2$ (rayons $r$ et $r+2a$), on obtient des points de l'autre branche.
\item Faire varier $r$ et relier les points : on obtient les deux branches de l'hyperbole.
\end{enumerate}
\end{methode}

\begin{exemplebox}[Hyperbole de foyers distants de 10 cm]
Prenons $F_1F_2 = 2c = 10$~cm ($c=5$) et $2a = 6$~cm ($a=3$). La condition $c-a = 2$ est respectée.
Pour $r = 4$~cm : cercle de centre $F_1$ de rayon $4$ et cercle de centre $F_2$ de rayon $4+6=10$.
Les intersections $M$ vérifient $MF_2 - MF_1 = 10 - 4 = 6 = 2a$ : elles appartiennent à l'hyperbole.
Les sommets sont sur la droite des foyers, à distance $a=3$ du centre, et $b = \sqrt{c^2-a^2} = \sqrt{25-9} = 4$.
\end{exemplebox}

\begin{attention}[La construction donne une approximation]
Les constructions géométriques à la ficelle, à l'équerre ou au compas ne donnent qu'une \emph{approximation} de la conique : épaisseur du trait, élasticité de la ficelle, nombre fini de points tracés, imprécision des instruments. Seule l'\cle{équation analytique} (équation réduite $\dfrac{x^2}{a^2}+\dfrac{y^2}{b^2}=1$, $y^2=2px$, $\dfrac{x^2}{a^2}-\dfrac{y^2}{b^2}=1$) décrit la conique de manière \emph{exacte}. En technologie, on trace à la ficelle pour l'atelier, mais on calcule avec l'équation pour la conception (commande numérique, découpe laser).
\end{attention}

\courssub{Problèmes concrets modélisés par des coniques}

\textbf{Trajectoire d'un projectile : la parabole.} Un objet lancé (ballon, jet d'eau, obus) suit, en l'absence de frottements, une trajectoire parabolique. Si un ouvrier lance un sac de ciment depuis le sol avec une vitesse initiale $v_0$ sous un angle $\alpha$, la hauteur en fonction de la distance horizontale s'écrit $y = -\dfrac{g}{2v_0^2\cos^2\alpha}\,x^2 + x\tan\alpha$ : c'est bien une parabole d'axe vertical, dont le sommet correspond à la hauteur maximale.

\textbf{Miroir elliptique.} Une propriété remarquable de l'ellipse : tout rayon lumineux (ou sonore) émis depuis un foyer $F_1$ se réfléchit sur l'ellipse et passe par l'autre foyer $F_2$. On exploite cela pour concentrer de la lumière ou des ondes.

\textbf{Cible elliptique et jeu de hasard.} Un forain propose de lancer une fléchette au hasard sur un panneau : la probabilité de gagner dépend du rapport entre l'aire de la cible et l'aire du panneau. C'est l'objet du paragraphe suivant.

\begin{savaistu}[Les miroirs elliptiques en médecine]
La propriété de réflexion de l'ellipse est utilisée dans les \emph{lampes scialytiques} des blocs opératoires : une source lumineuse placée au foyer $F_1$ d'un réflecteur elliptique voit sa lumière concentrée au foyer $F_2$, exactement là où se trouve la zone à opérer. Le même principe sert dans les appareils de lithotripsie, qui brisent les calculs rénaux par des ondes de choc émises en $F_1$ et convergentes en $F_2$, sans incision. Les coniques sauvent littéralement des vies !
\end{savaistu}

\courssub{Probabilités géométriques : le lien avec les coniques}

\textbf{Intuition.} Dans les sections précédentes, nous avons calculé des probabilités avec des ensembles finis (lancers de dés, tirages de cartes) grâce à la formule $\dfrac{\text{cas favorables}}{\text{cas possibles}}$. Mais que faire lorsque l'expérience consiste à choisir un point \emph{au hasard} dans une région du plan (une fléchette sur une cible, une goutte de pluie sur une dalle) ? Il y a une infinité de positions possibles. L'idée naturelle : si le point est réparti de façon \cle{uniforme}, la probabilité de tomber dans une sous-région est proportionnelle à son \emph{aire}.

\begin{propriete}[Probabilité géométrique (distribution uniforme)]
Soit $\Omega$ une région du plan d'aire $\mathcal{A}(\Omega)$ non nulle, et $E$ une sous-région de $\Omega$ d'aire $\mathcal{A}(E)$. Si un point est choisi au hasard uniformément dans $\Omega$, alors
\[ P(E) = \dfrac{\mathcal{A}(E)}{\mathcal{A}(\Omega)}. \]
En particulier, pour une conique $E$ inscrite dans une région $\Omega$ : $P(\text{tomber dans la conique}) = \dfrac{\mathcal{A}(\text{conique})}{\mathcal{A}(\text{région contenant})}$.
\end{propriete}

\begin{aretenir}[Aires à connaître]
\begin{itemize}
\item Aire d'une ellipse de demi-axes $a$ et $b$ : $\mathcal{A} = \pi a b$.
\item Aire d'un rectangle de dimensions $L$ et $\ell$ : $\mathcal{A} = L\ell$.
\item Aire d'un disque de rayon $r$ : $\mathcal{A} = \pi r^2$ (cas particulier de l'ellipse avec $a=b=r$).
\end{itemize}
\end{aretenir}

\begin{popfigure}
\begin{center}
\begin{tikzpicture}[scale=2.2]
% rectangle 2 x 1.4, ellipse a=1 b=0.7 inscrite
\fill[popblueL] (-1,-0.7) rectangle (1,0.7);
\fill[pattern=north east lines, pattern color=poporange] (0,0) ellipse (1 and 0.7);
\draw[popink, thick] (-1,-0.7) rectangle (1,0.7);
\draw[poporange, very thick] (0,0) ellipse (1 and 0.7);
\draw[<->, popmuted] (-1,-0.95) -- (1,-0.95) node[midway, below] {$L$};
\draw[<->, popmuted] (1.25,-0.7) -- (1.25,0.7) node[midway, right] {$\ell$};
\draw[<->, popmuted] (0,0.15) -- (1,0.15) node[midway, above] {$a$};
\draw[<->, popmuted] (-0.15,0) -- (-0.15,0.7) node[midway, left] {$b$};
\node[poporange] at (0.55,-0.35) {cible};
\node[popink] at (-0.7,-0.5) {panneau};
\end{tikzpicture}
\end{center}
\legende{Ellipse de demi-axes $a$ et $b$ inscrite dans un rectangle $L\times\ell$. La zone hachurée représente l'événement « la fléchette touche la cible » : $P = \dfrac{\pi a b}{L\ell}$.}
\end{popfigure}

\begin{exoresolu}[Cible elliptique sur un panneau rectangulaire]
\textbf{Énoncé.} Une cible elliptique de demi-grand axe $a = \SI{30}{cm}$ et de demi-petit axe $b = \SI{20}{cm}$ est fixée sur un panneau rectangulaire de $\SI{1}{m}\times\SI{1}{m}$. Un joueur lance une fléchette qui atteint le panneau en un point uniformément réparti. Quelle est la probabilité qu'elle touche la cible ? On donnera la valeur exacte puis une valeur approchée à $10^{-3}$ près.
\tcblower
\textbf{Solution détaillée.}
\begin{enumerate}
\item \textbf{Aire de la cible (ellipse).} On applique la formule $\mathcal{A}_{\text{cible}} = \pi a b$ :
\[ \mathcal{A}_{\text{cible}} = \pi \times 30 \times 20 = 600\pi~\text{cm}^2. \]
\item \textbf{Aire du panneau (rectangle).} Le panneau mesure $\SI{1}{m}\times\SI{1}{m}$, soit $100~\text{cm}\times 100~\text{cm}$ :
\[ \mathcal{A}_{\text{panneau}} = 100 \times 100 = 10\,000~\text{cm}^2. \]
\item \textbf{Probabilité géométrique.} La répartition étant uniforme :
\[ P = \dfrac{\mathcal{A}_{\text{cible}}}{\mathcal{A}_{\text{panneau}}} = \dfrac{600\pi}{10\,000} = \dfrac{3\pi}{50}. \]
Numériquement : $P \approx \dfrac{3\times 3{,}1416}{50} \approx 0{,}188$.
\end{enumerate}
\textbf{Conclusion.} La probabilité de toucher la cible est $\boxed{P = \dfrac{3\pi}{50} \approx 0{,}188}$, soit environ $18{,}8\,\%$. Le forain garde donc environ $81\,\%$ de chances que le joueur perde !
\end{exoresolu}

\begin{attention}[Erreurs fréquentes en probabilité géométrique]
\begin{itemize}
\item \textbf{Unités incohérentes} : avant de faire le quotient des aires, convertir toutes les longueurs dans la \emph{même} unité (ici, tout en cm). Oublier la conversion $1~\text{m} = 100$~cm donne un résultat 10\,000 fois trop grand.
\item \textbf{Confondre $a$ et $2a$} : dans la formule $\pi a b$, $a$ et $b$ sont les \emph{demi}-axes. Si l'énoncé donne les axes complets (grand axe $=60$~cm), il faut diviser par 2.
\item \textbf{Condition d'uniformité} : la formule n'est valable que si le point est réparti uniformément. Un joueur entraîné qui vise le centre ne suit pas ce modèle !
\end{itemize}
\end{attention}

\begin{exercice}[Trajectoire et cible]
Un jet d'eau sortant d'une fontaine suit une trajectoire parabolique d'équation $y = -\dfrac{1}{5}x^2 + 2x$ (en mètres). 
\begin{enumerate}
\item Déterminer la hauteur maximale atteinte par le jet (sommet de la parabole).
\item À quelle distance du point de départ le jet retombe-t-il au sol ($y=0$) ?
\item Une bassine circulaire de rayon $0{,}5$~m est placée au point de chute, sur une dalle carrée de $4$~m de côté. Une goutte tombant uniformément sur la dalle a quelle probabilité de tomber dans la bassine ?
\end{enumerate}
\end{exercice}

\begin{corrige}[Exercice : Trajectoire et cible]
\begin{enumerate}
\item La parabole $y = -\dfrac{1}{5}x^2+2x$ a son sommet d'abscisse $x_S = -\dfrac{b}{2a} = -\dfrac{2}{2\times(-1/5)} = 5$. Alors $y_S = -\dfrac{25}{5}+10 = 5$~m. Hauteur maximale : $5$~m.
\item $y=0 \iff x\left(-\dfrac{x}{5}+2\right)=0 \iff x=0$ ou $x=10$. Le jet retombe à $10$~m.
\item Aire de la bassine : $\pi \times 0{,}5^2 = 0{,}25\pi$~m$^2$. Aire de la dalle : $16$~m$^2$. Probabilité : $P = \dfrac{0{,}25\pi}{16} = \dfrac{\pi}{64} \approx 0{,}049$, soit environ $4{,}9\,\%$.
\end{enumerate}
\end{corrige}

\begin{savaistu}[Buffon et le calcul de $\pi$ par le hasard]
Au XVIII\textsuperscript{e} siècle, le comte de Buffon a montré qu'en laissant tomber une aiguille au hasard sur un parquet à lattes parallèles, la probabilité que l'aiguille coupe une rainure fait intervenir $\pi$ : en répétant l'expérience des milliers de fois, on peut \emph{estimer} $\pi$ ! C'est l'ancêtre des « méthodes de Monte-Carlo », utilisées aujourd'hui en ingénierie et en finance pour calculer des aires et des probabilités par simulation aléatoire — exactement l'esprit de notre cible elliptique.
\end{savaistu}

\begin{aretenir}[Bilan de la section]
\begin{itemize}
\item \textbf{Ellipse} : ficelle de longueur $2a$ fixée aux deux foyers (méthode du jardinier) ; condition $2a > 2c$.
\item \textbf{Parabole} : points équidistants du foyer $F$ et de la directrice $(d)$, tracés à la règle et à l'équerre.
\item \textbf{Hyperbole} : points vérifiant $|MF_1 - MF_2| = 2a$, tracés à la règle et au compas ; condition $2a < 2c$.
\item Les constructions instrumentales approchent la conique ; seule l'équation réduite la décrit exactement.
\item \textbf{Probabilité géométrique} (répartition uniforme) : $P(E) = \dfrac{\mathcal{A}(E)}{\mathcal{A}(\Omega)}$ ; aire de l'ellipse : $\pi a b$.
\end{itemize}
\end{aretenir}

\newpage

\coursec{Activité d'intégration}

\courssub{Objectifs de l'activité}

À l'issue de cette activité d'intégration, l'élève sera capable de :
\begin{itemize}
\item déterminer l'équation réduite d'une parabole à partir de trois points de sa trajectoire ;
\item reconnaître un schéma de Bernoulli et utiliser la loi binomiale pour calculer des probabilités ;
\item représenter graphiquement une conique et l'histogramme d'une loi binomiale ;
\item mobiliser conjointement les savoirs « coniques » et « probabilités » dans une situation concrète de la vie courante ;
\item rédiger et justifier une conclusion argumentée à partir de résultats mathématiques.
\end{itemize}

\courssub{Situation}

Lors du championnat inter-lycées de basket-ball organisé au gymnase de Yaoundé, l'équipe du Lycée Technique de Nkol-Eton s'entraîne pour la finale. Le coach, ancien élève de la filière Génie Mécanique, veut utiliser les mathématiques pour analyser les performances de son meilleur tireur, Étienne.

À l'aide d'une caméra et d'un logiciel de pointage, le coach a relevé la trajectoire idéale du ballon lors d'un lancer réussi. Dans un repère orthonormé $(O\,;\vec{i},\vec{j})$ dont l'origine est le point de départ du ballon (les mains du joueur) et dont l'unité est le mètre, la trajectoire du ballon est assimilée à un arc de \cle{parabole} passant par les trois points suivants :
\begin{itemize}
\item le point de lancer $A(0\,;0)$ ;
\item le sommet de la trajectoire $S(3\,;4)$ (point le plus haut atteint par le ballon) ;
\item le centre du panier $P(6\,;3)$.
\end{itemize}

Par ailleurs, les statistiques de la saison montrent qu'Étienne réussit $70\,\%$ de ses lancers francs. Pendant la séance d'entraînement, il effectue une série de $10$ lancers francs, supposés indépendants les uns des autres.

Le coach pose alors deux questions à ses joueurs, tous élèves de Terminale technique :
\begin{enumerate}
\item Quelle est l'équation de la trajectoire idéale du ballon ?
\item Quelle est la probabilité qu'Étienne marque exactement $7$ paniers sur ses $10$ lancers ? Au moins $8$ paniers ?
\end{enumerate}

\courssub{Consignes}

Travail en groupes de 3 à 4 élèves. Matériel : papier millimétré, calculatrice, règle, crayons de couleur. Durée : 2 séances. Chaque groupe affichera ses résultats au tableau et les présentera à la classe.

\textbf{Partie A : modélisation géométrique de la trajectoire}

\begin{enumerate}
\item On cherche la trajectoire sous la forme $y = ax^2 + bx + c$. En utilisant le fait que la parabole passe par $A(0\,;0)$, déterminer la valeur de $c$.
\item Écrire les deux équations obtenues en traduisant l'appartenance des points $S(3\,;4)$ et $P(6\,;3)$ à la parabole.
\item Résoudre le système obtenu et donner l'équation réduite de la parabole.
\item Vérifier que le sommet de la parabole obtenue est voisin du point $S(3\,;4)$ relevé expérimentalement (on pourra calculer l'abscisse du sommet $-\dfrac{b}{2a}$).
\item Sur le papier millimétré, représenter la parabole pour $x$ compris entre $0$ et $6$, puis placer le panier $P$.
\end{enumerate}

\textbf{Partie B : étude probabiliste de la série de lancers}

\begin{enumerate}
\setcounter{enumi}{5}
\item Justifier que le nombre $X$ de paniers marqués sur les $10$ lancers suit une loi binomiale dont on précisera les paramètres.
\item Calculer la probabilité qu'Étienne marque exactement $7$ paniers : $P(X = 7)$.
\item Calculer la probabilité qu'il marque au moins $8$ paniers : $P(X \geq 8)$.
\item Calculer l'espérance mathématique $E(X)$ et interpréter le résultat.
\item Construire l'histogramme de la loi binomiale $\mathcal{B}(10\,;0{,}7)$ pour les valeurs de $k$ de $4$ à $10$ (on donnera les probabilités arrondies au millième).
\end{enumerate}

\textbf{Partie C : synthèse}

\begin{enumerate}
\setcounter{enumi}{10}
\item Rédiger une conclusion argumentée (10 lignes maximum) à destination du coach : la trajectoire est-elle cohérente avec un lancer réussi ? Étienne est-il un tireur fiable ? Faut-il lui confier le dernier lancer de la finale ?
\end{enumerate}

\courssub{Ressources à mobiliser}

\begin{aretenir}[Boîte à outils de l'activité]
\begin{itemize}
\item \textbf{Parabole} d'équation $y = ax^2 + bx + c$ ($a \neq 0$) : sommet d'abscisse $x_S = -\dfrac{b}{2a}$ ; axe de symétrie vertical.
\item \textbf{Schéma de Bernoulli} : répétition de $n$ épreuves identiques et indépendantes à deux issues (succès de probabilité $p$, échec de probabilité $q = 1-p$).
\item \textbf{Loi binomiale} $\mathcal{B}(n\,;p)$ : $P(X = k) = \dbinom{n}{k} p^k (1-p)^{n-k}$.
\item \textbf{Espérance} : $E(X) = np$.
\item \textbf{Coefficient binomial} : $\dbinom{n}{k} = \dfrac{n!}{k!(n-k)!}$.
\end{itemize}
\end{aretenir}

\courssub{Démarche et corrigé commenté}

\begin{exoresolu}[Le panier probabiliste : corrigé complet]
Énoncé : déterminer l'équation de la trajectoire, étudier la loi de $X$, représenter les graphiques et conclure.
\tcblower
\textbf{Partie A. Équation de la parabole.}

\textbf{Étape 1.} La parabole passe par $A(0\,;0)$ : en remplaçant $x$ par $0$ et $y$ par $0$ dans $y = ax^2 + bx + c$, on obtient $c = 0$. L'équation s'écrit donc $y = ax^2 + bx$.

\textbf{Étape 2.} Passage par $S(3\,;4)$ et $P(6\,;3)$ :
\[
\begin{cases}
9a + 3b = 4 \\
36a + 6b = 3
\end{cases}
\]

\textbf{Étape 3.} Multiplions la première équation par $2$ : $18a + 6b = 8$. Retranchons-la de la seconde :
\[
(36a + 6b) - (18a + 6b) = 3 - 8 \iff 18a = -5 \iff a = -\frac{5}{18}.
\]
En reportant dans $9a + 3b = 4$ : $-\dfrac{45}{18} + 3b = 4$, soit $3b = 4 + \dfrac{5}{2} = \dfrac{13}{2}$, d'où $b = \dfrac{13}{6}$.

L'équation réduite de la trajectoire est :
\[
y = -\frac{5}{18}x^2 + \frac{13}{6}x
\]
\emph{Vérification :} pour $x = 3$ : $y = -\dfrac{5}{18}\times 9 + \dfrac{13}{6}\times 3 = -\dfrac{5}{2} + \dfrac{13}{2} = 4$ ; pour $x = 6$ : $y = -\dfrac{5}{18}\times 36 + \dfrac{13}{6}\times 6 = -10 + 13 = 3$. La parabole passe bien par $S$ et $P$.

\textbf{Étape 4.} Abscisse du sommet : $x_S = -\dfrac{b}{2a} = -\dfrac{13/6}{2 \times (-5/18)} = \dfrac{13/6}{5/9} = \dfrac{13}{6} \times \dfrac{9}{5} = \dfrac{117}{30} = 3{,}9$.

L'ordonnée du sommet est $y(3{,}9) = -\dfrac{5}{18}\times (3{,}9)^2 + \dfrac{13}{6}\times 3{,}9 = -4{,}225 + 8{,}45 = 4{,}225$.

\emph{Remarque :} le sommet exact de la parabole est $(3{,}9\,;\ 4{,}225)$, légèrement différent du point $S(3\,;4)$ relevé par le logiciel : c'est normal, les relevés expérimentaux comportent de petites imprécisions. La parabole passe bien par les trois points donnés et son sommet est voisin de $S$, ce qui valide le modèle.

\textbf{Étape 5.} Tableau de valeurs pour le tracé :
\begin{center}
\begin{tabular}{|c|ccccccc|}
\hline
$x$ & $0$ & $1$ & $2$ & $3$ & $4$ & $5$ & $6$ \\
\hline
$y$ & $0$ & $1{,}89$ & $3{,}22$ & $4$ & $4{,}22$ & $3{,}89$ & $3$ \\
\hline
\end{tabular}
\end{center}

\begin{popfigure}
\begin{tikzpicture}[scale=0.85]
\draw[->,popline] (-0.3,0) -- (7,0) node[below]{$x$ (m)};
\draw[->,popline] (0,-0.3) -- (0,5) node[left]{$y$ (m)};
\draw[popblue,thick,domain=0:6,samples=100] plot (\x,{-5/18*\x*\x + 13/6*\x});
\fill[poporange] (0,0) circle (2pt) node[below left]{$A$};
\fill[popblue] (3,4) circle (2pt) node[above left]{$S$};
\fill[poppink] (6,3) circle (2pt) node[right]{$P$ (panier)};
\draw[dashed,popmuted] (3.9,0) -- (3.9,4.225);
\fill[popgreen] (3.9,4.225) circle (2pt) node[above right]{sommet};
\end{tikzpicture}
\legende{Trajectoire parabolique du ballon, du point de lancer $A$ au panier $P$.}
\end{popfigure}

\textbf{Partie B. Loi binomiale.}

\textbf{Étape 6.} Chaque lancer est une épreuve de Bernoulli : succès « panier marqué » de probabilité $p = 0{,}7$, échec de probabilité $q = 1 - p = 0{,}3$. Les $10$ lancers sont identiques et indépendants : c'est un schéma de Bernoulli. Donc $X$ suit la loi binomiale $\mathcal{B}(10\,;\ 0{,}7)$.

\textbf{Étape 7.} Probabilité de marquer exactement $7$ paniers :
\begin{align*}
P(X = 7) &= \binom{10}{7}(0{,}7)^7(0{,}3)^3 \\
&= 120 \times 0{,}0823543 \times 0{,}027 \\
&\approx 0{,}267.
\end{align*}

\textbf{Étape 8.} Probabilité de marquer au moins $8$ paniers :
\begin{align*}
P(X \geq 8) &= P(X=8) + P(X=9) + P(X=10) \\
&= \binom{10}{8}(0{,}7)^8(0{,}3)^2 + \binom{10}{9}(0{,}7)^9(0{,}3) + (0{,}7)^{10} \\
&\approx 0{,}2335 + 0{,}1211 + 0{,}0282 \\
&\approx 0{,}383.
\end{align*}

\textbf{Étape 9.} Espérance : $E(X) = np = 10 \times 0{,}7 = 7$. En moyenne, Étienne marque $7$ paniers sur $10$ lancers.

\textbf{Étape 10.} Tableau des probabilités (arrondies au millième) :
\begin{center}
\begin{tabular}{|c|ccccccc|}
\hline
$k$ & $4$ & $5$ & $6$ & $7$ & $8$ & $9$ & $10$ \\
\hline
$P(X=k)$ & $0{,}037$ & $0{,}103$ & $0{,}200$ & $0{,}267$ & $0{,}233$ & $0{,}121$ & $0{,}028$ \\
\hline
\end{tabular}
\end{center}

\begin{popfigure}
\begin{tikzpicture}[yscale=8,xscale=0.85]
\draw[->,popline] (3.5,0) -- (10.5,0) node[below]{$k$};
\draw[->,popline] (3.5,0) -- (3.5,0.32) node[left]{$P(X=k)$};
\foreach \x/\h in {4/0.037,5/0.103,6/0.200,7/0.267,8/0.233,9/0.121,10/0.028}
  \fill[popblue!70] (\x-0.3,0) rectangle (\x+0.3,\h);
\foreach \x in {4,5,6,7,8,9,10} \node[below] at (\x,0) {\x};
\draw[popline] (3.5,0.1) node[left]{$0{,}1$} -- (3.6,0.1);
\draw[popline] (3.5,0.2) node[left]{$0{,}2$} -- (3.6,0.2);
\draw[popline] (3.5,0.3) node[left]{$0{,}3$} -- (3.6,0.3);
\end{tikzpicture}
\legende{Histogramme de la loi binomiale $\mathcal{B}(10\,;\,0{,}7)$ : la valeur la plus probable est $k = 7$.}
\end{popfigure}

\textbf{Partie C. Conclusion (exemple de rédaction attendue).}

La trajectoire du ballon est bien modélisée par la parabole d'équation $y = -\dfrac{5}{18}x^2 + \dfrac{13}{6}x$, qui passe par le point de lancer, culmine à environ $4{,}2$ m et retombe exactement au centre du panier : le modèle géométrique est cohérent avec un lancer réussi. Côté probabilités, Étienne marque en moyenne $7$ paniers sur $10$, la probabilité de marquer exactement $7$ paniers est d'environ $0{,}267$ et celle d'en marquer au moins $8$ est d'environ $0{,}383$. Étienne est donc un tireur fiable : le coach peut lui confier le dernier lancer de la finale, tout en sachant qu'un échec reste possible (environ $3$ chances sur $10$ à chaque tir).
\end{exoresolu}

\begin{attention}[Erreurs fréquentes à signaler pendant l'activité]
\begin{itemize}
\item Oublier que le passage par l'origine impose $c = 0$ et se lancer dans un système de trois équations à trois inconnues.
\item Confondre $\dbinom{10}{7}$ avec $10^7$ : rappeler que $\dbinom{10}{7} = \dbinom{10}{3} = 120$.
\item Calculer $P(X \geq 8)$ en oubliant un des trois termes $P(X=8)$, $P(X=9)$, $P(X=10)$.
\item Oublier de vérifier l'hypothèse d'\cle{indépendance} des lancers avant d'annoncer la loi binomiale : c'est une condition d'application exigée dans la justification.
\item Dans les graphiques sur papier millimétré, mal choisir l'échelle des ordonnées de l'histogramme (les probabilités sont toutes inférieures à $0{,}3$).
\end{itemize}
\end{attention}

\courssub{Bilan de l'activité}

Cette activité a permis de mettre en évidence que :
\begin{itemize}
\item une \cle{conique} (ici la parabole) n'est pas un objet abstrait : elle modélise des trajectoires réelles (ballon, jet d'eau, projecteur) et son équation se détermine à partir de points expérimentaux ;
\item les \cle{probabilités} et la \cle{loi binomiale} permettent de quantifier la fiabilité d'un sportif, d'une machine ou d'un procédé de production : savoir qu'un tireur réussit $70\,\%$ de ses lancers ne suffit pas, il faut savoir calculer ce que cela implique sur une série de tirs ;
\item l'espérance $E(X) = np$ donne une valeur moyenne de référence, tandis que l'histogramme montre la dispersion des résultats autour de cette moyenne ;
\item les mathématiques de Terminale technique fournissent des outils de décision : la conclusion rédigée à destination du coach illustre le savoir-faire « rédiger et justifier une démarche, une réponse ou une conclusion » exigé au Baccalauréat technique.
\end{itemize}

\begin{savaistu}[Les maths dans le sport camerounais]
L'analyse de données sportives (\emph{data scouting}) se développe rapidement : les trajectoires de ballons sont modélisées par des coniques et les performances des joueurs par des lois de probabilité. Au Cameroun, les fédérations de basket-ball et de football commencent à recruter des techniciens capables de croiser géométrie et statistiques : un débouché concret pour les titulaires d'un Baccalauréat technique !
\end{savaistu}

\newpage

\coursec{Méthodes et exercices résolus}

\coursec{Probabilités et Coniques}

\courssub{Vocabulaire et expériences aléatoires}

\begin{definition}[Expérience aléatoire]
Une \cle{expérience aléatoire} est une expérience qui remplit trois conditions :
\begin{enumerate}
    \item Elle admet plusieurs résultats possibles.
    \item On ne peut pas prédire avec certitude lequel se produira.
    \item Elle est reproductible à l'identique (mêmes conditions initiales).
\end{enumerate}
\end{definition}

\begin{definition}[Univers et événements]
L'\cle{univers} $\Omega$ est l'ensemble de \emph{tous} les résultats élémentaires possibles. Un \cle{événement} est une partie de $\Omega$ (sous-ensemble). Un événement \cle{élémentaire} ne contient qu'un seul résultat.
\end{definition}

\begin{propriete}[Opérations sur les événements]
Soient $A$ et $B$ deux événements de $\Omega$.
\begin{itemize}
    \item \textbf{Union} $A \cup B$ : « $A$ ou $B$ » (au moins l'un des deux se réalise).
    \item \textbf{Intersection} $A \cap B$ : « $A$ et $B$ » (les deux se réalisent simultanément).
    \item \textbf{Complémentaire} $\overline{A}$ ou $A^c$ : « non $A$ » ($A$ ne se réalise pas).
    \item \textbf{Événements incompatibles} : $A \cap B = \varnothing$ (ils ne peuvent pas se réaliser ensemble).
    \item \textbf{Événements contraires} : $A \cap B = \varnothing$ et $A \cup B = \Omega$ (l'un se réalise si et seulement si l'autre ne se réalise pas). Notation : $B = \overline{A}$.
    \item \textbf{Partition} : une famille $(A_i)_{1 \le i \le n}$ d'événements deux à deux incompatibles telle que $\bigcup_{i=1}^n A_i = \Omega$.
\end{itemize}
\end{propriete}

\begin{definition}[Probabilité — Équiprobabilité]
Si $\Omega$ est fini et que tous les résultats élémentaires ont la même chance de se produire (\cle{équiprobabilité}), alors pour tout événement $A \subset \Omega$ :
\[ P(A) = \frac{\text{Card}(A)}{\text{Card}(\Omega)}. \]
En général, une \cle{probabilité} est une application $P : \mathcal{P}(\Omega) \to [0,1]$ telle que $P(\Omega)=1$ et $\sigma$-additive (au programme : additivité finie pour événements incompatibles).
\end{definition}

\begin{propriete}[Propriétés fondamentales]
\begin{enumerate}
    \item $P(\varnothing) = 0$ ; $P(\overline{A}) = 1 - P(A)$.
    \item Si $A \subset B$ alors $P(A) \le P(B)$.
    \item $P(A \cup B) = P(A) + P(B) - P(A \cap B)$.
    \item Si $A$ et $B$ incompatibles : $P(A \cup B) = P(A) + P(B)$.
\end{enumerate}
\end{propriete}

\begin{definition}[Indépendance]
Deux événements $A$ et $B$ (avec $P(A)>0, P(B)>0$) sont \cle{indépendants} si et seulement si :
\[ P(A \cap B) = P(A) \times P(B). \]
Équivalent à $P(A|B) = P(A)$ ou $P(B|A) = P(B)$.
\end{definition}

\begin{exemplebox}[Lancer de deux dés]
$\Omega = \{(1,1), (1,2), \dots, (6,6)\}$, $\text{Card}(\Omega)=36$. Équiprobabilité.
$A$ : « somme = 7 » $\to A = \{(1,6),(2,5),(3,4),(4,3),(5,2),(6,1)\}$, $P(A)=6/36=1/6$.
$B$ : « premier dé = 4 » $\to \text{Card}(B)=6$, $P(B)=1/6$.
$A \cap B = \{(4,3)\}$, $P(A \cap B)=1/36 = P(A)P(B) \Rightarrow A,B$ indépendants.
\end{exemplebox}

\courssub{Probabilité d'un événement et probabilités conditionnelles}

\begin{definition}[Probabilité conditionnelle]
Soient $A, B$ deux événements avec $P(B) > 0$. La \cle{probabilité conditionnelle} de $A$ sachant $B$ est :
\[ P(A|B) = \frac{P(A \cap B)}{P(B)}. \]
C'est la probabilité de $A$ dans l'univers réduit à $B$.
\end{definition}

\begin{propriete}[Formule des probabilités composées]
\[ P(A \cap B) = P(B) \times P(A|B) = P(A) \times P(B|A). \]
Généralisation à $n$ événements : $P(A_1 \cap \dots \cap A_n) = P(A_1) P(A_2|A_1) \dots P(A_n | A_1 \cap \dots \cap A_{n-1})$.
\end{propriete}

\begin{propriete}[Loi des probabilités totales]
Si $(B_i)_{1 \le i \le n}$ est une partition de $\Omega$ avec $P(B_i) > 0$, alors pour tout événement $A$ :
\[ P(A) = \sum_{i=1}^n P(B_i) P(A|B_i). \]
\end{propriete}

\begin{propriete}[Formule de Bayes]
Dans les mêmes conditions :
\[ P(B_i|A) = \frac{P(B_i) P(A|B_i)}{\sum_{j=1}^n P(B_j) P(A|B_j)}. \]
Permet d'« inverser » le conditionnement : connaître la cause ($B_i$) sachant l'effet ($A$).
\end{propriete}

\begin{aretenir}[Critère d'indépendance]
$A$ et $B$ indépendants $\iff P(A \cap B) = P(A)P(B)$.
\emph{Attention :} Incompatibles ($\cap = \varnothing$) $\not\Rightarrow$ Indépendants (sauf si $P(A)=0$ ou $P(B)=0$).
\end{aretenir}

\begin{exemplebox}[Test médical — Bayes]
Maladie $M$ : $P(M)=0,01$. Test positif $T$ : $P(T|M)=0,99$ (sensibilité), $P(T|\overline{M})=0,05$ (faux positifs).
$P(T) = 0,01\times0,99 + 0,99\times0,05 = 0,0594$.
$P(M|T) = \frac{0,01\times0,99}{0,0594} \approx 0,167$. Malgré un test performant, la rareté de la maladie fait que $P(M|T)$ reste faible.
\end{exemplebox}

\courssub{Variables aléatoires et épreuves de Bernoulli}

\begin{definition}[Variable aléatoire discrète]
Une \cle{variable aléatoire} (v.a.) discrète $X$ est une application de $\Omega$ dans $\mathbb{N}$ (ou $\mathbb{Z}$). Sa \cle{loi de probabilité} est le tableau des couples $(x_i, p_i)$ où $x_i$ sont les valeurs prises et $p_i = P(X=x_i)$, avec $\sum p_i = 1$.
\end{definition}

\begin{propriete}[Espérance, Variance, Écart-type]
\begin{align*}
    E(X) &= \sum_i x_i p_i \quad \text{(moyenne théorique)}. \\
    E(X^2) &= \sum_i x_i^2 p_i. \\
    V(X) &= E(X^2) - [E(X)]^2 = E[(X-E(X))^2] \quad \text{(dispersion)}. \\
    \sigma(X) &= \sqrt{V(X)} \quad \text{(écart-type, même unité que $X$)}.
\end{align*}
\textbf{Propriétés linéaires :} $E(aX+b) = aE(X)+b$ ; $V(aX+b) = a^2 V(X)$.
\end{propriete}

\begin{definition}[Épreuve de Bernoulli et Loi Binomiale]
Une \cle{épreuve de Bernoulli} est une expérience aléatoire à deux issues : \cle{Succès} ($S$, prob. $p$) et \cle{Échec} ($E$, prob. $q=1-p$).
Un \cle{schéma de Bernoulli} de paramètres $(n,p)$ est la répétition de $n$ épreuves de Bernoulli \emph{indépendantes} de même probabilité $p$.
La v.a. $X$ = « nombre de succès » suit la \cle{loi binomiale} $\mathcal{B}(n,p)$ :
\[ P(X=k) = \binom{n}{k} p^k (1-p)^{n-k}, \quad k=0,\dots,n. \]
\[ E(X) = np, \quad V(X) = np(1-p), \quad \sigma(X) = \sqrt{np(1-p)}. \]
\end{definition}

\begin{aretenir}[Calculs usuels avec la Binomiale]
\begin{itemize}
    \item « Au moins $k$ » : $P(X \ge k) = 1 - P(X \le k-1)$.
    \item « Au plus $k$ » : $P(X \le k) = \sum_{i=0}^k P(X=i)$.
    \item « Plus de $k$ » : $P(X > k) = 1 - P(X \le k)$.
    \item Utiliser la calculatrice (fonctions \texttt{binompdf} et \texttt{binomcdf}) ou les tables si $n$ petit.
\end{itemize}
\end{aretenir}

\courssub{Présentation générale des coniques et propriétés}

\begin{definition}[Conique — Définition par locus]
Une \cle{conique} est l'ensemble des points $M$ du plan tels que le rapport de la distance à un point fixe $F$ (foyer) sur la distance à une droite fixe $\Delta$ (directrice) est une constante $e > 0$ (excentricité) :
\[ \frac{MF}{d(M,\Delta)} = e. \]
\begin{itemize}
    \item $0 < e < 1$ : \cle{Ellipse} (deux foyers $F,F'$, deux directrices).
    \item $e = 1$ : \cle{Parabole} (un foyer, une directrice).
    \item $e > 1$ : \cle{Hyperbole} (deux foyers, deux directrices).
\end{itemize}
\end{definition}

\begin{propriete}[Équations réduites dans un repère orthonormé $(O;\vec{i},\vec{j})$]
\emph{On suppose les axes de symétries alignés sur les axes du repère (pas de terme $xy$).}
\begin{center}
\begin{tabularx}{\linewidth}{|l|X|X|X|}\hline
\textbf{Conique} & \textbf{Équation réduite} & \textbf{Paramètres ($>0$)} & \textbf{Éléments caractéristiques} \\ \hline
\textbf{Ellipse} & $\dfrac{x^2}{a^2} + \dfrac{y^2}{b^2} = 1$ & $a$ : demi-grand axe\par $b$ : demi-petit axe\par $c^2 = |a^2-b^2|$ & Sommets $(\pm a,0), (0,\pm b)$\par Foyers $(\pm c,0)$ si $a>b$\par Directrices $x = \pm a/e$\par $e = c/a < 1$ \\ \hline
\textbf{Parabole} & $y^2 = 2px$ (axe $Ox$)\par ou $x^2 = 2py$ (axe $Oy$) & $p \neq 0$ (paramètre) & Sommet $O(0,0)$\par Foyer $F(p/2, 0)$\par Directrice $x = -p/2$\par Axe focal : $Ox$ \\ \hline
\textbf{Hyperbole} & $\dfrac{x^2}{a^2} - \dfrac{y^2}{b^2} = 1$ (foyers sur $Ox$)\par ou $\dfrac{y^2}{a^2} - \dfrac{x^2}{b^2} = 1$ (foyers sur $Oy$) & $a,b > 0$\par $c^2 = a^2+b^2$ & Sommets $(\pm a,0)$\par Foyers $(\pm c,0)$\par Asymptotes $y = \pm \frac{b}{a}x$\par $e = c/a > 1$ \\ \hline
\end{tabularx}
\end{center}
\end{propriete}

\begin{popfigure}
\begin{tikzpicture}[scale=0.8, >=Stealth]
% Ellipse
\begin{scope}[xshift=-7cm]
\draw[popline] (-3.5,0) -- (3.5,0) node[right] {$x$};
\draw[popline] (0,-2.5) -- (0,2.5) node[above] {$y$};
\draw[popblue, thick] (0,0) ellipse (3 and 2);
\fill[popblue] (-2.236,0) circle (2pt) node[below] {$F$};
\fill[popblue] (2.236,0) circle (2pt) node[below] {$F'$};
\draw[popmuted, dashed] (-4.5,0) -- (-4.5,2) node[left] {$\Delta$};
\draw[popmuted, dashed] (4.5,0) -- (4.5,2);
\node[popdark] at (0,-2.8) {Ellipse $e<1$};
\end{scope}
% Parabole
\begin{scope}
\draw[popline] (-1,0) -- (5,0) node[right] {$x$};
\draw[popline] (0,-3) -- (0,3) node[above] {$y$};
\draw[poporange, thick, domain=-2.5:2.5, samples=100] plot (\x^2/2, \x);
\fill[poporange] (0.5,0) circle (2pt) node[below] {$F$};
\draw[popmuted, dashed] (-1, -3) -- (-1, 3) node[left] {$\Delta$};
\node[popdark] at (2,-3.3) {Parabole $e=1$};
\end{scope}
% Hyperbole
\begin{scope}[xshift=7cm]
\draw[popline] (-3.5,0) -- (3.5,0) node[right] {$x$};
\draw[popline] (0,-2.5) -- (0,2.5) node[above] {$y$};
\draw[poppurple, thick, domain=-2.5:2.5, samples=100] plot ({2/cos(\x r)}, {1.5*tan(\x r)});
\draw[poppurple, thick, domain=-2.5:2.5, samples=100] plot ({-2/cos(\x r)}, {1.5*tan(\x r)});
\draw[popmuted, thin, domain=-3:3] plot (\x, {0.75*\x});
\draw[popmuted, thin, domain=-3:3] plot (\x, {-0.75*\x});
\fill[poppurple] (-2.5,0) circle (2pt) node[below] {$F$};
\fill[poppurple] (2.5,0) circle (2pt) node[below] {$F'$};
\node[popdark] at (0,-2.8) {Hyperbole $e>1$};
\end{scope}
\end{tikzpicture}
\legende{Les trois coniques : ellipse, parabole, hyperbole (foyers $F$, directrices $\Delta$, asymptotes en pointillés).}
\end{popfigure}

\begin{propriete}[Propriétés optiques (exigibles au Probatoire)]
\begin{itemize}
    \item \textbf{Ellipse} : Un rayon issu d'un foyer $F$ se réfléchit sur l'ellipse vers l'autre foyer $F'$. (La tangente bissecte l'angle $\widehat{FMF'}$).
    \item \textbf{Parabole} : Un rayon parallèle à l'axe se réfléchit vers le foyer $F$ (et réciproquement). (La tangente bissecte l'angle entre la parallèle à l'axe et $(MF)$).
    \item \textbf{Hyperbole} : Un rayon issu d'un foyer $F$ se réfléchit comme s'il venait de l'autre foyer $F'$. (La tangente bissecte l'angle $\widehat{FMF'}$).
\end{itemize}
\end{propriete}

\courssub{Tangente en un point à une conique}

\begin{propriete}[Équation de la tangente — Méthode du remplacement]
Soit une conique d'équation réduite $F(x,y)=0$ et $M_0(x_0,y_0)$ un point de cette conique. L'équation de la tangente en $M_0$ s'obtient par les remplacements systématiques :
\[ x^2 \to x_0 x,\quad y^2 \to y_0 y,\quad x \to \frac{x+x_0}{2},\quad y \to \frac{y+y_0}{2},\quad xy \to \frac{x_0 y + y_0 x}{2} \text{ (si terme } xy \text{ présent, hors programme réduit)}. \]
\emph{Formules directes à connaître :}
\begin{itemize}
    \item \textbf{Ellipse} $\frac{x^2}{a^2}+\frac{y^2}{b^2}=1$ : $\frac{x_0 x}{a^2} + \frac{y_0 y}{b^2} = 1$.
    \item \textbf{Parabole} $y^2 = 2px$ : $y_0 y = p(x + x_0)$.
    \item \textbf{Hyperbole} $\frac{x^2}{a^2}-\frac{y^2}{b^2}=1$ : $\frac{x_0 x}{a^2} - \frac{y_0 y}{b^2} = 1$.
\end{itemize}
\end{propriete}

\begin{methode}[Déterminer la tangente en un point]
\begin{enumerate}
    \item Vérifier que $M_0(x_0,y_0)$ appartient à la conique (remplacer dans l'équation).
    \item Appliquer la formule de remplacement adaptée au type de conique.
    \item Simplifier pour obtenir la forme cartésienne $ax+by+c=0$ ou réduite $y=mx+p$.
    \item (Optionnel) Vérifier la propriété optique : angle avec $(M_0F)$ = angle avec $(M_0F')$ ou l'axe.
\end{enumerate}
\end{methode}

\courssub{Construction d'une conique et applications intégrées}

\begin{methode}[Construction de l'ellipse — Méthode du jardinier]
Données : centre $O$, axes $2a, 2b$ (ou foyers $F,F'$ et $2a$).
\begin{enumerate}
    \item Calculer $c = \sqrt{|a^2-b^2|}$. Placer les foyers $F(-c,0), F'(c,0)$.
    \item Planter deux piquets en $F, F'$. Nœud une corde de longueur $2a$ autour des piquets.
    \item Tendre la corde avec un crayon et faire tourner : le crayon trace l'ellipse (propriété $MF+MF'=2a$).
\end{enumerate}
\end{methode}

\begin{methode}[Construction de la parabole — Règle et équerre]
Données : foyer $F$, directrice $\Delta$.
\begin{enumerate}
    \item Placer une règle le long de $\Delta$. Une équerre glisse sur la règle.
    \item Un fil tendu de $F$ au sommet de l'équerre (côté mobile), longueur égale à la distance du crayon à $\Delta$.
    \item En faisant glisser l'équerre le long de la règle, le crayon trace la parabole (propriété $MF = d(M,\Delta)$).
\end{enumerate}
\end{methode}

\begin{exemplebox}[Application intégrée — Voûte en arc parabolique]
Un pont en arc parabolique : portée $L=40$ m, flèche $f=10$ m. Repère : origine au milieu de la portée, $Ox$ horizontal, $Oy$ vertical vers le haut. Sommet $S(0,10)$, appuis $A(-20,0), B(20,0)$.
Équation : $y = -\frac{x^2}{40} + 10$ (ou $x^2 = -40(y-10)$).
Poussée horizontale aux appuis (charge uniforme $q$) : $H = \frac{q L^2}{8f}$.
\end{exemplebox}

\courssub{Méthodes et exercices résolus}

\begin{methode}[Analyser une expérience aléatoire : vocabulaire et modélisation]
    \begin{enumerate}
        \item \textbf{Identifier l'expérience} : vérifier qu'elle admet plusieurs résultats possibles, qu'on ne peut pas prédire avec certitude lequel se produira, et qu'elle est reproductible à l'identique.
        \item \textbf{Définir l'univers $\Omega$} : l'ensemble de \emph{tous} les résultats élémentaires possibles. Le noter explicitement (liste, ensemble cartésien, arbre).
        \item \textbf{Identifier les événements étudiés} : ce sont des parties de $\Omega$. Utiliser la notation ensembliste ($A \subset \Omega$, $A \cup B$, $A \cap B$, $\overline{A}$).
        \item \textbf{Préciser le modèle de probabilité} : équiprobabilité (hasard « symétrique ») ou loi donnée (fréquences, tableau). Vérifier $\sum p_i = 1$.
        \item \textbf{Traduire la question} en langage ensembliste avant de calculer : « au moins un » $\to$ complémentaire de « aucun » ; « l'un ou l'autre mais pas les deux » $\to$ différence symétrique.
    \end{enumerate}
\end{methode}

\begin{exoresolu}[Usinage et contrôle qualité — Univers et événements]
Un tourneur effectue une série de 3 pièces sur une machine à commande numérique. Pour chaque pièce, l'événement « pièce conforme » est noté $C$, l'événement « pièce non conforme » est noté $N$. On suppose que les résultats des 3 pièces sont indépendants et que la probabilité d'avoir une pièce conforme est $0,95$.
\begin{enumerate}
    \item Décrire l'univers $\Omega$ de l'expérience « contrôle des 3 pièces » et donner son cardinal.
    \item Écrire en lettres et en ensembliste les événements :
    \begin{itemize}
        \item $A$ : « Les 3 pièces sont conformes »,
        \item $B$ : « Au moins une pièce est non conforme »,
        \item $C_1$ : « La première pièce est conforme ».
    \end{itemize}
    \item Calculer $P(A)$, $P(B)$ et $P(C_1)$.
    \item Vérifier que $A$ et $B$ sont contraires.
\end{enumerate}
\tcblower
\textbf{Solution détaillée}
\begin{enumerate}
    \item L'expérience consiste à observer la conformité de 3 pièces successives. Chaque pièce a 2 issues ($C$ ou $N$). L'univers est l'ensemble des triplets :
    \[ \Omega = \{ (C,C,C),\ (C,C,N),\ (C,N,C),\ (N,C,C),\ (C,N,N),\ (N,C,N),\ (N,N,C),\ (N,N,N) \}. \]
    $\text{Card}(\Omega) = 2^3 = 8$.
    \item 
    \begin{itemize}
        \item $A$ : « Les 3 pièces sont conformes » $\to A = \{ (C,C,C) \}$.
        \item $B$ : « Au moins une pièce est non conforme » $\to$ C'est le complémentaire de $A$ dans $\Omega$.
        $B = \overline{A} = \Omega \setminus \{ (C,C,C) \} = \{ (C,C,N),\ (C,N,C),\ (N,C,C),\ (C,N,N),\ (N,C,N),\ (N,N,C),\ (N,N,N) \}$.
        \item $C_1$ : « La première pièce est conforme » $\to C_1 = \{ (C,C,C),\ (C,C,N),\ (C,N,C),\ (C,N,N) \}$.
    \end{itemize}
    \item On utilise l'indépendance et $P(C)=0,95$, $P(N)=0,05$.
    \begin{align*}
        P(A) &= P(C \cap C \cap C) = 0,95 \times 0,95 \times 0,95 = 0,95^3 = 0,857375. \\
        P(B) &= P(\overline{A}) = 1 - P(A) = 1 - 0,857375 = 0,142625. \\
        P(C_1) &= P(\text{1ère conforme}) = 0,95 \quad \text{(indépendance, les 2 autres n'importent pas)}.
    \end{align*}
    \item $A \cap B = \varnothing$ (on ne peut pas avoir toutes conformes ET au moins une non conforme) et $A \cup B = \Omega$ (soit toutes conformes, soit au moins une non conforme). Donc $B = \overline{A}$, ils sont bien événements contraires.
\end{enumerate}
\textbf{Conclusion :} $P(A) \approx 0,857$, $P(B) \approx 0,143$, $P(C_1)=0,95$. La modélisation par triplet ordonné est indispensable car l'ordre des pièces compte pour l'indépendance.
\end{exoresolu}

\begin{methode}[Calculer des probabilités conditionnelles et utiliser la formule de Bayes]
    \begin{enumerate}
        \item \textbf{Repérer le conditionnement} : phrase « sachant que », « si on sait que », « parmi les... ». Noter $P(B|A)$ : probabilité de $B$ \emph{sachant} $A$ réalisé.
        \item \textbf{Rappeler la définition} : $P(B|A) = \dfrac{P(A \cap B)}{P(A)}$ (si $P(A) > 0$). En déduire \textbf{la formule des probabilités composées} : $P(A \cap B) = P(A) \times P(B|A)$.
        \item \textbf{Construire un arbre pondéré} (ou un tableau de contingence) : branches = événements, poids = probabilités. Vérifier que la somme des probabilités des issues d'un même nœud vaut 1.
        \item \textbf{Appliquer la loi des probabilités totales} si on cherche $P(B)$ et qu'on a une partition $(A_i)$ : $P(B) = \sum P(A_i)P(B|A_i)$.
        \item \textbf{Appliquer Bayes} pour « inverser » le conditionnement : $P(A_i|B) = \dfrac{P(A_i)P(B|A_i)}{\sum_j P(A_j)P(B|A_j)}$.
        \item \textbf{Conclure par une phrase} interprétant le nombre trouvé dans le contexte.
    \end{enumerate}
\end{methode}

\begin{exoresolu}[Maintenance prédictive — Probabilités conditionnelles et Bayes]
Une usine dispose de deux lignes de production $L_1$ et $L_2$. $L_1$ fournit $60\%$ de la production, $L_2$ les $40\%$ restants.
Le taux de pièces défectueuses est de $2\%$ sur $L_1$ et de $5\%$ sur $L_2$.
On prélève une pièce au hasard dans la production totale.
\begin{enumerate}
    \item Construire l'arbre pondéré de la situation.
    \item Calculer la probabilité que la pièce soit défectueuse.
    \item Sachant que la pièce est défectueuse, quelle est la probabilité qu'elle vienne de la ligne $L_1$ ?
    \item L'usine décide de ne contrôler que les pièces de $L_1$. Quelle proportion de pièces défectueuses échappent au contrôle ?
\end{enumerate}
\tcblower
\textbf{Solution détaillée}
\begin{enumerate}
    \item \textbf{Arbre pondéré} :
    \begin{itemize}
        \item Niveau 1 (Choix ligne) : $L_1$ ($p=0,6$), $L_2$ ($p=0,4$).
        \item Niveau 2 (État pièce | Ligne) :
        \begin{itemize}
            \item Depuis $L_1$ : $D$ (défectueuse) $p=0,02$ ; $\overline{D}$ (conforme) $p=0,98$.
            \item Depuis $L_2$ : $D$ $p=0,05$ ; $\overline{D}$ $p=0,95$.
        \end{itemize}
    \end{itemize}
    \item \textbf{Loi des probabilités totales} pour $P(D)$ :
    \[ P(D) = P(L_1)P(D|L_1) + P(L_2)P(D|L_2) = 0,6 \times 0,02 + 0,4 \times 0,05 = 0,012 + 0,020 = 0,032. \]
    \item \textbf{Formule de Bayes} pour $P(L_1|D)$ :
    \[ P(L_1|D) = \frac{P(L_1 \cap D)}{P(D)} = \frac{P(L_1)P(D|L_1)}{P(D)} = \frac{0,012}{0,032} = \frac{12}{32} = 0,375. \]
    \item Les pièces défectueuses échappant au contrôle sont celles de $L_2$ défectueuses.
    Proportion parmi la production totale : $P(L_2 \cap D) = 0,4 \times 0,05 = 0,02$.
    Proportion \emph{parmi les pièces défectueuses} : $P(L_2|D) = 1 - P(L_1|D) = 1 - 0,375 = 0,625$.
    \textbf{Conclusion :} $62,5\%$ des pièces défectueuses échappent au contrôle si on ne contrôle que $L_1$.
\end{enumerate}
\end{exoresolu}

\begin{methode}[Étudier une variable aléatoire discrète : loi, espérance, variance — Épreuves de Bernoulli]
    \begin{enumerate}
        \item \textbf{Identifier la variable aléatoire $X$} : quelle grandeur numérique ? Quelles valeurs $x_i$ prend-elle ? (Entiers le plus souvent).
        \item \textbf{Établir le tableau de loi} : lignes $x_i$ et $P(X=x_i)$. Vérifier $\sum P(X=x_i) = 1$.
        \item \textbf{Calculer l'espérance} $E(X) = \sum x_i P(X=x_i)$. Interprétation : moyenne théorique sur un grand nombre de répétitions.
        \item \textbf{Calculer la variance} $V(X) = E(X^2) - [E(X)]^2$ avec $E(X^2) = \sum x_i^2 P(X=x_i)$. Écart-type $\sigma(X) = \sqrt{V(X)}$.
        \item \textbf{Repérer un schéma de Bernoulli} : $n$ épreuves indépendantes, même probabilité de succès $p$, variable $X$ = « nombre de succès ». Alors $X \sim \mathcal{B}(n,p)$.
        \item \textbf{Formules Binomiale} : $P(X=k) = \binom{n}{k} p^k (1-p)^{n-k}$ ; $E(X)=np$ ; $V(X)=np(1-p)$.
        \item \textbf{Calculer « au moins / au plus »} : utiliser le complémentaire ou la somme. Pour $n$ grand, utiliser la calculatrice (fonction \texttt{binomcdf}) ou les tables.
    \end{enumerate}
\end{methode}

\begin{exoresolu}[Réseau électrique — Variable aléatoire binomiale]
Un câble souterrain comporte 20 jonctions. La probabilité qu'une jonction soit défectueuse est de $0,03$, indépendamment des autres. On note $X$ le nombre de jonctions défectueuses sur le câble.
\begin{enumerate}
    \item Justifier que $X$ suit une loi binomiale et donner ses paramètres.
    \item Calculer $P(X=0)$ et $P(X=1)$ (donner une valeur approchée à $10^{-4}$).
    \item Calculer la probabilité qu'il y ait au plus 2 jonctions défectueuses.
    \item Déterminer l'espérance et l'écart-type de $X$. Interpréter l'espérance.
\end{enumerate}
\tcblower
\textbf{Solution détaillée}
\begin{enumerate}
    \item On répète $n=20$ fois l'épreuve « tester une jonction » (indépendantes, même $p=0,03$). $X$ compte le nombre de succès (jonction défectueuse). Donc $X \sim \mathcal{B}(n=20,\ p=0,03)$.
    \item 
    \begin{align*}
        P(X=0) &= \binom{20}{0} 0,03^0 \times 0,97^{20} = 1 \times 1 \times 0,97^{20} \approx 0,5438. \\
        P(X=1) &= \binom{20}{1} 0,03^1 \times 0,97^{19} = 20 \times 0,03 \times 0,97^{19} \approx 0,3364.
    \end{align*}
    \item « Au plus 2 » $\iff X \in \{0, 1, 2\}$.
    \begin{align*}
        P(X=2) &= \binom{20}{2} 0,03^2 \times 0,97^{18} = 190 \times 0,0009 \times 0,97^{18} \approx 0,0988. \\
        P(X \le 2) &= P(X=0)+P(X=1)+P(X=2) \approx 0,5438 + 0,3364 + 0,0988 = 0,9790.
    \end{align*}
    \item 
    \begin{align*}
        E(X) &= n \times p = 20 \times 0,03 = 0,6. \\
        V(X) &= n \times p \times (1-p) = 20 \times 0,03 \times 0,97 = 0,582. \\
        \sigma(X) &= \sqrt{0,582} \approx 0,763.
    \end{align*}
    \textbf{Interprétation :} Sur un grand nombre de câbles de 20 jonctions, le nombre moyen de jonctions défectueuses est de $0,6$ (soit moins d'une par câble en moyenne).
\end{enumerate}
\end{exoresolu}

\begin{methode}[Identifier une conique, écrire son équation réduite et en déduire ses éléments caractéristiques]
    \begin{enumerate}
        \item \textbf{Repérer le type} d'équation cartésienne (degré 2 en $x,y$ sans terme $xy$ dans le programme réduit) :
        \begin{itemize}
            \item Ellipse : $\dfrac{x^2}{a^2} + \dfrac{y^2}{b^2} = 1$ ($a>b>0$ ou $b>a>0$).
            \item Parabole : $y^2 = 2px$ (axe $Ox$) ou $x^2 = 2py$ (axe $Oy$) ($p \neq 0$).
            \item Hyperbole : $\dfrac{x^2}{a^2} - \dfrac{y^2}{b^2} = 1$ (foyers sur $Ox$) ou $\dfrac{y^2}{a^2} - \dfrac{x^2}{b^2} = 1$ (foyers sur $Oy$).
        \end{itemize}
        \item \textbf{Extraire les paramètres} $a, b, p$ (toujours positifs dans l'équation réduite).
        \item \textbf{Déduire les éléments géométriques} :
        \begin{itemize}
            \item \textbf{Ellipse} : Sommets $(\pm a, 0)$ et $(0, \pm b)$ ; Foyers $(\pm c, 0)$ avec $c^2 = |a^2-b^2|$ ; Excentricité $e = c/a < 1$ ; Directrices $x = \pm a/e$.
            \item \textbf{Parabole} $y^2=2px$ : Sommet $O(0,0)$ ; Foyer $F(p/2, 0)$ ; Directrice $x = -p/2$ ; Axe $Ox$.
            \item \textbf{Hyperbole} $\frac{x^2}{a^2}-\frac{y^2}{b^2}=1$ : Sommets $(\pm a, 0)$ ; Foyers $(\pm c, 0)$ avec $c^2 = a^2+b^2$ ; Asymptotes $y = \pm \frac{b}{a}x$ ; Excentricité $e = c/a > 1$.
        \end{itemize}
        \item \textbf{Tracer un croquis} soigné : repère, asymptotes (hyperbole), foyers, directrices, sommet(s).
    \end{enumerate}
\end{methode}

\begin{exoresolu}[Antenne parabolique — Équation réduite et éléments]
Une antenne parabolique a un diamètre de $4$ m et une profondeur de $0,5$ m. On place son sommet à l'origine $O$, son axe sur $Ox$ (vers l'intérieur de l'antenne) et son plan d'ouverture perpendiculaire à $Ox$.
\begin{enumerate}
    \item Déterminer l'équation réduite de la parabole génératrice dans ce repère.
    \item Calculer les coordonnées du foyer $F$. Quelle est la distance du foyer au sommet ? (C'est la distance focale).
    \item Écrire l'équation de la directrice.
    \item Le récepteur doit être placé au foyer. Vérifier qu'il se trouve bien à l'intérieur de l'antenne (profondeur $0,5$ m).
\end{enumerate}
\tcblower
\textbf{Solution détaillée}
\begin{enumerate}
    \item L'antenne est un paraboloïde de révolution. Dans le plan $(Oxy)$, la parabole génératrice a pour équation $y^2 = 2px$ ($p>0$ car axe $Ox$ orienté vers l'intérieur).
    Le bord de l'antenne correspond à $x = 0,5$ (profondeur) et $y = \pm 2$ (rayon = diamètre/2).
    Le point $M(0,5 ; 2)$ appartient à la parabole :
    \[ 2^2 = 2p \times 0,5 \iff 4 = p \iff p = 4. \]
    L'équation réduite est \textbf{$y^2 = 8x$}.
    \item Pour $y^2 = 2px$, le foyer est $F(\frac{p}{2}, 0)$. Ici $F(2, 0)$.
    La distance focale (sommet au foyer) est $\frac{p}{2} = 2$ m.
    \item La directrice a pour équation $x = -\frac{p}{2}$, soit \textbf{$x = -2$}.
    \item Le foyer $F$ a pour abscisse $x_F = 2$ m. La profondeur de l'antenne est $0,5$ m (plan $x=0,5$).
    Comme $2 > 0,5$, le foyer se trouve \textbf{à l'extérieur} de l'antenne (derrière le réflecteur).
    \textbf{Attention :} Cela signifie qu'avec ces dimensions ($D=4$m, prof=$0,5$m), le foyer est trop loin. Pour un foyer à l'intérieur, il faudrait $p/2 < 0,5 \iff p < 1$. Ici $p=4$. L'antenne est trop « plate » pour ce diamètre.
\end{enumerate}
\textbf{Conclusion :} Équation $y^2=8x$, Foyer $F(2,0)$, Directrice $x=-2$. Le foyer est à 2 m du sommet, hors de l'antenne (profondeur 0,5 m).
\end{exoresolu}

\begin{methode}[Déterminer l'équation de la tangente en un point à une conique]
    \begin{enumerate}
        \item \textbf{Vérifier que le point $M_0(x_0,y_0)$ appartient à la conique} (remplacer dans l'équation).
        \item \textbf{Appliquer la formule de la tangente} (méthode du « remplacement ») selon le type :
        \begin{itemize}
            \item Ellipse $\frac{x^2}{a^2}+\frac{y^2}{b^2}=1$ : $\frac{x_0 x}{a^2} + \frac{y_0 y}{b^2} = 1$.
            \item Parabole $y^2 = 2px$ : $y_0 y = p(x + x_0)$.
            \item Hyperbole $\frac{x^2}{a^2}-\frac{y^2}{b^2}=1$ : $\frac{x_0 x}{a^2} - \frac{y_0 y}{b^2} = 1$.
        \end{itemize}
        \item \textbf{Mettre sous forme réduite} $y = ax + b$ ou $ax+by+c=0$ si demandé.
        \item \textbf{Propriété optique (si demandée)} : La tangente fait des angles égaux avec les droites $(MF)$ et $(MF')$ (ellipse/hyperbole) ou avec l'axe et la parallèle à l'axe passant par $M$ (parabole).
    \end{enumerate}
\end{methode}

\begin{exoresolu}[Trajectoire satellite — Tangente à l'ellipse]
Un satellite décrit une orbite elliptique d'équation $\dfrac{x^2}{25} + \dfrac{y^2}{9} = 1$ (distances en milliers de km). Au point $M_0(4 ; \frac{9}{5})$, on veut déterminer la direction de la poussée tangentielle pour corriger l'orbite.
\begin{enumerate}
    \item Vérifier que $M_0$ appartient à l'ellipse.
    \item Déterminer l'équation cartésienne de la tangente à l'ellipse en $M_0$.
    \item Calculer l'angle (en degrés, arrondi au degré) que cette tangente fait avec l'axe $Ox$.
\end{enumerate}
\tcblower
\textbf{Solution détaillée}
\begin{enumerate}
    \item On vérifie : $\frac{4^2}{25} + \frac{(9/5)^2}{9} = \frac{16}{25} + \frac{81/25}{9} = \frac{16}{25} + \frac{9}{25} = \frac{25}{25} = 1$. $M_0$ appartient bien à l'ellipse.
    \item Ellipse : $a^2=25$, $b^2=9$. Formule de la tangente en $M_0(x_0,y_0)$ :
    \[ \frac{x_0 x}{a^2} + \frac{y_0 y}{b^2} = 1 \iff \frac{4x}{25} + \frac{\frac{9}{5}y}{9} = 1 \iff \frac{4x}{25} + \frac{y}{5} = 1. \]
    On multiplie par 25 : $4x + 5y = 25 \iff 5y = -4x + 25 \iff \textbf{y = -0,8x + 5}$.
    \item Le coefficient directeur est $m = -0,8 = -\frac{4}{5}$.
    L'angle $\alpha$ avec $Ox$ (orienté positivement) tel que $\tan \alpha = m$.
    $\alpha = \arctan(-0,8) \approx -38,66^\circ$.
    L'angle aigu avec $Ox$ est $39^\circ$ (arrondi). La tangente est descendante.
\end{enumerate}
\textbf{Conclusion :} Tangente : $4x+5y=25$. Angle avec $Ox \approx 39^\circ$ (pente négative).
\end{exoresolu}

\begin{methode}[Construire une conique par la méthode du jardinier (ellipse) ou à la règle/équerre (parabole) — Application intégrée]
    \begin{enumerate}
        \item \textbf{Ellipse (Jardinier)} : Données : centre $O$, axes $2a, 2b$ ou foyers $F,F'$ et $2a$.
        \begin{itemize}
            \item Calculer $c = \sqrt{|a^2-b^2|}$. Placer $F(-c,0), F'(c,0)$.
            \item Planter deux piquets en $F, F'$. Nœud une corde de longueur $2a$ autour des piquets.
            \item Tendre la corde avec un crayon et tourner : le crayon trace l'ellipse ($MF+MF'=2a$).
        \end{itemize}
        \item \textbf{Parabole (Règle/Équerre)} : Données : foyer $F$, directrice $\Delta$.
        \begin{itemize}
            \item Placer une règle le long de $\Delta$. Une équerre glisse sur la règle.
            \item Un fil tendu de $F$ au sommet de l'équerre (côté mobile), longueur égale à la distance du crayon à $\Delta$.
            \item En faisant glisser l'équerre, le crayon trace la parabole ($MF = d(M,\Delta)$).
        \end{itemize}
        \item \textbf{Application intégrée (Pont/Voûte)} : Modéliser l'arc par une parabole (charges uniformes horizontales) ou une chaînette (poids propre, hors programme, mais souvent approximée par parabole). Utiliser l'équation réduite pour trouver la flèche, la poussée, la hauteur à un point.
    \end{enumerate}
\end{methode}

\begin{exoresolu}[Pont en arc parabolique — Construction et application]
On veut construire un pont en arc parabolique. La portée (distance entre les deux appuis) est de $40$ m. La flèche (hauteur max au milieu) est de $10$ m. On place l'origine au milieu de la portée, au niveau des appuis, $Ox$ horizontal, $Oy$ vertical vers le haut.
\begin{enumerate}
    \item Déterminer l'équation de la parabole dans ce repère.
    \item Calculer la hauteur de l'arc à $5$ m d'un appui.
    \item Pour tracer l'arc au sol (méthode du jardinier adaptée ou quadrillage), on repère des points. Donner les coordonnées des points pour $x = 0, \pm 5, \pm 10, \pm 15, \pm 20$.
    \item L'arc supporte une charge uniformément répartie horizontale de $50$ kN/m. La poussée horizontale $H$ aux appuis vaut $H = \frac{q L^2}{8f}$ ($q$ charge, $L$ portée, $f$ flèche). Calculer $H$.
\end{enumerate}
\tcblower
\textbf{Solution détaillée}
\begin{enumerate}
    \item Sommet au milieu, hauteur $10$ m $\to$ Sommet $S(0, 10)$. Appuis : $A(-20, 0)$ et $B(20, 0)$.
    Parabole à axe vertical, sommet $S(0,10)$, concavité vers le bas $\to$ équation $x^2 = -2p(y-10)$ ou $y = -\frac{1}{2p}x^2 + 10$.
    Passage par $B(20,0)$ : $0 = -\frac{1}{2p}(400) + 10 \iff \frac{400}{2p} = 10 \iff 2p = 40 \iff p = 20$.
    Équation : \textbf{$x^2 = -40(y-10)$} ou \textbf{$y = 10 - \frac{x^2}{40}$}.
    \item À $5$ m d'un appui (disons $B(20,0)$), l'abscisse est $x = 15$ (ou $-15$).
    $y = 10 - \frac{15^2}{40} = 10 - \frac{225}{40} = 10 - 5,625 = \textbf{4,375 m}$.
    \item Tableau de points :
    \begin{center}
    \begin{tabular}{|c|c|c|c|c|c|c|c|c|c|}\hline
    $x$ & $-20$ & $-15$ & $-10$ & $-5$ & $0$ & $5$ & $10$ & $15$ & $20$ \\ \hline
    $y$ & $0$ & $4,375$ & $7,5$ & $9,375$ & $10$ & $9,375$ & $7,5$ & $4,375$ & $0$ \\ \hline
    \end{tabular}
    \end{center}
    \item Poussée horizontale : $q = 50$ kN/m, $L = 40$ m, $f = 10$ m.
    \[ H = \frac{50 \times 40^2}{8 \times 10} = \frac{50 \times 1600}{80} = 50 \times 20 = \textbf{1000 kN} = \textbf{1 MN}. \]
\end{enumerate}
\textbf{Conclusion :} Équation $y=10-x^2/40$. Hauteur à 5m de l'appui : 4,375 m. Poussée horizontale : 1000 kN.
\end{exoresolu}

\begin{attention}[Les 5 erreurs qui coûtent des points au Probatoire]
    \begin{enumerate}
        \item \textbf{Confondre $P(A|B)$ et $P(B|A)$} : Ne jamais écrire $P(A|B) = P(A \cap B)/P(B)$ si on cherche $P(B|A)$. Toujours écrire la définition : $P(\text{cherché}|\text{donné}) = \frac{P(\text{cherché} \cap \text{donné})}{P(\text{donné})}$. Perte de 2 à 3 points.
        \item \textbf{Oublier de vérifier l'appartenance du point à la conique} avant d'écrire la tangente. Si $M_0$ n'est pas sur la courbe, la formule du remplacement donne une droite quelconque, pas la tangente. Point de méthode perdu.
        \item \textbf{Mauvaise identification des paramètres $a, b, p$ dans l'équation réduite} : Écrire $y^2 = 2px$ avec $p$ négatif sans changer l'orientation de l'axe, ou confondre $a$ (demi-grand axe) et $c$ (demi-distance focale) pour l'ellipse/hyperbole. Entraîne des erreurs en cascade sur foyers, directrices, asymptotes.
        \item \textbf{Calculer $P(X \ge k)$ au lieu de $P(X > k)$ (ou inversement) pour la binomiale} : « Au moins $k$ » = $P(X \ge k) = 1 - P(X \le k-1)$. « Plus de $k$ » = $P(X > k) = 1 - P(X \le k)$. Lire attentivement l'énoncé.
        \item \textbf{Notation décimale vs fractionnaire / Arrondi} : Donner une probabilité sous forme de fraction non simplifiée ($\frac{12}{32}$ au lieu de $\frac{3}{8}$ ou $0,375$) ou arrondir trop tôt dans les calculs intermédiaires (garder les valeurs exactes ou 5 décimales minimum). Perte de point de présentation/précision.
    \end{enumerate}
\end{attention}

\courssub{Exercices d'entraînement}

\begin{exercice}[Probabilités conditionnelles — Contrôle qualité]
Une usine fabrique des boulons sur trois machines $M_1, M_2, M_3$ qui produisent respectivement $50\%, 30\%, 20\%$ de la production. Les taux de défaut sont $1\%, 2\%, 3\%$.
\begin{enumerate}
    \item Construire l'arbre pondéré.
    \item Calculer la probabilité qu'un boulon tiré au hasard soit défectueux.
    \item Sachant qu'un boulon est défectueux, quelle est la probabilité qu'il vienne de $M_3$ ?
\end{enumerate}
\end{exercice}

\begin{corrige}[Exercice 1]
\begin{enumerate}
    \item Arbre : Niveau 1 $M_i$ ($0,5; 0,3; 0,2$). Niveau 2 $D|\overline{D}$ avec prob conditionnelles.
    \item $P(D) = 0,5\times0,01 + 0,3\times0,02 + 0,2\times0,03 = 0,005+0,006+0,006 = 0,017$.
    \item $P(M_3|D) = \frac{0,2\times0,03}{0,017} = \frac{0,006}{0,017} \approx 0,353$.
\end{enumerate}
\end{corrige}

\begin{exercice}[Variable aléatoire — Gestion de stock]
Un magasin reçoit une commande de 10 articles. La probabilité qu'un article soit abîmé à la livraison est de $0,05$. Soit $X$ le nombre d'articles abîmés.
\begin{enumerate}
    \item Quelle loi suit $X$ ? Préciser les paramètres.
    \item Calculer $P(X=0)$ et $P(X \ge 2)$.
    \item Calculer $E(X)$ et $\sigma(X)$.
\end{enumerate}
\end{exercice}

\begin{corrige}[Exercice 2]
\begin{enumerate}
    \item $X \sim \mathcal{B}(n=10, p=0,05)$.
    \item $P(X=0) = 0,95^{10} \approx 0,5987$. $P(X \ge 2) = 1 - P(X=0) - P(X=1) = 1 - 0,5987 - 10\times0,05\times0,95^9 \approx 1 - 0,5987 - 0,3151 = 0,0862$.
    \item $E(X) = 10\times0,05 = 0,5$. $V(X) = 10\times0,05\times0,95 = 0,475$. $\sigma(X) \approx 0,689$.
\end{enumerate}
\end{corrige}

\begin{exercice}[Conique — Ellipse et tangente]
On considère l'ellipse $\mathcal{E}$ d'équation $\frac{x^2}{16} + \frac{y^2}{7} = 1$.
\begin{enumerate}
    \item Déterminer les coordonnées des foyers $F, F'$, l'excentricité $e$ et les équations des directrices.
    \item Déterminer l'équation de la tangente $\mathcal{T}$ en $M_0(2; \frac{\sqrt{21}}{2})$.
    \item Vérifier que $\mathcal{T}$ bissecte l'angle $\widehat{F M_0 F'}$.
\end{enumerate}
\end{exercice}

\begin{corrige}[Exercice 3]
\begin{enumerate}
    \item $a^2=16, b^2=7 \Rightarrow a=4, b=\sqrt{7}$. $c^2 = a^2-b^2 = 9 \Rightarrow c=3$. Foyers $F(-3,0), F'(3,0)$. $e = c/a = 3/4 = 0,75$. Directrices $x = \pm a/e = \pm 16/3$.
    \item Tangente : $\frac{2x}{16} + \frac{\frac{\sqrt{21}}{2}y}{7} = 1 \iff \frac{x}{8} + \frac{y}{2\sqrt{7}} = 1 \iff \sqrt{7}x + 4y = 8\sqrt{7}$.
    \item Vecteurs $\overrightarrow{M_0F}=(-5, -\frac{\sqrt{21}}{2})$, $\overrightarrow{M_0F'}=(1, -\frac{\sqrt{21}}{2})$. Vecteur directeur de $\mathcal{T}$ : $\vec{u}=(4, -\sqrt{7})$. Calculer les cosinus des angles $(\vec{u}, \overrightarrow{M_0F})$ et $(\vec{u}, \overrightarrow{M_0F'})$ : ils sont égaux.
\end{enumerate}
\end{corrige}

\begin{exercice}[Construction — Parabole]
On veut tracer une parabole de foyer $F(0,2)$ et de directrice $\Delta : y = -2$.
\begin{enumerate}
    \item Écrire l'équation réduite de cette parabole.
    \item Décrire la construction à la règle et à l'équerre.
    \item Déterminer les coordonnées du point de la parabole d'abscisse $x=4$.
\end{enumerate}
\end{exercice}

\begin{corrige}[Exercice 4]
\begin{enumerate}
    \item Sommet à l'origine (milieu de $F$ et $\Delta$), $p/2 = 2 \Rightarrow p=4$. Axe $Oy$ vers le haut. Équation : $x^2 = 8y$ ou $y = x^2/8$.
    \item Règle sur $\Delta$ ($y=-2$). Équerre glissant sur la règle. Fil de $F$ au sommet de l'équerre, crayon au bout du fil tendu.
    \item Pour $x=4$, $y = 16/8 = 2$. Point $P(4,2)$. Vérification : $PF = \sqrt{4^2+0^2}=4$, $d(P,\Delta)=2-(-2)=4$.
\end{enumerate}
\end{corrige}

\newpage

\coursec{Exercices}

\courssub{Je vérifie mes connaissances}

\begin{exercice}[QCM : vocabulaire des probabilités]
Pour chaque question, une seule réponse est exacte. Recopie la lettre correspondante.
\begin{enumerate}
\item On lance un dé équilibré à six faces. L'univers $\Omega$ est :
\begin{enumerate}
\item[(a)] $\{1;2;3;4;5\}$ \quad (b) $\{1;2;3;4;5;6\}$ \quad (c) $\{0;1;2;3;4;5;6\}$
\end{enumerate}
\item Si $A$ et $B$ sont deux événements incompatibles, alors :
\begin{enumerate}
\item[(a)] $P(A\cup B)=P(A)+P(B)$ \quad (b) $P(A\cup B)=P(A)\times P(B)$ \quad (c) $P(A\cap B)=P(A)\times P(B)$
\end{enumerate}
\item La probabilité de l'événement contraire de $A$ est :
\begin{enumerate}
\item[(a)] $1+P(A)$ \quad (b) $1-P(A)$ \quad (c) $P(A)-1$
\end{enumerate}
\item Une variable aléatoire $X$ suit la loi binomiale $\mathcal{B}(n,p)$. Son espérance est :
\begin{enumerate}
\item[(a)] $np$ \quad (b) $np(1-p)$ \quad (c) $n(1-p)$
\end{enumerate}
\end{enumerate}
\end{exercice}

\begin{exercice}[Vrai ou faux, en justifiant]
Réponds par vrai ou faux en justifiant chaque réponse.
\begin{enumerate}
\item Une probabilité peut être égale à $1,2$.
\item Si $P(A)=0,4$, alors $P(\overline{A})=0,6$.
\item Deux événements indépendants sont toujours incompatibles.
\item L'ellipse d'équation réduite $\dfrac{x^2}{25}+\dfrac{y^2}{9}=1$ a pour demi-grand axe $a=5$.
\item Une parabole possède deux foyers.
\end{enumerate}
\end{exercice}

\begin{exercice}[Définitions à compléter]
Recopie et complète les phrases suivantes avec les mots : \emph{univers, événement, issue, équiprobabilité, épreuve de Bernoulli}.
\begin{enumerate}
\item Une expérience dont le résultat dépend du hasard est appelée \ldots ; chaque résultat possible est une \ldots
\item L'ensemble de toutes les issues possibles est appelé \ldots ; toute partie de cet ensemble est un \ldots
\item On dit qu'il y a \ldots lorsque toutes les issues ont la même probabilité.
\item Une expérience aléatoire qui n'a que deux issues (succès et échec) est une \ldots
\end{enumerate}
\end{exercice}

\begin{exercice}[Calculs directs]
\begin{enumerate}
\item On lance une pièce équilibrée deux fois de suite. Calculer la probabilité d'obtenir deux fois « Face ».
\item On tire une carte au hasard dans un jeu de 32 cartes. Calculer la probabilité de tirer un cœur, puis celle de tirer un roi.
\item Une variable aléatoire $X$ suit la loi $\mathcal{B}(10\,;\,0,3)$. Donner $E(X)$ et $V(X)$.
\item Soit l'hyperbole d'équation $\dfrac{x^2}{16}-\dfrac{y^2}{9}=1$. Donner les valeurs de $a$, $b$, $c$ et les coordonnées des foyers.
\end{enumerate}
\end{exercice}

\courssub{Je m'entraîne}

\begin{exercice}[Tirage dans une urne]
Une urne contient 4 boules blanches et 6 boules noires, indiscernables au toucher. On tire une boule au hasard.
\begin{enumerate}
\item Calculer la probabilité de tirer une boule blanche.
\item On remet la boule dans l'urne et on effectue un second tirage. Calculer la probabilité d'obtenir deux boules noires.
\item Reprendre la question 2 dans le cas d'un tirage sans remise.
\end{enumerate}
\end{exercice}

\begin{exercice}[Probabilité conditionnelle en atelier]
Dans un atelier de mécanique, $60\,\%$ des pièces proviennent de la machine $M_1$ et $40\,\%$ de la machine $M_2$. On sait que $2\,\%$ des pièces de $M_1$ et $5\,\%$ des pièces de $M_2$ sont défectueuses. On prélève une pièce au hasard.
\begin{enumerate}
\item Construire un arbre de probabilités traduisant la situation.
\item Calculer la probabilité que la pièce soit défectueuse.
\item La pièce prélevée est défectueuse. Calculer la probabilité qu'elle provienne de la machine $M_1$.
\end{enumerate}
\end{exercice}

\begin{exercice}[Loi binomiale]
Un technicien contrôle des composants électroniques. La probabilité qu'un composant soit conforme est $0,9$. On prélève au hasard 8 composants (le stock est assez grand pour assimiler les tirages à des tirages avec remise). Soit $X$ le nombre de composants conformes.
\begin{enumerate}
\item Justifier que $X$ suit une loi binomiale et préciser ses paramètres.
\item Calculer $P(X=8)$ et $P(X=0)$.
\item Calculer $P(X\geq 1)$.
\item Donner l'espérance $E(X)$ et interpréter le résultat.
\end{enumerate}
\end{exercice}

\begin{exercice}[Équations réduites de coniques]
Pour chaque conique, donner sa nature, ses éléments caractéristiques (centre, sommets, foyers) et tracer une esquisse.
\begin{enumerate}
\item $\dfrac{x^2}{36}+\dfrac{y^2}{16}=1$ ;
\item $\dfrac{x^2}{9}-\dfrac{y^2}{16}=1$ ;
\item $y^2=12x$.
\end{enumerate}
\end{exercice}

\courssub{Je me prépare au Probatoire}

\begin{exercice}[Type Probatoire : contrôle qualité et loi binomiale]
Une entreprise de Douala fabrique des interrupteurs. On admet que $5\,\%$ des interrupteurs produits sont défectueux.

\textbf{Partie A.} On prélève un interrupteur au hasard dans la production, puis un second après remise du premier.
\begin{enumerate}
\item Calculer la probabilité que les deux interrupteurs soient défectueux.
\item Calculer la probabilité qu'au moins l'un des deux soit défectueux.
\end{enumerate}

\textbf{Partie B.} Pour un contrôle, on prélève un échantillon de 12 interrupteurs au hasard. On assimile ce prélèvement à un tirage avec remise. Soit $X$ la variable aléatoire égale au nombre d'interrupteurs défectueux dans l'échantillon.
\begin{enumerate}
\item Justifier que $X$ suit la loi binomiale $\mathcal{B}(12\,;\,0,05)$.
\item Calculer $P(X=0)$, arrondi à $10^{-3}$.
\item En déduire $P(X\geq 1)$.
\item Calculer $E(X)$ et $V(X)$. Interpréter $E(X)$ dans le contexte de l'exercice.
\end{enumerate}

\textbf{Partie C.} Un client exige qu'un lot soit refusé si l'échantillon contient plus de 2 interrupteurs défectueux. Sachant que $P(X\leq 2)\approx 0,980$, calculer la probabilité que le lot soit refusé.
\end{exercice}

\begin{exercice}[Type Probatoire : urne et probabilités conditionnelles]
Une urne contient 5 boules rouges et 3 boules vertes, indiscernables au toucher. On tire successivement et sans remise 2 boules de l'urne.
\begin{enumerate}
\item Construire un arbre de probabilités décrivant cette expérience, en indiquant sur chaque branche la probabilité conditionnelle correspondante.
\item Calculer la probabilité d'obtenir deux boules rouges.
\item Calculer la probabilité d'obtenir deux boules de même couleur.
\item Calculer la probabilité d'obtenir au moins une boule verte.
\item Sachant que la deuxième boule tirée est verte, calculer la probabilité que la première ait été rouge.
\item On répète maintenant l'expérience suivante : on tire une boule, on note sa couleur et on la remet dans l'urne, et cela 4 fois de suite. Soit $Y$ le nombre de boules rouges obtenues. Déterminer la loi de $Y$, son espérance et $P(Y=2)$.
\end{enumerate}
\end{exercice}

\begin{exercice}[Type Probatoire : coniques et application technique]
Le plan est muni d'un repère orthonormé $(O\,;\,\vec{i},\vec{j})$ (unité : 1 cm). Une entreprise doit concevoir un miroir elliptique dont le contour est le lieu des points $M$ tels que $MF+MF'=10$, où $F(-3\,;\,0)$ et $F'(3\,;\,0)$ sont les foyers.
\begin{enumerate}
\item Montrer que ce lieu est une ellipse et déterminer son équation réduite.
\item Donner les coordonnées des sommets de cette ellipse et calculer son excentricité.
\item Tracer l'ellipse dans le repère (on placera les foyers et les sommets).
\item On considère maintenant la parabole $\mathcal{P}$ d'équation $y^2=8x$.
\begin{enumerate}
\item Donner le foyer et la directrice de $\mathcal{P}$.
\item Déterminer une équation de la tangente à $\mathcal{P}$ au point $A(2\,;\,4)$.
\item Vérifier si cette tangente passe par le point $F(-2\,;\,0)$ et conclure.
\end{enumerate}
\item Le technicien affirme : « le lieu des points $M$ tels que $|MF_1-MF_2|=6$, avec $F_1(-5\,;\,0)$ et $F_2(5\,;\,0)$, est une hyperbole ». Justifier cette affirmation, donner l'équation réduite de cette hyperbole et déterminer les équations de ses asymptotes.
\end{enumerate}
\end{exercice}

\newpage

\coursec{Corrigés des exercices}

\coursec{Vocabulaire et expériences aléatoires}

\begin{corrige}[Exercice 1 — QCM : vocabulaire des probabilités]
\begin{enumerate}
\item Réponse \textbf{(b)} : $\Omega=\{1;2;3;4;5;6\}$. L'univers contient toutes les issues possibles du lancer de dé, et seulement celles-là.
\item Réponse \textbf{(a)} : $P(A\cup B)=P(A)+P(B)$. C'est la définition des événements incompatibles : $A\cap B=\varnothing$, donc $P(A\cap B)=0$.
\item Réponse \textbf{(b)} : $P(\overline{A})=1-P(A)$.
\item Réponse \textbf{(a)} : $E(X)=np$. (La variance est $np(1-p)$, à ne pas confondre.)
\end{enumerate}
\end{corrige}

\begin{corrige}[Exercice 2 — Vrai ou faux, en justifiant]
\begin{enumerate}
\item \textbf{Faux.} Une probabilité est toujours comprise entre $0$ et $1$ : $0\leq P(A)\leq 1$. Or $1,2>1$.
\item \textbf{Vrai.} $P(\overline{A})=1-P(A)=1-0,4=0,6$.
\item \textbf{Faux.} L'indépendance signifie $P(A\cap B)=P(A)\times P(B)$ ; l'incompatibilité signifie $P(A\cap B)=0$. Si $P(A)>0$ et $P(B)>0$, deux événements indépendants ne sont jamais incompatibles. Exemple : deux lancers de pièce indépendants, « Face au premier » et « Face au second » ne sont pas incompatibles.
\item \textbf{Vrai.} L'équation réduite est $\dfrac{x^2}{a^2}+\dfrac{y^2}{b^2}=1$ avec $a^2=25$, donc $a=5$ (et $b=3$).
\item \textbf{Faux.} Une parabole possède un \emph{seul} foyer (et une directrice). Ce sont l'ellipse et l'hyperbole qui ont deux foyers.
\end{enumerate}
\end{corrige}

\begin{corrige}[Exercice 3 — Définitions à compléter]
\begin{enumerate}
\item Une expérience dont le résultat dépend du hasard est appelée \cle{expérience aléatoire} ; chaque résultat possible est une \cle{issue}.
\item L'ensemble de toutes les issues possibles est appelé \cle{univers} ; toute partie de cet ensemble est un \cle{événement}.
\item On dit qu'il y a \cle{équiprobabilité} lorsque toutes les issues ont la même probabilité.
\item Une expérience aléatoire qui n'a que deux issues (succès et échec) est une \cle{épreuve de Bernoulli}.
\end{enumerate}
\end{corrige}

\coursec{Probabilité d'un événement et probabilités conditionnelles}

\begin{corrige}[Exercice 4 — Calculs directs]
\begin{enumerate}
\item Les deux lancers sont indépendants et $P(F)=\dfrac12$ à chaque lancer :
\[P(F\text{ puis }F)=\dfrac12\times\dfrac12=\dfrac14=0,25.\]
\item Un jeu de 32 cartes contient 8 cœurs et 4 rois. Par équiprobabilité :
\[P(\text{cœur})=\dfrac{8}{32}=\dfrac14=0,25 \qquad ; \qquad P(\text{roi})=\dfrac{4}{32}=\dfrac18=0,125.\]
\item Pour $X\sim\mathcal{B}(10\,;\,0,3)$ :
\[E(X)=np=10\times 0,3=3 \qquad ; \qquad V(X)=np(1-p)=10\times0,3\times0,7=2,1.\]
\item L'équation est de la forme $\dfrac{x^2}{a^2}-\dfrac{y^2}{b^2}=1$ avec $a^2=16$ et $b^2=9$, donc $a=4$ et $b=3$. Pour une hyperbole : $c^2=a^2+b^2=16+9=25$, donc $c=5$. Les foyers sont $F(-5\,;\,0)$ et $F'(5\,;\,0)$.
\end{enumerate}
\end{corrige}

\begin{corrige}[Exercice 5 — Tirage dans une urne]
L'urne contient $4+6=10$ boules.
\begin{enumerate}
\item Par équiprobabilité : $P(B)=\dfrac{4}{10}=\dfrac25=0,4$.
\item Avec remise, les deux tirages sont indépendants et $P(N)=\dfrac{6}{10}=0,6$ à chaque tirage :
\[P(N\text{ puis }N)=0,6\times0,6=0,36.\]
\item Sans remise : au premier tirage $P(N)=\dfrac{6}{10}$ ; sachant qu'une noire est sortie, il reste 5 noires sur 9 boules, donc $P_{N_1}(N_2)=\dfrac59$ :
\[P(N_1\cap N_2)=\dfrac{6}{10}\times\dfrac59=\dfrac{30}{90}=\dfrac13\approx0,333.\]
\end{enumerate}
\end{corrige}

\begin{corrige}[Exercice 6 — Probabilité conditionnelle en atelier]
On note $D$ l'événement « la pièce est défectueuse ».
\begin{enumerate}
\item Arbre de probabilités :
\begin{popfigure}
\begin{tikzpicture}[grow=right, level distance=3.2cm,
level 1/.style={sibling distance=2.6cm}, level 2/.style={sibling distance=1.3cm},
every node/.style={font=\small}]
\node {Pièce}
  child { node {$M_2$}
    child { node {$\overline{D}$} edge from parent node[below]{$0,95$} }
    child { node {$D$} edge from parent node[above]{$0,05$} }
    edge from parent node[below]{$0,4$} }
  child { node {$M_1$}
    child { node {$\overline{D}$} edge from parent node[below]{$0,98$} }
    child { node {$D$} edge from parent node[above]{$0,02$} }
    edge from parent node[above]{$0,6$} };
\end{tikzpicture}
\legende{Arbre de probabilités de la production de l'atelier}
\end{popfigure}
\item Par la formule des probabilités totales ($M_1$ et $M_2$ forment une partition de l'univers) :
\begin{align*}
P(D)&=P(M_1)\times P_{M_1}(D)+P(M_2)\times P_{M_2}(D)\\
&=0,6\times0,02+0,4\times0,05=0,012+0,020=0,032.
\end{align*}
\item Par la formule des probabilités conditionnelles :
\[P_D(M_1)=\dfrac{P(M_1\cap D)}{P(D)}=\dfrac{0,012}{0,032}=0,375.\]
La probabilité que la pièce défectueuse provienne de $M_1$ est $0,375$.
\end{enumerate}
\end{corrige}

\coursec{Variables aléatoires et épreuves de Bernoulli}

\begin{corrige}[Exercice 7 — Loi binomiale]
\begin{enumerate}
\item On répète $n=8$ fois la même épreuve de Bernoulli (succès : « le composant est conforme », de probabilité $p=0,9$), de façon identique et indépendante (tirages assimilés à des tirages avec remise). Donc $X$ suit la loi binomiale $\mathcal{B}(8\,;\,0,9)$.
\item 
\begin{align*}
P(X=8)&=\binom{8}{8}(0,9)^8(0,1)^0=(0,9)^8\approx0,4305.\\
P(X=0)&=\binom{8}{0}(0,9)^0(0,1)^8=(0,1)^8=10^{-8}.
\end{align*}
\item On passe par l'événement contraire :
\[P(X\geq1)=1-P(X=0)=1-10^{-8}\approx1.\]
\item $E(X)=np=8\times0,9=7,2$. Interprétation : sur un grand nombre de prélèvements de 8 composants, on trouve en moyenne $7,2$ composants conformes par échantillon.
\end{enumerate}
\end{corrige}

\begin{corrige}[Exercice 9 — Type Probatoire : contrôle qualité et loi binomiale]
\textbf{Partie A.} On note $D$ : « l'interrupteur est défectueux », $P(D)=0,05$. Les tirages sont indépendants (remise).
\begin{enumerate}
\item $P(D\cap D)=0,05\times0,05=0,0025$.
\item L'événement contraire de « au moins un défectueux » est « aucun défectueux » :
\[P(\text{au moins un})=1-(0,95)^2=1-0,9025=0,0975.\]
\end{enumerate}
\textbf{Partie B.}
\begin{enumerate}
\item On répète $n=12$ fois, de façon identique et indépendante (tirage assimilé à un tirage avec remise), l'épreuve de Bernoulli de succès « interrupteur défectueux » de probabilité $p=0,05$. Donc $X\sim\mathcal{B}(12\,;\,0,05)$.
\item $P(X=0)=\binom{12}{0}(0,05)^0(0,95)^{12}=(0,95)^{12}\approx0,540$ (arrondi à $10^{-3}$).
\item $P(X\geq1)=1-P(X=0)\approx1-0,540=0,460$.
\item $E(X)=np=12\times0,05=0,6$ et $V(X)=np(1-p)=12\times0,05\times0,95=0,57$. Interprétation : en moyenne, un échantillon de 12 interrupteurs contient $0,6$ interrupteurs défectueux (soit environ 6 interrupteurs défectueux pour 10 échantillons).
\end{enumerate}
\textbf{Partie C.} Le lot est refusé si $X>2$, c'est-à-dire $X\geq3$ :
\[P(X\geq3)=1-P(X\leq2)\approx1-0,980=0,020.\]
La probabilité que le lot soit refusé est environ $0,02$.

\textbf{Barème indicatif (sur 10)} : Partie A : 2 pts ; B.1 (justification complète) : 2 pts ; B.2–B.3 : 2 pts ; B.4 (calculs + interprétation) : 2 pts ; Partie C : 2 pts.

\textbf{Points de vigilance} : citer explicitement les trois conditions de la loi binomiale (répétitions identiques, indépendance, deux issues) ; écrire la formule $\binom{n}{k}p^k(1-p)^{n-k}$ avant l'application numérique ; pour « au moins un », justifier le passage à l'événement contraire ; l'interprétation de $E(X)$ est exigée, ne pas l'oublier.
\end{corrige}

\begin{corrige}[Exercice 10 — Type Probatoire : urne et probabilités conditionnelles]
\begin{enumerate}
\item Arbre (tirage sans remise, 8 boules au départ) :
\begin{popfigure}
\begin{tikzpicture}[grow=right, level distance=3.4cm,
level 1/.style={sibling distance=2.6cm}, level 2/.style={sibling distance=1.3cm},
every node/.style={font=\small}]
\node {Départ}
  child { node {$V_1$}
    child { node {$V_2$} edge from parent node[below]{$2/7$} }
    child { node {$R_2$} edge from parent node[above]{$5/7$} }
    edge from parent node[below]{$3/8$} }
  child { node {$R_1$}
    child { node {$V_2$} edge from parent node[below]{$3/7$} }
    child { node {$R_2$} edge from parent node[above]{$4/7$} }
    edge from parent node[above]{$5/8$} };
\end{tikzpicture}
\legende{Tirage successif sans remise de 2 boules}
\end{popfigure}
\item $P(R_1\cap R_2)=\dfrac58\times\dfrac47=\dfrac{20}{56}=\dfrac{5}{14}\approx0,357$.
\item « Même couleur » $=(R_1\cap R_2)\cup(V_1\cap V_2)$, événements incompatibles :
\[P=\dfrac58\times\dfrac47+\dfrac38\times\dfrac27=\dfrac{20}{56}+\dfrac{6}{56}=\dfrac{26}{56}=\dfrac{13}{28}\approx0,464.\]
\item « Au moins une verte » est le contraire de « deux rouges » :
\[P=1-\dfrac{5}{14}=\dfrac{9}{14}\approx0,643.\]
\item On cherche $P_{V_2}(R_1)$. D'abord, par probabilités totales :
\[P(V_2)=\dfrac58\times\dfrac37+\dfrac38\times\dfrac27=\dfrac{15}{56}+\dfrac{6}{56}=\dfrac{21}{56}=\dfrac38.\]
Puis :
\[P_{V_2}(R_1)=\dfrac{P(R_1\cap V_2)}{P(V_2)}=\dfrac{15/56}{21/56}=\dfrac{15}{21}=\dfrac57\approx0,714.\]
\item On répète $n=4$ fois, de façon identique et indépendante (avec remise), l'épreuve de Bernoulli de succès « boule rouge » de probabilité $p=\dfrac58$. Donc $Y\sim\mathcal{B}\left(4\,;\,\dfrac58\right)$.
\[E(Y)=np=4\times\dfrac58=\dfrac52=2,5.\]
\[P(Y=2)=\binom{4}{2}\left(\dfrac58\right)^2\left(\dfrac38\right)^2=6\times\dfrac{25}{64}\times\dfrac{9}{64}=\dfrac{1350}{4096}=\dfrac{675}{2048}\approx0,330.\]
\end{enumerate}
\textbf{Barème indicatif (sur 10)} : arbre : 2 pts ; questions 2 et 3 : 2 pts ; question 4 : 1 pt ; question 5 : 2 pts ; question 6 : 3 pts.

\textbf{Points de vigilance} : bien distinguer les probabilités du second tirage selon que le tirage est avec ou sans remise ; en question 5, calculer d'abord $P(V_2)$ par la formule des probabilités totales ; en question 6, justifier la loi binomiale en citant l'indépendance due à la remise ; simplifier les fractions.
\end{corrige}

\coursec{Présentation générale des coniques, tangentes et applications}

\begin{corrige}[Exercice 8 — Équations réduites de coniques]
\begin{enumerate}
\item $\dfrac{x^2}{36}+\dfrac{y^2}{16}=1$ : \textbf{ellipse} de centre $O(0\,;\,0)$, avec $a^2=36$ donc $a=6$, $b^2=16$ donc $b=4$, et $c^2=a^2-b^2=20$ donc $c=2\sqrt5$. Sommets : $A(6\,;\,0)$, $A'(-6\,;\,0)$, $B(0\,;\,4)$, $B'(0\,;\,-4)$. Foyers : $F(2\sqrt5\,;\,0)$ et $F'(-2\sqrt5\,;\,0)$.
\item $\dfrac{x^2}{9}-\dfrac{y^2}{16}=1$ : \textbf{hyperbole} de centre $O$, avec $a=3$, $b=4$, $c^2=a^2+b^2=25$ donc $c=5$. Sommets : $(3\,;\,0)$ et $(-3\,;\,0)$. Foyers : $(5\,;\,0)$ et $(-5\,;\,0)$. Asymptotes : $y=\pm\dfrac{4}{3}x$.
\item $y^2=12x$ : \textbf{parabole} de sommet $O(0\,;\,0)$, d'axe $(Ox)$, de la forme $y^2=2px$ avec $2p=12$ donc $p=6$. Foyer $F(3\,;\,0)$, directrice $x=-3$.
\end{enumerate}
\begin{popfigure}
\begin{tikzpicture}[scale=0.32]
\draw[->,gray] (-7.5,0)--(7.5,0); \draw[->,gray] (0,-5.5)--(0,5.5);
\draw[popblue,thick] (0,0) ellipse (6 and 4);
\fill[popink] (4.47,0) circle (2.5pt) node[below right]{\scriptsize $F$};
\fill[popink] (-4.47,0) circle (2.5pt) node[below left]{\scriptsize $F'$};
\node[popblue] at (0,4.6) {\scriptsize ellipse};
\end{tikzpicture}\hfill
\begin{tikzpicture}[scale=0.32]
\draw[->,gray] (-7.5,0)--(7.5,0); \draw[->,gray] (0,-5.5)--(0,5.5);
\draw[gray,dashed,domain=-5:5] plot(\x,{4/3*\x});
\draw[gray,dashed,domain=-5:5] plot(\x,{-4/3*\x});
\draw[poporange,thick,domain=3:6.5,samples=40] plot(\x,{4/3*sqrt(max(0,\x*\x-9))});
\draw[poporange,thick,domain=3:6.5,samples=40] plot(\x,{-4/3*sqrt(max(0,\x*\x-9))});
\draw[poporange,thick,domain=-6.5:-3,samples=40] plot(\x,{4/3*sqrt(max(0,\x*\x-9))});
\draw[poporange,thick,domain=-6.5:-3,samples=40] plot(\x,{-4/3*sqrt(max(0,\x*\x-9))});
\fill[popink] (5,0) circle (2.5pt); \fill[popink] (-5,0) circle (2.5pt);
\node[poporange] at (0,5) {\scriptsize hyperbole};
\end{tikzpicture}\hfill
\begin{tikzpicture}[scale=0.32]
\draw[->,gray] (-4.5,0)--(7.5,0); \draw[->,gray] (0,-5.5)--(0,5.5);
\draw[popgreen,thick,domain=0:6.5,samples=60] plot(\x,{sqrt(max(0,12*\x))});
\draw[popgreen,thick,domain=0:6.5,samples=60] plot(\x,{-sqrt(max(0,12*\x))});
\draw[gray,dashed] (-3,-5.5)--(-3,5.5);
\fill[popink] (3,0) circle (2.5pt) node[below right]{\scriptsize $F$};
\node[popgreen] at (4.5,4.5) {\scriptsize parabole};
\end{tikzpicture}
\legende{Esquisses des trois coniques}
\end{popfigure}
\end{corrige}

\begin{corrige}[Exercice 11 — Type Probatoire : coniques et application technique]
\begin{enumerate}
\item Le lieu des points $M$ tels que $MF+MF'=10$ (constante) avec $F$ et $F'$ fixes est, par définition, une \cle{ellipse} de foyers $F$ et $F'$. On a $2a=10$ donc $a=5$, et $FF'=2c=6$ donc $c=3$. Pour une ellipse : $b^2=a^2-c^2=25-9=16$, donc $b=4$. Le centre est le milieu de $[FF']$, c'est $O$. L'équation réduite est :
\[\boxed{\dfrac{x^2}{25}+\dfrac{y^2}{16}=1}\]
\item Sommets : $A(5\,;\,0)$, $A'(-5\,;\,0)$, $B(0\,;\,4)$, $B'(0\,;\,-4)$. Excentricité :
\[e=\dfrac{c}{a}=\dfrac35=0,6.\]
\item Tracé :
\begin{popfigure}
\begin{tikzpicture}[scale=0.55]
\draw[->,gray] (-6.5,0)--(6.5,0) node[right]{\scriptsize $x$};
\draw[->,gray] (0,-5)--(0,5) node[above]{\scriptsize $y$};
\draw[popblue,thick] (0,0) ellipse (5 and 4);
\fill[popink] (3,0) circle (2.5pt) node[below right]{\scriptsize $F'(3;0)$};
\fill[popink] (-3,0) circle (2.5pt) node[below left]{\scriptsize $F(-3;0)$};
\fill[poporange] (5,0) circle (2.5pt) node[above right]{\scriptsize $A$};
\fill[poporange] (-5,0) circle (2.5pt) node[above left]{\scriptsize $A'$};
\fill[poporange] (0,4) circle (2.5pt) node[above right]{\scriptsize $B$};
\fill[poporange] (0,-4) circle (2.5pt) node[below right]{\scriptsize $B'$};
\end{tikzpicture}
\legende{Ellipse $\dfrac{x^2}{25}+\dfrac{y^2}{16}=1$ : miroir elliptique}
\end{popfigure}
\item 
\begin{enumerate}
\item $\mathcal{P}$ : $y^2=8x$ est de la forme $y^2=2px$ avec $2p=8$, donc $p=4$. Foyer $F(2\,;\,0)$, directrice d'équation $x=-2$.
\item La tangente à la parabole $y^2=2px$ au point $A(x_0\,;\,y_0)$ a pour équation $y\,y_0=p(x+x_0)$. Ici $p=4$, $x_0=2$, $y_0=4$ :
\[4y=4(x+2)\quad\Longleftrightarrow\quad y=x+2.\]
(Vérification : $A(2\,;\,4)$ appartient bien à $\mathcal{P}$ car $4^2=16=8\times2$, et $4=2+2$ donc $A$ est sur la tangente.)
\item Pour $F(-2\,;\,0)$ : $y_F=0$ et $x_F+2=0$, donc $0=0$ : le point $F(-2\,;\,0)$ appartient à la tangente. Conclusion : la tangente en $A$ à $\mathcal{P}$ passe par le point $(-2\,;\,0)$, intersection de la directrice avec l'axe de la parabole.
\end{enumerate}
\item Le lieu des points $M$ tels que $|MF_1-MF_2|=6$ (constante inférieure à $F_1F_2=10$) est, par définition, une \cle{hyperbole} de foyers $F_1$ et $F_2$ : l'affirmation du technicien est exacte. On a $2a=6$ donc $a=3$, et $2c=10$ donc $c=5$. Pour une hyperbole : $b^2=c^2-a^2=25-9=16$, donc $b=4$. Équation réduite :
\[\dfrac{x^2}{9}-\dfrac{y^2}{16}=1.\]
Les asymptotes ont pour équations :
\[y=\dfrac{b}{a}x=\dfrac43x \qquad\text{et}\qquad y=-\dfrac43x.\]
\end{enumerate}
\textbf{Barème indicatif (sur 10)} : question 1 (définition + calculs + équation) : 2,5 pts ; question 2 : 1,5 pts ; question 3 (tracé soigné) : 1 pt ; question 4 : 3 pts ; question 5 : 2 pts.

\textbf{Points de vigilance} : citer la définition bifocale de la conique ; ne pas confondre les relations $b^2=a^2-c^2$ (ellipse) et $c^2=a^2+b^2$ (hyperbole) ; pour la tangente, vérifier d'abord que le point appartient à la parabole ; pour l'hyperbole, vérifier que la constante $2a$ est strictement inférieure à la distance focale $2c$ (ici $6<10$), condition d'existence.
\end{corrige}

\newpage

\coursec{Fiche bilan}

\begin{popfigure}
\begin{tikzpicture}[
    mindmap,
    grow cyclic,
    every node/.style={concept, execute at begin node=\hskip0pt},
    concept color=popblue!80,
    level 1/.style={level distance=4.5cm, sibling angle=60, font=\bfseries\footnotesize},
    level 2/.style={level distance=3cm, sibling angle=45, font=\scriptsize},
    branch1/.style={concept color=popblue, font=\bfseries\small},
    branch2/.style={concept color=poporange, font=\bfseries\small},
    branch3/.style={concept color=popgreen, font=\bfseries\small},
    branch4/.style={concept color=poppurple, font=\bfseries\small},
    branch5/.style={concept color=poppink, font=\bfseries\small},
    branch6/.style={concept color=popgold, font=\bfseries\small},
    text=popdark
]
\node[concept, text=white, font=\bfseries\large, align=center]{Probabilités \& Coniques\\Terminale Tech}
    child[branch1] { node[concept] {1. Vocabulaire \& Expériences} 
        child { node {Univers $\Omega$, Événements} }
        child { node {Opérations : $\cup, \cap, \bar{\cdot}$} }
        child { node {Incompatibles, Contraires} }
        child { node {Expérience aléatoire} }
    }
    child[branch2] { node[concept] {2. Probabilités \& Conditionnelles} 
        child { node {Équiprobabilité : $P(A)=\frac{\text{Card}A}{\text{Card}\Omega}$} }
        child { node {Propriétés : $0\le P\le 1$, $P(\bar A)=1-P(A)$} }
        child { node {Formule $P(A\cup B)$} }
        child { node {Conditionnelle $P_A(B)=\frac{P(A\cap B)}{P(A)}$} }
        child { node {Indépendance : $P(A\cap B)=P(A)P(B)$} }
    }
    child[branch3] { node[concept] {3. Variables Aléatoires \& Bernoulli} 
        child { node {Loi : $x_i, p_i$ ; $\sum p_i=1$} }
        child { node {Espérance $E(X)=\sum x_i p_i$} }
        child { node {Variance $V(X)=E(X^2)-E(X)^2$} }
        child { node {Bernoulli $B(n,p)$ : $P(X=k)=\binom{n}{k}p^k q^{n-k}$} }
        child { node {$E(X)=np$, $V(X)=npq$} }
    }
    child[branch4] { node[concept] {4. Coniques : Généralités} 
        child { node {Définition : Foyer $F$, Directrice $D$, $e$} }
        child { node {Ellipse ($e<1$), Parabole ($e=1$), Hyperbole ($e>1$)} }
        child { node {Équations réduites} }
        child { node {Axes, Sommets, Foyers, Excentricité} }
    }
    child[branch5] { node[concept] {5. Tangante en un point} 
        child { node {Ellipse $\frac{xx_0}{a^2}+\frac{yy_0}{b^2}=1$} }
        child { node {Parabole $y^2=2px \to yy_0=p(x+x_0)$} }
        child { node {Hyperbole $\frac{xx_0}{a^2}-\frac{yy_0}{b^2}=1$} }
        child { node {Propriété optique (bissectrice)} }
    }
    child[branch6] { node[concept] {6. Construction \& Applications} 
        child { node {Méthode du jardinier (ellipse)} }
        child { node {Plis de papier (parabole)} }
        child { node {Règle \& Compas (hyperbole)} }
        child { node {Applications : Optique, Mécanique, Architecture} }
    };
\end{tikzpicture}
\legende{Carte mentale : Leçon 5 -- Probabilités et Coniques (Terminale Technique)}
\end{popfigure}

\begin{bilan}
\courssub{Vocabulaire et expériences aléatoires}
\begin{itemize}
    \item \cle{Univers $\Omega$} : ensemble des issues possibles. \cle{Événement} : partie de $\Omega$.
    \item Opérations : réunion $A\cup B$, intersection $A\cap B$, complémentaire $\bar A = \Omega \setminus A$.
    \item \cle{Incompatibles} : $A\cap B = \emptyset$. \cle{Contraires} : $A\cup B=\Omega$ et $A\cap B=\emptyset$ ($B=\bar A$).
    \item \cle{Expérience aléatoire} : issues non déterministes, reproductibles, issues connues.
\end{itemize}

\courssub{Probabilité d'un événement et probabilités conditionnelles}
\begin{itemize}
    \item \textbf{Équiprobabilité} : $P(A) = \dfrac{\text{Card}(A)}{\text{Card}(\Omega)}$.
    \item \textbf{Propriétés fondamentales} :
    \begin{align*}
        & 0 \le P(A) \le 1,\quad P(\Omega)=1,\quad P(\emptyset)=0 \\
        & P(\bar A) = 1 - P(A) \\
        & P(A \cup B) = P(A) + P(B) - P(A \cap B) \quad (\text{si incompatibles : } P(A\cup B)=P(A)+P(B))
    \end{align*}
    \item \textbf{Probabilité conditionnelle} : $P_A(B) = \dfrac{P(A \cap B)}{P(A)}$ ($P(A)>0$).
    \item \textbf{Indépendance} : $A$ et $B$ indépendants $\iff P(A \cap B) = P(A)P(B) \iff P_A(B)=P(B)$.
    \item \textbf{Formule des probabilités totales} : si $(B_i)$ partition de $\Omega$, $P(A) = \sum P(B_i)P_{B_i}(A)$.
\end{itemize}

\courssub{Variables aléatoires et épreuves de Bernoulli}
\begin{itemize}
    \item \cle{Variable aléatoire $X$} : application $\Omega \to \mathbb{R}$. Loi : tableau $(x_i, p_i)$ avec $\sum p_i = 1$.
    \item \textbf{Espérance} : $E(X) = \sum x_i p_i$. \textbf{Variance} : $V(X) = E(X^2) - [E(X)]^2 = \sum (x_i - E(X))^2 p_i$. \textbf{Écart-type} : $\sigma(X) = \sqrt{V(X)}$.
    \item \textbf{Propriétés linéaires} : $E(aX+b)=aE(X)+b$, $V(aX+b)=a^2V(X)$.
    \item \cle{Épreuves de Bernoulli} : $n$ épreuves indépendantes, même prob. de succès $p$, échec $q=1-p$.
    \item \cle{Loi Binomiale $B(n,p)$} : $X \sim B(n,p) \iff P(X=k) = \binom{n}{k} p^k q^{n-k}$.
    \item $E(X) = np$, $V(X) = npq$, $\sigma(X) = \sqrt{npq}$.
    \item \textbf{Intervalle de fluctuation} (seuil 5\% ou 95\%) : $[np - 1,96\sqrt{npq} ; np + 1,96\sqrt{npq}]$ (arrondi).
\end{itemize}

\courssub{Présentation générale des coniques et propriétés}
\begin{itemize}
    \item \cle{Définition unifiée} : Lieu des points $M$ tel que $MF = e \cdot d(M,D)$ ($F$ foyer, $D$ directrice, $e$ excentricité).
    \item \textbf{Ellipse} ($0<e<1$) : $\dfrac{x^2}{a^2} + \dfrac{y^2}{b^2} = 1$ ($a>b>0$, $c=\sqrt{a^2-b^2}$, $e=c/a$). Foyers $F(\pm c,0)$. $MF+MF'=2a$.
    \item \textbf{Parabole} ($e=1$) : $y^2 = 2px$ ($p>0$). Foyer $F(p/2,0)$, Directrice $x=-p/2$. $MF = d(M,D)$.
    \item \textbf{Hyperbole} ($e>1$) : $\dfrac{x^2}{a^2} - \dfrac{y^2}{b^2} = 1$ ($c=\sqrt{a^2+b^2}$, $e=c/a$). Foyers $F(\pm c,0)$. $|MF-MF'|=2a$. Asymptotes $y=\pm\frac{b}{a}x$.
\end{itemize}

\courssub{Tangente en un point à une conique}
\begin{itemize}
    \item \textbf{Ellipse} $\frac{x^2}{a^2}+\frac{y^2}{b^2}=1$ en $M(x_0,y_0)$ : $\frac{xx_0}{a^2}+\frac{yy_0}{b^2}=1$.
    \item \textbf{Parabole} $y^2=2px$ en $M(x_0,y_0)$ : $yy_0 = p(x+x_0)$.
    \item \textbf{Hyperbole} $\frac{x^2}{a^2}-\frac{y^2}{b^2}=1$ en $M(x_0,y_0)$ : $\frac{xx_0}{a^2}-\frac{yy_0}{b^2}=1$.
    \item \textbf{Propriété optique} : La tangente est la bissectrice de l'angle formé par les droites joignant le point de contact aux foyers (bisectrice intérieure pour l'ellipse, extérieure pour l'hyperbole, bissectrice de l'angle entre la focale et la parallèle à l'axe pour la parabole).
\end{itemize}

\courssub{Construction d'une conique et applications intégrées}
\begin{itemize}
    \item \textbf{Ellipse} : Méthode du jardinier (ficelle de longueur $2a$, pointes aux foyers).
    \item \textbf{Parabole} : Plis de papier (point $F$ sur bord $D$, plis tangents) ou méthode du rectangle.
    \item \textbf{Hyperbole} : Règle et compas (différence de distances constante $2a$) ou asymptotes + point.
    \item \textbf{Applications techniques} :
    \begin{itemize}
        \item Optique : Réflecteurs paraboliques (phares, antennes), miroirs elliptiques.
        \item Mécanique : Vilebrequins (ellipse), engrenages coniques.
        \item Architecture : Arches (parabole/ellipse), structures tendues (hyperboloïdes).
        \item Probabilités : Modélisation aléa (contrôle qualité $B(n,p)$, files d'attente).
    \end{itemize}
\end{itemize}
\end{bilan}

\begin{aretenir}[Checklist avant le Probatoire]
\begin{itemize}
    \item $\square$ Définir univers, événement, opérations ensemblistes, expérience aléatoire.
    \item $\square$ Calculer $P(A)$ en équiprobabilité ; appliquer $P(\bar A)$, $P(A\cup B)$, $P(A\cap B)$.
    \item $\square$ Manipuler $P_A(B)$ et tester l'indépendance ($P(A\cap B)=P(A)P(B)$).
    \item $\square$ Construire la loi d'une v.a. $X$ ; calculer $E(X)$, $V(X)$, $\sigma(X)$.
    \item $\square$ Identifier un schéma de Bernoulli ; écrire $P(X=k)$ pour $X\sim B(n,p)$ ; donner $E(X), V(X)$.
    \item $\square$ Déterminer un intervalle de fluctuation (seuil 95\%).
    \item $\square$ Reconnaître une conique (équation, $e$, foyers, directrices, asymptotes).
    \item $\square$ Écrire l'équation de la tangente en $M(x_0,y_0)$ pour les 3 coniques.
    \item $\square$ Utiliser la propriété optique (bissectrice) pour justifier une construction.
    \item $\square$ Décrire une méthode de construction (jardinier, plis, règle-compas).
    \item $\square$ Résoudre un problème mêlant probabilités (loi binomiale) et géométrie (conique) ou application technique.
    \item $\square$ Rédiger une démonstration complète (théorème probabilités totales, formule tangente, propriété focale).
\end{itemize}
\end{aretenir}

\findecours

\end{document}
